\begin{center}\rule{0.5\linewidth}{0.5pt}\end{center}

\hypertarget{bhm_aqmes_med_plan0.5.0}{%
\subsection{bhm\_aqmes\_med}\label{bhm_aqmes_med_plan0.5.0}}

\noindent\textbf{target\_selection plan:} 0.5.0

\noindent\textbf{target\_selection tag:}
\href{https://github.com/sdss/target_selection/tree/0.3.0/}{0.3.0}

\noindent\textbf{crossmatch plan:} 0.5.0

\noindent\textbf{Summary:} Spectroscopically confirmed optically bright SDSS
QSOs, selected from the SDSS QSO catalog (DR16Q,
\citealt{Lyke2020}). Located in 36 mostly disjoint fields within the SDSS QSO
footprint that were pre-selected to contain higher than average numbers
of bright QSOs and CSC targets. The list of field centres can be found
within
\href{https://github.com/sdss/target_selection/blob/0.3.0/python/target_selection/masks/candidate_target_fields_bhm_aqmes_med_v0.3.1.fits}{the
target\_selection repository}.

\noindent\textbf{Simplified description of selection criteria:} Select all
objects from SDSS DR16 QSO catalog that have
16.0\textless{}\texttt{sdss\_psfmag\_i}\textless19.1~AB, that lie within
1.49~degrees of at least one AQMES-medium field location.


\noindent\textbf{Target priority options:} 1100

\noindent\textbf{Cadence options:} \texttt{dark\_10x4\_4yr}

\noindent\textbf{Implementation:}
\href{https://github.com/sdss/target_selection/blob/0.3.0/python/target_selection/cartons/bhm_aqmes.py}{bhm\_aqmes.py}

\noindent\textbf{Number of targets:} 2663

\begin{center}\rule{0.5\linewidth}{0.5pt}\end{center}

\hypertarget{bhm_aqmes_med_faint_plan0.5.0}{%
\subsection{bhm\_aqmes\_med\_faint}\label{bhm_aqmes_med_faint_plan0.5.0}}

\noindent\textbf{target\_selection plan:} 0.5.0

\noindent\textbf{target\_selection tag:}
\href{https://github.com/sdss/target_selection/tree/0.3.0/}{0.3.0}

\noindent\textbf{crossmatch plan:} 0.5.0

\noindent\textbf{Summary:} Spectroscopically confirmed optically faint SDSS QSOs,
selected from the SDSS QSO catalog (DR16Q,
\citealt{Lyke2020}). Located in 36 mostly disjoint fields within the SDSS QSO
footprint that were pre-selected to contain higher than average numbers
of bright QSOs and CSC targets. The list of field centres can be found
within
\href{https://github.com/sdss/target_selection/blob/0.3.0/python/target_selection/masks/candidate_target_fields_bhm_aqmes_med_v0.3.1.fits}{the
target\_selection repository}.

\noindent\textbf{Simplified description of selection criteria:} Select all
objects from SDSS DR16 QSO catalog that have
19.1\textless{}\texttt{sdss\_psfmag\_i}\textless21.0~AB, that lie within
1.49~degrees of at least one AQMES-medium field location.


\noindent\textbf{Target priority options:} 3100

\noindent\textbf{Cadence options:} \texttt{dark\_10x4\_4yr}

\noindent\textbf{Implementation:}
\href{https://github.com/sdss/target_selection/blob/0.3.0/python/target_selection/cartons/bhm_aqmes.py}{bhm\_aqmes.py}

\noindent\textbf{Number of targets:} 16853

\begin{center}\rule{0.5\linewidth}{0.5pt}\end{center}

\hypertarget{bhm_aqmes_wide2_plan0.5.4}{%
\subsection{bhm\_aqmes\_wide2}\label{bhm_aqmes_wide2_plan0.5.4}}

\noindent\textbf{target\_selection plan:} 0.5.4

\noindent\textbf{target\_selection tag:}
\href{https://github.com/sdss/target_selection/tree/0.3.5/}{0.3.5}

\noindent\textbf{crossmatch plan:} 0.5.0

\noindent\textbf{Summary:} Spectroscopically confirmed optically bright SDSS
QSOs, selected from the SDSS QSO catalog (DR16Q,
\citealt{Lyke2020}). Located in 425 fields within the SDSS QSO footprint,
where the choice of survey area prioritized field that overlapped with
the SPIDERS footprint (approx 180\textless{}\emph{b}\textless360~deg),
and/or had higher than average numbers of bright QSOs and CSC targets.
The list of field centres can be found
\href{https://github.com/sdss/target_selection/blob/0.3.0/python/target_selection/masks/candidate_target_fields_bhm_aqmes_wide_v0.3.1.fits}{within
the target\_selection repository}.

\noindent\textbf{Simplified description of selection criteria:} Select all
objects from SDSS DR16 QSO catalog that have
16.0\textless{}\texttt{sdss\_psfmag\_i}\textless19.1~AB, and that lie
within 1.49~degrees of at least one AQMES-wide field location.


\noindent\textbf{Target priority options:} 1210, 1211

\noindent\textbf{Cadence options:} \texttt{dark\_2x4}

\noindent\textbf{Implementation:}
\href{https://github.com/sdss/target_selection/blob/0.3.5/python/target_selection/cartons/bhm_aqmes.py}{bhm\_aqmes.py}

\noindent\textbf{Number of targets:} 24142

\begin{center}\rule{0.5\linewidth}{0.5pt}\end{center}

\hypertarget{bhm_aqmes_wide2_faint_plan0.5.4}{%
\subsection{bhm\_aqmes\_wide2\_faint}\label{bhm_aqmes_wide2_faint_plan0.5.4}}

\noindent\textbf{target\_selection plan:} 0.5.4

\noindent\textbf{target\_selection tag:}
\href{https://github.com/sdss/target_selection/tree/0.3.5/}{0.3.5}

\noindent\textbf{crossmatch plan:} 0.5.0

\noindent\textbf{Summary:} Spectroscopically confirmed optically faint SDSS QSOs,
selected from the SDSS QSO catalog (DR16Q,
\citealt{Lyke2020}). Located in 425 fields within the SDSS QSO footprint,
where the choice of survey area prioritized field that overlapped with
the SPIDERS footprint (approx 180\textless{}\emph{b}\textless360~deg),
and/or had higher than average numbers of bright QSOs and CSC targets.
The list of field centres can be found
\href{https://github.com/sdss/target_selection/blob/0.3.0/python/target_selection/masks/candidate_target_fields_bhm_aqmes_wide_v0.3.1.fits}{within
the target\_selection repository}.

\noindent\textbf{Simplified description of selection criteria:} Select all
objects from SDSS DR16 QSO catalog that have
19.1\textless{}\texttt{sdss\_psfmag\_i}\textless21.0~AB, and that lie
within 1.49~degrees of at least one AQMES-wide field location.


\noindent\textbf{Target priority options:} 3210, 3211

\noindent\textbf{Cadence options:} \texttt{dark\_2x4}

\noindent\textbf{Implementation:}
\href{https://github.com/sdss/target_selection/blob/0.3.5/python/target_selection/cartons/bhm_aqmes.py}{bhm\_aqmes.py}

\noindent\textbf{Number of targets:} 99586

\begin{center}\rule{0.5\linewidth}{0.5pt}\end{center}

\hypertarget{bhm_aqmes_bonus_core_plan0.5.4}{%
\subsection{bhm\_aqmes\_bonus\_core}\label{bhm_aqmes_bonus_core_plan0.5.4}}

\noindent\textbf{target\_selection plan:} 0.5.4

\noindent\textbf{target\_selection tag:}
\href{https://github.com/sdss/target_selection/tree/0.3.5/}{0.3.5}

\noindent\textbf{crossmatch plan:} 0.5.0

\noindent\textbf{Summary:} Spectroscopically confirmed optically bright SDSS
QSOs, selected from the SDSS QSO catalog (DR16Q,
\citealt{Lyke2020}). Located anywhere within the SDSS DR16Q footprint.

\noindent\textbf{Simplified description of selection criteria:} Select all
objects from SDSS DR16 QSO catalog that have
16.0\textless{}\texttt{sdss\_psfmag\_i}\textless19.1~AB


\noindent\textbf{Target priority options:} 3300,3301

\noindent\textbf{Cadence options:} \texttt{dark\_1x4}

\noindent\textbf{Implementation:}
\href{https://github.com/sdss/target_selection/blob/0.3.5/python/target_selection/cartons/bhm_aqmes.py}{bhm\_aqmes.py}

\noindent\textbf{Number of targets:} 83163

\begin{center}\rule{0.5\linewidth}{0.5pt}\end{center}

\hypertarget{bhm_aqmes_bonus_faint_plan0.5.4}{%
\subsection{bhm\_aqmes\_bonus\_faint}\label{bhm_aqmes_bonus_faint_plan0.5.4}}

\noindent\textbf{target\_selection plan:} 0.5.4

\noindent\textbf{target\_selection tag:}
\href{https://github.com/sdss/target_selection/tree/0.3.5/}{0.3.5}

\noindent\textbf{crossmatch plan:} 0.5.0

\noindent\textbf{Summary:} Spectroscopically confirmed optically faint SDSS QSOs,
selected from the SDSS QSO catalog (DR16Q,
\citealt{Lyke2020}). Located anywhere within the SDSS DR16Q footprint.

\noindent\textbf{Simplified description of selection criteria:} Select all
objects from SDSS DR16 QSO catalog that have
19.1\textless{}\texttt{sdss\_psfmag\_i}\textless21.0~AB


\noindent\textbf{Target priority options:} 3302,3303

\noindent\textbf{Cadence options:} \texttt{dark\_1x4}

\noindent\textbf{Implementation:}
\href{https://github.com/sdss/target_selection/blob/0.3.5/python/target_selection/cartons/bhm_aqmes.py}{bhm\_aqmes.py}

\noindent\textbf{Number of targets:} 424163

\begin{center}\rule{0.5\linewidth}{0.5pt}\end{center}

\hypertarget{bhm_aqmes_bonus_bright_plan0.5.4}{%
\subsection{bhm\_aqmes\_bonus\_bright}\label{bhm_aqmes_bonus_bright_plan0.5.4}}

\noindent\textbf{target\_selection plan:} 0.5.4

\noindent\textbf{target\_selection tag:}
\href{https://github.com/sdss/target_selection/tree/0.3.5/}{0.3.5}

\noindent\textbf{crossmatch plan:} 0.5.0

\noindent\textbf{Summary:} Spectroscopically confirmed, extremely optically
bright SDSS QSOs, selected from the SDSS QSO catalog (DR16Q,
\citealt{Lyke2020}). Located anywhere within the SDSS DR16Q footprint.

\noindent\textbf{Simplified description of selection criteria:} Select all
objects from SDSS DR16 QSO catalog that have
14.0\textless{}\texttt{sdss\_psfmag\_i}\textless18.0~AB


\noindent\textbf{Target priority options:} 4040,4041

\noindent\textbf{Cadence options:} \texttt{bright\_3x1}

\noindent\textbf{Implementation:}
\href{https://github.com/sdss/target_selection/blob/0.3.5/python/target_selection/cartons/bhm_aqmes.py}{bhm\_aqmes.py}

\noindent\textbf{Number of targets:} 10848

\begin{center}\rule{0.5\linewidth}{0.5pt}\end{center}

\hypertarget{bhm_rm_ancillary_plan0.5.0}{%
\subsection{bhm\_rm\_ancillary}\label{bhm_rm_ancillary_plan0.5.0}}

\noindent\textbf{target\_selection plan:} 0.5.0

\noindent\textbf{target\_selection tag:}
\href{https://github.com/sdss/target_selection/tree/0.3.0/}{0.3.0}

\noindent\textbf{crossmatch plan:} 0.5.0

\noindent\textbf{Summary:} A supporting sample of candidate QSOs which have been
selected by the Gaia-unWISE AGN catalog
(\citealt{Shu2019}) and/or the SDSS XDQSO catalog
(\citealt{Bovy2011}). These targets are located within five (+1 backup) well
known survey fields (SDSS-RM, COSMOS, XMM-LSS, ECDFS, CVZ-S/SEP, and
ELIAS-S1).

\noindent\textbf{Simplified description of selection criteria:} Starting from a
parent catalog of optically selected objects in the RM fields (as
presented by
\citealt{Yang2022}), select candidate QSOs that satisfy all of the
following: i) are identified via external ancillary methods
(\texttt{photo\_bitmask\ \&\ 3\ !=\ 0}); ii) have
15\textless{}\texttt{psfmag\_i}\textless21.5~AB
(16\textless{}\emph{G}\textless21.7~AB in the CVZ-S/SEP field); iii) do
not have significant detections (\textgreater3$\sigma$) of parallax and/or
proper motion in Gaia DR2; iv) are not vetoed due to results of visual
inspections of recent spectroscopy; and v) do not lie in the SDSS-RM
field


\noindent\textbf{Target priority options:} 900-1050

\noindent\textbf{Cadence options:} \texttt{dark\_174x8,\ dark\_100x8}

\noindent\textbf{Implementation:}
\href{https://github.com/sdss/target_selection/blob/0.3.0/python/target_selection/cartons/bhm_rm.py}{bhm\_rm.py}

\noindent\textbf{Number of targets:} 943

\begin{center}\rule{0.5\linewidth}{0.5pt}\end{center}

\hypertarget{bhm_rm_core_plan0.5.0}{%
\subsection{bhm\_rm\_core}\label{bhm_rm_core_plan0.5.0}}

\noindent\textbf{target\_selection plan:} 0.5.0

\noindent\textbf{target\_selection tag:}
\href{https://github.com/sdss/target_selection/tree/0.3.0/}{0.3.0}

\noindent\textbf{crossmatch plan:} 0.5.0

\noindent\textbf{Summary:} A sample of candidate QSOs selected via the methods
presented by
\citet{Yang2022}. These targets are located within five (+1 backup) well
known survey fields (SDSS-RM, COSMOS, XMM-LSS, ECDFS, CVZ-S/SEP, and
ELIAS-S1).

\noindent\textbf{Simplified description of selection criteria:} Starting from a
parent catalog of optically selected objects in the RM fields (as
presented by
\citealt{Yang2022}), select candidate QSOs that satisfy all of the
following: i) are identified via the Skew-T algorithm
(\texttt{skewt\_qso\ ==\ 1}); ii) have
17\textless{}\texttt{psfmag\_i}\textless21.5~AB
(16\textless{}\emph{G}\textless21.7~AB in the CVZ-S/SEP field); iii) do
not have significant detections (\textgreater3$\sigma$) of parallax and/or
proper motion in Gaia DR2; iv) are not vetoed due to results of visual
inspections of recent spectroscopy; vi) have detections in all of the
\emph{gri} bands (a Gaia detection is sufficient in the CVZ-S/SEP
field);and vii) do not lie in the SDSS-RM field


\noindent\textbf{Target priority options:} 900-1050

\noindent\textbf{Cadence options:} \texttt{dark\_174x8,\ dark\_100x8}

\noindent\textbf{Implementation:}
\href{https://github.com/sdss/target_selection/blob/0.3.0/python/target_selection/cartons/bhm_rm.py}{bhm\_rm.py}

\noindent\textbf{Number of targets:} 3721

\begin{center}\rule{0.5\linewidth}{0.5pt}\end{center}

\hypertarget{bhm_rm_var_plan0.5.0}{%
\subsection{bhm\_rm\_var}\label{bhm_rm_var_plan0.5.0}}

\noindent\textbf{target\_selection plan:} 0.5.0

\noindent\textbf{target\_selection tag:}
\href{https://github.com/sdss/target_selection/tree/0.3.0/}{0.3.0}

\noindent\textbf{crossmatch plan:} 0.5.0

\noindent\textbf{Summary:} A sample of candidate QSOs selected via their optical
variability properties, as presented by
\citet{Yang2022}. These targets are located within five (+1 backup) well
known survey fields (SDSS-RM, COSMOS, XMM-LSS, ECDFS, CVZ-S/SEP, and
ELIAS-S1).

\noindent\textbf{Simplified description of selection criteria:} Starting from a
parent catalog of optically selected objects in the RM fields (as
presented by
\citealt{Yang2022}), select candidate QSOs that satisfy all of the
following: i) have significant variability in the DES or PanSTARRS1
multi-epoch photometry (\texttt{var\_sn{[}g{]}}\textgreater3 and
\texttt{var\_rms{[}g{]}}\textgreater0.05); ii) have
17\textless{}\texttt{psfmag\_i}\textless20.5~AB
(16\textless{}\emph{G}\textless21.7~AB in the CVZ-S/SEP field); iii) do
not have significant detections (\textgreater3$\sigma$) of parallax and/or
proper motion in Gaia DR2; iv) are not vetoed due to results of visual
inspections of recent spectroscopy; and vi) do not lie in the SDSS-RM
field


\noindent\textbf{Target priority options:} 900-1050

\noindent\textbf{Cadence options:} \texttt{dark\_174x8,\ dark\_100x8}

\noindent\textbf{Implementation:}
\href{https://github.com/sdss/target_selection/blob/0.3.0/python/target_selection/cartons/bhm_rm.py}{bhm\_rm.py}

\noindent\textbf{Number of targets:} 934

\begin{center}\rule{0.5\linewidth}{0.5pt}\end{center}

\hypertarget{bhm_rm_known_spec_plan0.5.0}{%
\subsection{bhm\_rm\_known\_spec}\label{bhm_rm_known_spec_plan0.5.0}}

\noindent\textbf{target\_selection plan:} 0.5.0

\noindent\textbf{target\_selection tag:}
\href{https://github.com/sdss/target_selection/tree/0.3.0/}{0.3.0}

\noindent\textbf{crossmatch plan:} 0.5.0

\noindent\textbf{Summary:} A sample of known QSOs identified through optical
spectroscopy from various projects, as collated by
\citet{Yang2022}. These targets are located within five (+1 backup) well
known survey fields (SDSS-RM, COSMOS, XMM-LSS, ECDFS, CVZ-S/SEP, and
ELIAS-S1).

\noindent\textbf{Simplified description of selection criteria:} Starting from a
parent catalog of optically selected objects in the RM fields (as
presented by
\citealt{Yang2022}), select targets which satisfy all of the following: i)
are flagged as having a spectroscopic identification (in the parent
catalog or in the bhm\_rm\_tweaks table); ii) have
15\textless{}\texttt{psfmag\_i}\textless21.7~AB (SDSS-RM, CDFS, ELIAS-S1
field), 16\textless{}\emph{G}\textless21.7~Vega (CVZ-S/SEP field),
15\textless{}\texttt{psfmag\_i}\textless21.5~AB (COSMOS and XMM-LSS
fields); iii) have a spectroscopic redshift in the range
0.005\textless{}\emph{z}\textless7; iv) are not vetoed due to results of
visual inspections of recent spectroscopy


\noindent\textbf{Target priority options:} 900-1050

\noindent\textbf{Cadence options:} \texttt{dark\_174x8,\ dark\_100x8}

\noindent\textbf{Implementation:}
\href{https://github.com/sdss/target_selection/blob/0.3.0/python/target_selection/cartons/bhm_rm.py}{bhm\_rm.py}

\noindent\textbf{Number of targets:} 3022

\begin{center}\rule{0.5\linewidth}{0.5pt}\end{center}

\hypertarget{bhm_csc_apogee_plan0.5.15}{%
\subsection{bhm\_csc\_apogee}\label{bhm_csc_apogee_plan0.5.15}}

\noindent\textbf{target\_selection plan:} 0.5.15

\noindent\textbf{target\_selection tag:}
\href{https://github.com/sdss/target_selection/tree/0.3.14/}{0.3.14}

\noindent\textbf{crossmatch plan:} 0.5.0

\noindent\textbf{Summary:} X-ray sources from the CSC2.0 source catalog with NIR
counterparts in 2MASS PSC

\noindent\textbf{Simplified description of selection criteria:} Starting from the
parent catalog of CSC sources with optical/IR counterparts
(\texttt{bhm\_csc\_v2}). Select entries satisfying the following
criteria: i) NIR counterpart is from the 2MASS catalog, ii) 2MASS
\emph{H}-band magnitude measurement is not null and in the accepted
range for SDSS-V: 7.0\textless{}\emph{H\&}lt;14.0. Allocate cadence
(exposure time) requests based on \emph{H}-band magnitude.


\noindent\textbf{Target priority options:} 2930-2939

\noindent\textbf{Cadence options:} \texttt{bright\_1x1,bright\_3x1}

\noindent\textbf{Implementation:}
\href{https://github.com/sdss/target_selection/blob/0.3.14/python/target_selection/cartons/bhm_csc.py}{bhm\_csc.py}

\noindent\textbf{Number of targets:} 48928

\begin{center}\rule{0.5\linewidth}{0.5pt}\end{center}

\hypertarget{bhm_csc_boss_plan0.5.15}{%
\subsection{bhm\_csc\_boss}\label{bhm_csc_boss_plan0.5.15}}

\noindent\textbf{target\_selection plan:} 0.5.15

\noindent\textbf{target\_selection tag:}
\href{https://github.com/sdss/target_selection/tree/0.3.14/}{0.3.14}

\noindent\textbf{crossmatch plan:} 0.5.0

\noindent\textbf{Summary:} X-ray sources from the CSC2.0 source catalog with
counterparts in Panstarrs1-DR1 or Gaia DR2.

\noindent\textbf{Simplified description of selection criteria:} Starting from the
parent catalog of CSC sources with optical/IR counterparts
(\texttt{bhm\_csc\_v2}). Select entries satisfying the following
criteria: i) optical counterpart is from the PanSTARRS1 or Gaia DR2
catalogs, ii) optical flux/magnitude is in the accepted range for
SDSS-V: \texttt{psfmag\_{[}g,r,i,z{]}}\textgreater13.5~AB and non-Null
\texttt{psfmag\_i} (objects with PanSTARRS1 counterparts);
\emph{G,RP}\textgreater13.0~Vega (Gaia DR2 counterparts). Deprioritize
targets which already have good quality SDSS spectroscopy. Allocate
cadence (exposure time) requests based on optical brightness (PS1
\texttt{psfmag\_i} or Gaia \emph{G}).


\noindent\textbf{Target priority options:} 1920-1939, 2920-2939

\noindent\textbf{Cadence options:} \texttt{bright\_1x1,dark\_1x2,dark\_1x4}

\noindent\textbf{Implementation:}
\href{https://github.com/sdss/target_selection/blob/0.3.14/python/target_selection/cartons/bhm_csc.py}{bhm\_csc.py}

\noindent\textbf{Number of targets:} 122731

\begin{center}\rule{0.5\linewidth}{0.5pt}\end{center}

\hypertarget{bhm_gua_bright_plan0.5.0}{%
\subsection{bhm\_gua\_bright}\label{bhm_gua_bright_plan0.5.0}}

\noindent\textbf{target\_selection plan:} 0.5.0

\noindent\textbf{target\_selection tag:}
\href{https://github.com/sdss/target_selection/tree/0.3.0/}{0.3.0}

\noindent\textbf{crossmatch plan:} 0.5.0

\noindent\textbf{Summary:} A sample of optically bright candidate AGN lacking
spectroscopic confirmations, derived from the parent sample presented by
\citet{Shu2019}, who applied a machine-learning approach to select QSO
candidates from a combination of the Gaia DR2 and unWISE catalogs.

\noindent\textbf{Simplified description of selection criteria:} Starting with the
\citet{Shu2019} catalog, select targets which satisfy the following
criteria: i) have a Random Forest probability of being a QSO
of\textgreater0.8, ii) are in the (dereddened) magnitude range suitable
for BOSS spectroscopy in bright time
(\emph{G}\textsubscript{dered}\textgreater13.0 and
\emph{RP}\textsubscript{dered}\textgreater13.5, as well as
\emph{G}\textsubscript{dered}\textless18.5 or
\emph{RP}\textsubscript{dered}\textless18.5,~Vega mags), iii) do not
have good optical spectroscopic measurements from a previous iteration
of SDSS.


\noindent\textbf{Target priority options:} 4040

\noindent\textbf{Cadence options:} \texttt{bright\_2x1}

\noindent\textbf{Implementation:}
\href{https://github.com/sdss/target_selection/blob/0.3.0/python/target_selection/cartons/bhm_gua.py}{bhm\_gua.py}

\noindent\textbf{Number of targets:} 254601

\begin{center}\rule{0.5\linewidth}{0.5pt}\end{center}

\hypertarget{bhm_gua_dark_plan0.5.0}{%
\subsection{bhm\_gua\_dark}\label{bhm_gua_dark_plan0.5.0}}

\noindent\textbf{target\_selection plan:} 0.5.0

\noindent\textbf{target\_selection tag:}
\href{https://github.com/sdss/target_selection/tree/0.3.0/}{0.3.0}

\noindent\textbf{crossmatch plan:} 0.5.0

\noindent\textbf{Summary:} A sample of optically faint candidate AGN lacking
spectroscopic confirmations, derived from the parent sample presented by
\citet{Shu2019}, who applied a machine-learning approach to select QSO
candidates from a combination of the Gaia DR2 and unWISE catalogs.

\noindent\textbf{Simplified description of selection criteria:} Starting with the
\citet{Shu2019} catalog, select targets which satisfy the following
criteria: i) have a Random Forest probability of being a QSO
of\textgreater0.8, ii) are in the (dereddened) magnitude range suitable
for BOSS spectroscopy in dark time
(\emph{G}\textsubscript{dered}\textgreater16.5 and
\emph{RP}\textsubscript{dered}\textgreater16.5, as well as
\emph{G}\textsubscript{dered}\textless21.2 or
\emph{RP}\textsubscript{dered}\textless21.0,~Vega mag), iii) do not have
good optical spectroscopic measurements from a previous iteration of
SDSS.


\noindent\textbf{Target priority options:} 3400

\noindent\textbf{Cadence options:} \texttt{dark\_1x4}

\noindent\textbf{Implementation:}
\href{https://github.com/sdss/target_selection/blob/0.3.0/python/target_selection/cartons/bhm_gua.py}{bhm\_gua.py}

\noindent\textbf{Number of targets:} 2156582

\begin{center}\rule{0.5\linewidth}{0.5pt}\end{center}

\hypertarget{bhm_colr_galaxies_lsdr8_plan0.5.16}{%
\subsection{bhm\_colr\_galaxies\_lsdr8}\label{bhm_colr_galaxies_lsdr8_plan0.5.16}}

\noindent\textbf{target\_selection plan:} 0.5.16

\noindent\textbf{target\_selection tag:}
\href{https://github.com/sdss/target_selection/tree/0.3.13/}{0.3.13}

\noindent\textbf{crossmatch plan:} 0.5.0

\noindent\textbf{Summary:} A supplementary magnitude limited sample of optically
bright galaxies selected from the legacysurvey.org/dr8 optical/IR
imaging catalog. Selection is based on optical morphology, lack of Gaia
DR2 parallax, and several magnitude cuts.

\noindent\textbf{Simplified description of selection criteria:} Starting from the
\texttt{legacy\_survey\_dr8} catalog (lsdr8), select entries satisfying
all of the following criteria: i) lsdr8 morphological \texttt{type} !=
'PSF', ii) zero or Null parallax in Gaia DR2, iii) dereddened
\emph{z}-band model mag\textless19.0~AB, and dereddened \emph{z}-band
fiber mag\textless19.5~AB, and 16\textless apparent \emph{z}-band fiber
mag\textless19.0~AB, and apparent \emph{r}-band fiber
mag\textgreater16.0~AB, and Gaia \emph{G}\textgreater15.0~Vega, and Gaia
\emph{RP}\textgreater15.0~Vega;


\noindent\textbf{Target priority options:} 7100

\noindent\textbf{Cadence options:} \texttt{bright\_1x1,\ dark\_1x1,\ dark\_1x4}

\noindent\textbf{Implementation:}
\href{https://github.com/sdss/target_selection/blob/0.3.13/python/target_selection/cartons/bhm_galaxies.py}{bhm\_galaxies.py}

\noindent\textbf{Number of targets:} 7320203

\begin{center}\rule{0.5\linewidth}{0.5pt}\end{center}

\hypertarget{bhm_spiders_agn_lsdr8_plan0.5.0}{%
\subsection{bhm\_spiders\_agn\_lsdr8}\label{bhm_spiders_agn_lsdr8_plan0.5.0}}

\noindent\textbf{target\_selection plan:} 0.5.0

\noindent\textbf{target\_selection tag:}
\href{https://github.com/sdss/target_selection/tree/0.3.0/}{0.3.0}

\noindent\textbf{crossmatch plan:} 0.5.0

\noindent\textbf{Summary:} This is the highest priority carton for SPIDERS AGN
wide area follow up. The carton provides optical counterparts to
point-like (unresolved) X-ray sources detected in early reductions of
the first 6~months of eROSITA all sky survey data (eRASS:1). The sample
is expected to contain a mixture of QSOs, AGN, stars and compact
objects. The X-ray sources have been cross-matched by the eROSITA-DE
team to \href{https://www.legacysurvey.org/dr8/}{legacysurvey.org/dr8}
optical/IR counterparts. All targets are located in the sky hemisphere
where MPE controls the data rights (approx.
180\textless{}\emph{l}\textless360~deg). Due to the footprint of lsdr8,
nearly all targets in this carton are located at high Galactic latitudes
\textbar{}\emph{b}\textbar\textgreater15~deg.

\noindent\textbf{Simplified description of selection criteria:} Starting from a
parent catalog of eRASS:1 point source $\rightarrow$ legacysurvey.org/dr8
associations (method: NWAY assisted by optical/IR priors computed via a
pre-trained Random Forest, building on
\citealt{Salvato2022}), select targets which meet all of the following criteria:
i) have eROSITA detection likelihood\textgreater6.0, ii) have an X-ray $\rightarrow$
optical/IR cross-match probability (NWAY) of
\texttt{p\_any}\textgreater0.1, iii) have
13.5\textless{}\texttt{fibertotmag\_r}\textless22.5 or
13.5\textless{}\texttt{fibertotmag\_z}\textless21.0, iv) are not
saturated in LegacySurvey imaging, v) have at least one observation in
\emph{r}-band and at least one observation in \emph{g}- or
\emph{z}-band, vi) if detected by Gaia DR2 then have
\emph{G}\textgreater13.5 and \emph{RP}\textgreater13.5~Vega. We
deprioritise targets if any of the following criteria are met: a) the
target already has existing good quality SDSS spectroscopy, b) the X-ray
detection likelihood is \textless8.0, c) the target is a secondary
X-ray$\rightarrow$optical/IR association. We assign cadences (exposure time
requests) based on optical brightness.


\noindent\textbf{Target priority options:} 1520-1523, 1720-1723

\noindent\textbf{Cadence options:} \texttt{bright\_2x1,\ dark\_1x2,\ dark\_1x4}

\noindent\textbf{Implementation:}
\href{https://github.com/sdss/target_selection/blob/0.3.0/python/target_selection/cartons/bhm_spiders_agn.py}{bhm\_spiders\_agn.py}

\noindent\textbf{Number of targets:} 235745

\begin{center}\rule{0.5\linewidth}{0.5pt}\end{center}

\hypertarget{bhm_spiders_agn_ps1dr2_plan0.5.0}{%
\subsection{bhm\_spiders\_agn\_ps1dr2}\label{bhm_spiders_agn_ps1dr2_plan0.5.0}}

\noindent\textbf{target\_selection plan:} 0.5.0

\noindent\textbf{target\_selection tag:}
\href{https://github.com/sdss/target_selection/tree/0.3.0/}{0.3.0}

\noindent\textbf{crossmatch plan:} 0.5.0

\noindent\textbf{Summary:} This is the second highest priority carton for SPIDERS
AGN wide area follow up, it is included to expand the survey footprint
beyond legacysurvey/dr8. The carton provides optical counterparts to
point-like (unresolved) X-ray sources detected in early reductions of
the first 6~months of eROSITA all sky survey data (eRASS:1). The sample
is expected to contain a mixture of QSOs, AGN, stars and compact
objects. The X-ray sources have been cross-matched by the eROSITA-DE
team, first to
CatWISE2020 \citep{Marocco2021}
mid-IR sources, and then to optical counterparts from the
\href{https://outerspace.stsci.edu/display/PANSTARRS/}{PanSTARRS1 dr2}
catalog. All targets are located in the sky hemisphere where MPE
controls the data rights (approx.
180\textless{}\emph{l}\textless360~deg), and at Dec\textgreater-30~deg,
spanning a wide range of Galactic latitudes. Targets at low Galactic
latitudes \textbar{}\emph{b}\textbar\textless15~deg do not drive survey
strategy.

\noindent\textbf{Simplified description of selection criteria:} Starting from a
parent catalog of eRASS:1 point source $\rightarrow$ CatWISE2020 $\rightarrow$ Pan-STARRS1
associations (method: NWAY assisted by IR priors computed via a
pre-trained Random Forest, building on
\citealt{Salvato2022}), select targets which meet all of the following criteria:
i) have eROSITA detection likelihood\textgreater6.0, ii) have an
X-ray$\rightarrow$IR cross-match probability of \texttt{p\_any}\textgreater0.1, iii)
have \texttt{psfmag\_g}, \texttt{psfmag\_r}, \texttt{psfmag\_i},
\texttt{psfmag\_z}\textgreater13.5~AB and at least one of
\texttt{psfmag\_g}\textless22.5, \texttt{psfmag\_r}\textless22.0,
\texttt{psfmag\_i}\textless21.5 or \texttt{psfmag\_z}\textless20.5, iv)
are not associated with a bad Pan-STARRS1-dr2 image stack, v) have
non-null measurements of \texttt{psfmag\_g}, \texttt{psfmag\_r} and
\texttt{psfmag\_i}, vi) if detected by Gaia DR2 then have
\emph{G}\textgreater13.5 and \emph{RP}\textgreater13.5~Vega. We
deprioritise targets if any of the following criteria are met: a) the
target already has existing good quality SDSS spectroscopy, b) the X-ray
detection likelihood is \textless8.0, c) the target is a secondary
X-ray$\rightarrow$IR association. We assign cadences (exposure time requests) based
on optical brightness.


\noindent\textbf{Target priority options:}
1530-1533,1730-1733,3530-3533,3730-3732

\noindent\textbf{Cadence options:} \texttt{bright\_2x1,\ dark\_1x2,\ dark\_1x4}

\noindent\textbf{Implementation:}
\href{https://github.com/sdss/target_selection/blob/0.3.0/python/target_selection/cartons/bhm_spiders_agn.py}{bhm\_spiders\_agn.py}

\noindent\textbf{Number of targets:} 200681

\begin{center}\rule{0.5\linewidth}{0.5pt}\end{center}

\hypertarget{bhm_spiders_agn_gaiadr2_plan0.5.0}{%
\subsection{bhm\_spiders\_agn\_gaiadr2}\label{bhm_spiders_agn_gaiadr2_plan0.5.0}}

\noindent\textbf{target\_selection plan:} 0.5.0

\noindent\textbf{target\_selection tag:}
\href{https://github.com/sdss/target_selection/tree/0.3.0/}{0.3.0}

\noindent\textbf{crossmatch plan:} 0.5.0

\noindent\textbf{Summary:} This is the third highest priority carton for SPIDERS
AGN wide area follow up, it is included to expand the survey footprint
to the full hemisphere where X-ray sources are available (beyond
legacysurvey/dr8 and Pan-STARRS1). The carton provides optical
counterparts to point-like (unresolved) X-ray sources detected in early
reductions of the first 6~months of eROSITA all sky survey data
(eRASS:1). The sample is expected to contain a mixture of QSOs, AGN,
stars and compact objects. The X-ray sources have been cross-matched by
the eROSITA-DE team, first to
CatWISE2020 \citep{Marocco2021}
mid-IR sources, and then to optical counterparts from the Gaia DR2
catalog. All targets are located in the sky hemisphere where MPE
controls the data rights (approx.
180\textless{}\emph{l}\textless360~deg). The targets in this carton are
distributed over a wide range of Galactic latitudes, but targets at low
Galactic latitudes \textbar{}\emph{b}\textbar\textless15~deg do not
drive survey strategy.

\noindent\textbf{Simplified description of selection criteria:} Starting from a
parent catalog of eRASS:1 point source $\rightarrow$ CatWISE2020 $\rightarrow$ Gaia DR2
associations (method: NWAY assisted by IR priors computed via a
pre-trained Random Forest, building on
\citealt{Salvato2022}), select targets which meet all of the following criteria:
i) have eROSITA detection likelihood\textgreater6.0, ii) have an
X-ray$\rightarrow$IR cross-match probability of \texttt{p\_any}\textgreater0.1, iii)
have \emph{G}\textgreater13.5 and \emph{RP}\textgreater13.5~Vega. We
deprioritise targets if any of the following criteria are met: a) the
target already has existing good quality SDSS spectroscopy, b) the X-ray
detection likelihood is \textless8.0, c) the target is a secondary
X-ray$\rightarrow$IR association. We assign cadences (exposure time requests) based
on optical brightness.


\noindent\textbf{Target priority options:}
1540-1543,1740-1743,3540-3543,3740-3742

\noindent\textbf{Cadence options:} \texttt{bright\_2x1,\ dark\_1x2,\ dark\_1x4}

\noindent\textbf{Implementation:}
\href{https://github.com/sdss/target_selection/blob/0.3.0/python/target_selection/cartons/bhm_spiders_agn.py}{bhm\_spiders\_agn.py}

\noindent\textbf{Number of targets:} 324576

\begin{center}\rule{0.5\linewidth}{0.5pt}\end{center}

\hypertarget{bhm_spiders_agn_skymapperdr2_plan0.5.0}{%
\subsection{bhm\_spiders\_agn\_skymapperdr2}\label{bhm_spiders_agn_skymapperdr2_plan0.5.0}}

\noindent\textbf{target\_selection plan:} 0.5.0

\noindent\textbf{target\_selection tag:}
\href{https://github.com/sdss/target_selection/tree/0.3.0/}{0.3.0}

\noindent\textbf{crossmatch plan:} 0.5.0

\noindent\textbf{Summary:} This is a lower ranked carton for SPIDERS AGN wide
area follow up, it supplements the survey in areas which rely on Gaia
DR2 (beyond LegacySurvey DR8 and Pan-STARRS1) by recovering extended
targets (galaxies) that are missed by Gaia. The carton provides optical
counterparts to point-like (unresolved) X-ray sources detected in early
reductions of the first 6~months of eROSITA all sky survey data
(eRASS:1). The sample is expected to contain a mixture of QSOs, AGN,
stars and compact objects. The X-ray sources have been cross-matched by
the eROSITA-DE team, first to
CatWISE2020 \citep{Marocco2021}
mid-IR sources, and then to optical counterparts from the
SkyMapper-dr2 \citep{Onken2019}
catalog. All targets are located in the sky hemisphere where MPE
controls the data rights (approx.
180\textless{}\emph{l}\textless360~deg) and at Dec\textless0deg,
spanning a wide range of Galactic latitudes. Targets at low Galactic
latitudes \textbar{}\emph{b}\textbar\textless15~deg do not drive survey
strategy.

\noindent\textbf{Simplified description of selection criteria:} Starting from a
parent catalog of eRASS:1 point source $\rightarrow$ CatWISE2020 $\rightarrow$ SkyMapper-dr2
associations (method: NWAY assisted by IR priors computed via a
pre-trained Random Forest, building on
\citealt{Salvato2022}), select targets which meet all of the following criteria:
i) have eROSITA detection likelihood\textgreater6.0, ii) have an
X-ray$\rightarrow$IR cross-match probability of \texttt{p\_any}\textgreater0.1, iii)
have \texttt{psfmag\_g}, \texttt{psfmag\_r}, \texttt{psfmag\_i},
\texttt{psfmag\_z}\textgreater13.5~AB and at least one of
\texttt{psfmag\_g}\textless22.5, \texttt{psfmag\_r}\textless22.0,
\texttt{psfmag\_i}\textless21.5 or \texttt{psfmag\_z}\textless20.5, iv)
Are not associated with a bad SkyMapper source detection
(flags\textless4), v) have non-null measurements in at least one of
g\_psfmag, r\_psfmag and i\_psfmag, vi) if detected by Gaia DR2 then
have \emph{G}\textgreater13.5 and \emph{RP}\textgreater13.5~Vega. We
deprioritise targets if any of the following criteria are met: a) the
target already has existing good quality SDSS spectroscopy, b) the X-ray
detection likelihood is \textless8.0, c) the target is a secondary
X-ray$\rightarrow$IR association. We assign cadences (exposure time requests) based
on optical brightness.


\noindent\textbf{Target priority options:}
1550-1553,1750-1753,3550-3553,3750-3752

\noindent\textbf{Cadence options:} \texttt{bright\_2x1,\ dark\_1x2,\ dark\_1x4}

\noindent\textbf{Implementation:}
\href{https://github.com/sdss/target_selection/blob/0.3.0/python/target_selection/cartons/bhm_spiders_agn.py}{bhm\_spiders\_agn.py}

\noindent\textbf{Number of targets:} 82683

\begin{center}\rule{0.5\linewidth}{0.5pt}\end{center}

\hypertarget{bhm_spiders_agn_supercosmos_plan0.5.0}{%
\subsection{bhm\_spiders\_agn\_supercosmos}\label{bhm_spiders_agn_supercosmos_plan0.5.0}}

\noindent\textbf{target\_selection plan:} 0.5.0

\noindent\textbf{target\_selection tag:}
\href{https://github.com/sdss/target_selection/tree/0.3.0/}{0.3.0}

\noindent\textbf{crossmatch plan:} 0.5.0

\noindent\textbf{Summary:} This is a lower ranked carton for SPIDERS AGN wide
area follow up, it supplements the survey in areas which rely on
Gaia-DR2 (beyond LegacySurvey DR8 and Pan-STARRS1) by recovering
extended targets (galaxies) that are missed by Gaia. The carton provides
optical counterparts to point-like (unresolved) X-ray sources detected
in early reductions of the first 6~months of eROSITA all sky survey data
(eRASS:1). The sample is expected to contain a mixture of QSOs, AGN,
stars and compact objects. The X-ray sources have been cross-matched by
the eROSITA-DE team, first to
CatWISE2020 \citep{Marocco2021}
mid-IR sources, and then to optical counterparts from the
\href{http://www-wfau.roe.ac.uk/sss/intro.html}{SuperCosmos Sky Surveys}
catalog (derived from scans of photographic plates). All targets are
located in the sky hemisphere where MPE controls the data rights
(approx. 180\textless{}\emph{l}\textless360~deg), spanning a wide range
of Galactic latitudes. Targets at low Galactic latitudes
\textbar{}\emph{b}\textbar\textless15~deg do not drive survey strategy.

\noindent\textbf{Simplified description of selection criteria:} Starting from a
parent catalog of eRASS:1 point source $\rightarrow$ CatWISE2020 $\rightarrow$ SuperCosmos
associations (method: NWAY assisted by IR priors computed via a
pre-trained Random Forest, building on
\citealt{Salvato2022}), select targets which meet all of the following criteria:
i) have eROSITA detection likelihood\textgreater6.0, ii) have an
X-ray$\rightarrow$IR cross-match probability of \texttt{p\_any}\textgreater0.1, iii)
have \texttt{scormagb}, \texttt{scormagr2},
\texttt{scormagi}\textgreater13.5~Vega and at least one of
\texttt{scormagb}\textless22.5, \texttt{scormagr2}\textless22.0, or
\texttt{scormagi}\textless21.5, iv) if detected by Gaia DR2 then have
\emph{G}\textgreater13.5 and \emph{RP}\textgreater13.5~Vega. We
deprioritise targets if any of the following criteria are met: a) the
target already has existing good quality SDSS spectroscopy, b) the X-ray
detection likelihood is \textless8.0, c) the target is a secondary
X-ray$\rightarrow$IR association. We assign cadences (exposure time requests) based
on optical brightness.


\noindent\textbf{Target priority options:}
1560-1563,1760-1763,3560-3563,3760-3763

\noindent\textbf{Cadence options:} \texttt{bright\_2x1,\ dark\_1x2,\ dark\_1x4}

\noindent\textbf{Implementation:}
\href{https://github.com/sdss/target_selection/blob/0.3.0/python/target_selection/cartons/bhm_spiders_agn.py}{bhm\_spiders\_agn.py}

\noindent\textbf{Number of targets:} 430780

\begin{center}\rule{0.5\linewidth}{0.5pt}\end{center}

\hypertarget{bhm_spiders_agn_efeds_stragglers_plan0.5.0}{%
\subsection{bhm\_spiders\_agn\_efeds\_stragglers}\label{bhm_spiders_agn_efeds_stragglers_plan0.5.0}}

\noindent\textbf{target\_selection plan:} 0.5.0

\noindent\textbf{target\_selection tag:}
\href{https://github.com/sdss/target_selection/tree/0.3.0/}{0.3.0}

\noindent\textbf{crossmatch plan:} 0.5.0

\noindent\textbf{Summary:} This is an opportunistic supplementary carton for
SPIDERS AGN follow up, aiming, where fiber resources allow, to gather a
small additional number of spectra for targets in the eFEDS field (which
has been repeatedly surveyed in earlier iterations of SDSS). This carton
provides optical counterparts to point-like (unresolved) X-ray sources
detected in early reductions of the eROSITA/eFEDS performance validation
field. The sample is expected to contain a mixture of QSOs, AGN, stars
and compact objects. The X-ray sources have been cross-matched by the
eROSITA-DE team to
\href{https://www.legacysurvey.org/dr8/}{legacysurvey.org/dr8}
optical/IR counterparts. All targets in this carton are located within
the eFEDS field (approx 126\textless RA\textless146,
-3\textless Dec\textless+6~deg). These targets do not drive survey
strategy.

\noindent\textbf{Simplified description of selection criteria:} Starting from a
parent catalog of eFEDS point source $\rightarrow$ LegacySurvey DR8 associations
(method: NWAY assisted by optical/IR priors computed via a pre-trained
Random Forest, see
\citealt{Salvato2022}), select targets which meet all of the following criteria:
i) have eROSITA detection likelihood\textgreater6.0, ii) have an X-ray $\rightarrow$
optical/IR cross-match probability (NWAY) of
\texttt{p\_any}\textgreater0.1, iii) have
13.5\textless{}\texttt{fibertotmag\_r}\textless22.5 or
13.5\textless{}\texttt{fibertotmag\_z}\textless21.0, iv) are not
saturated in LegacySurvey imaging, v) have at least one observation in
r-band and at least one observation in g- or z-band, vi) if detected by
Gaia DR2 then have \emph{G}\textgreater13.5 and
\emph{RP}\textgreater13.5~Vega. We deprioritise targets if any of the
following criteria are met: a) the target already has existing good
quality SDSS spectroscopy, b) the X-ray detection likelihood is
\textless8.0, c) the target is a secondary X-ray$\rightarrow$optical/IR association.
We assign cadences (exposure time requests) based on optical brightness.


\noindent\textbf{Target priority options:} 1510-1514, 1710-1714

\noindent\textbf{Cadence options:} \texttt{bright\_2x1,\ dark\_1x2,\ dark\_1x4}

\noindent\textbf{Implementation:}
\href{https://github.com/sdss/target_selection/blob/0.3.0/python/target_selection/cartons/bhm_spiders_agn.py}{bhm\_spiders\_agn.py}

\noindent\textbf{Number of targets:} 15926

\begin{center}\rule{0.5\linewidth}{0.5pt}\end{center}

\hypertarget{bhm_spiders_agn_sep_plan0.5.0}{%
\subsection{bhm\_spiders\_agn\_sep}\label{bhm_spiders_agn_sep_plan0.5.0}}

\noindent\textbf{target\_selection plan:} 0.5.0

\noindent\textbf{target\_selection tag:}
\href{https://github.com/sdss/target_selection/tree/0.3.0/}{0.3.0}

\noindent\textbf{crossmatch plan:} 0.5.0

\noindent\textbf{Summary:} This special carton is dedicated to SPIDERS AGN follow
up in the CVZ-S/SEP field. The carton provides optical counterparts to
point-like (unresolved) X-ray sources detected in a dedicated analysis
of the first 6~months of eROSITA scanning data near the SEP. The X-ray
sources have been cross-matched by the eROSITA-DE team, first to
CatWISE2020 \citep{Marocco2021}
mid-IR sources, and then to optical counterparts from the Gaia DR2
catalog, using additional filtering (including Gaia EDR3 astrometric
info) to reduce the contamination from foreground stars located in the
LMC. All targets are located within 1.5~deg of the SEP (RA,Dec =
90.0,-66.56~deg).

\noindent\textbf{Simplified description of selection criteria:} Starting from a
parent catalog of eRASS:1/SEP point source $\rightarrow$ CatWISE2020 $\rightarrow$ Gaia DR2
associations (method: NWAY assisted by IR priors computed via a
pre-trained Random Forest, building on
\citealt{Salvato2022}), select targets which meet all of the following criteria:
i) have eROSITA detection likelihood\textgreater6.0, ii) have an
X-ray$\rightarrow$IR cross-match probability of \texttt{p\_any}\textgreater0.1, iii)
have \emph{G}\textgreater13.5 and \emph{RP}\textgreater13.5~AB. We
deprioritise targets if any of the following criteria are met: a) the
X-ray detection likelihood is \textless8.0, b) the target is a secondary
X-ray$\rightarrow$IR association. We assign cadences (exposure time requests) based
on optical brightness.


\noindent\textbf{Target priority options:} 1510,1512

\noindent\textbf{Cadence options:} \texttt{bright\_2x1,\ dark\_1x2,\ dark\_1x4}

\noindent\textbf{Implementation:}
\href{https://github.com/sdss/target_selection/blob/0.3.0/python/target_selection/cartons/bhm_spiders_agn.py}{bhm\_spiders\_agn.py}

\noindent\textbf{Number of targets:} 697

\begin{center}\rule{0.5\linewidth}{0.5pt}\end{center}

\hypertarget{bhm_spiders_clusters_lsdr8_plan0.5.0}{%
\subsection{bhm\_spiders\_clusters\_lsdr8}\label{bhm_spiders_clusters_lsdr8_plan0.5.0}}

\noindent\textbf{target\_selection plan:} 0.5.0

\noindent\textbf{target\_selection tag:}
\href{https://github.com/sdss/target_selection/tree/0.3.0/}{0.3.0}

\noindent\textbf{crossmatch plan:} 0.5.0

\noindent\textbf{Summary:} This is the highest priority carton for SPIDERS
Clusters wide area follow up. The carton provides a list of galaxies
which are candidate members of clusters selected from early reductions
of the first 6~months of eROSITA all sky survey data (eRASS:1). The
X-ray clusters have been associated by the eROSITA-DE team to
\href{https://www.legacysurvey.org/dr8/}{legacysurvey.org/dr8} (lsdr8)
optical/IR counterparts using the eROMAPPER red-sequence finder
algorithm
(\citealt{Rykoff2014};
\citealt{IderChitham2020}). All targets are located in the sky hemisphere
where MPE controls the data rights (approx.
180\textless{}\emph{l}\textless360~deg). Due to the footprint of lsdr8,
nearly all targets in this carton are located at high Galactic latitudes
\textbar{}\emph{b}\textbar\textgreater15~deg.

\noindent\textbf{Simplified description of selection criteria:} Starting from a
parent catalog of eRASS:1 $\rightarrow$ legacysurvey.org/dr8 eROMAPPER cluster
associations, select targets which meet all of the following criteria:
i) have 13.5\textless{}\texttt{fibertotmag\_r}\textless21.0 or
13.5\textless{}\texttt{fibertotmag\_z}\textless20.0, ii) if detected by
Gaia DR2 then have \emph{G}\textgreater13.5 and
\emph{RP}\textgreater13.5~Vega, iii) do not have existing good quality
SDSS spectroscopy. We assign a range of priorities to targets in this
carton, with BCGs top ranked, followed by candidate member galaxies
according their probability of membership. We assign cadences (exposure
time requests) based on optical brightness.


\noindent\textbf{Target priority options:} 1501,1630-1659

\noindent\textbf{Cadence options:} \texttt{bright\_2x1,\ dark\_1x2,\ dark\_1x4}

\noindent\textbf{Implementation:}
\href{https://github.com/sdss/target_selection/blob/0.3.0/python/target_selection/cartons/bhm_spiders_clusters.py}{bhm\_spiders\_clusters.py}

\noindent\textbf{Number of targets:} 87490

\begin{center}\rule{0.5\linewidth}{0.5pt}\end{center}

\hypertarget{bhm_spiders_clusters_ps1dr2_plan0.5.0}{%
\subsection{bhm\_spiders\_clusters\_ps1dr2}\label{bhm_spiders_clusters_ps1dr2_plan0.5.0}}

\noindent\textbf{target\_selection plan:} 0.5.0

\noindent\textbf{target\_selection tag:}
\href{https://github.com/sdss/target_selection/tree/0.3.0/}{0.3.0}

\noindent\textbf{crossmatch plan:} 0.5.0

\noindent\textbf{Summary:} This is the second highest priority carton for SPIDERS
Clusters wide area follow up, it is designed to expand the survey area
beyond the legacysurvey/dr8 footprint. The carton provides a list of
galaxies which are candidate members of clusters selected from early
reductions of the first 6~months of eROSITA all sky survey data
(eRASS:1). The X-ray clusters have been associated by the eROSITA-DE
team to the
\href{https://outerspace.stsci.edu/display/PANSTARRS/}{PanSTARRS1 dr2}
catalog using the eROMAPPER red-sequence finder algorithm
(\citealt{Rykoff2014};
\citealt{IderChitham2020}). All targets are located in the sky hemisphere
where MPE controls the data rights (approx.
180\textless{}\emph{l}\textless360~deg). Nearly all targets in this
carton are located at high Galactic latitudes
\textbar{}\emph{b}\textbar\textgreater15~deg.

\noindent\textbf{Simplified description of selection criteria:} Starting from a
parent catalog of eRASS:1 $\rightarrow$ Pan-STARRS1-dr2 eROMAPPER cluster
associations, select targets which meet all of the following criteria:
i) have \texttt{psfmag\_r}, \texttt{psfmag\_i},
\texttt{psfmag\_z}\textgreater13.5~AB and at least one of
\texttt{psfmag\_r}\textless21.5, \texttt{psfmag\_i}\textless21.0 or
\texttt{psfmag\_z}\textless20.5, ii) do not have existing good quality
SDSS spectroscopy. We assign a range of priorities to targets in this
carton, with BCGs top ranked, followed by candidate member galaxies
according their probability of membership. We assign cadences (exposure
time requests) based on optical brightness.


\noindent\textbf{Target priority options:} 1502,1660-1689

\noindent\textbf{Cadence options:} \texttt{bright\_2x1,dark\_1x2,dark\_1x4}

\noindent\textbf{Implementation:}
\href{https://github.com/sdss/target_selection/blob/0.3.0/python/target_selection/cartons/bhm_spiders_clusters.py}{bhm\_spiders\_clusters.py}

\noindent\textbf{Number of targets:} 86179

\begin{center}\rule{0.5\linewidth}{0.5pt}\end{center}

\hypertarget{bhm_spiders_clusters_efeds_stragglers_plan0.5.0}{%
\subsection{bhm\_spiders\_clusters\_efeds\_stragglers}\label{bhm_spiders_clusters_efeds_stragglers_plan0.5.0}}

\noindent\textbf{target\_selection plan:} 0.5.0

\noindent\textbf{target\_selection tag:}
\href{https://github.com/sdss/target_selection/tree/0.3.0/}{0.3.0}

\noindent\textbf{crossmatch plan:} 0.5.0

\noindent\textbf{Summary:} This is an opportunistic supplementary carton for
SPIDERS Clusters follow up, aiming, where fiber resources allow, to
gather a small additional number of spectra for targets in the eFEDS
field (which has been repeatedly surveyed in earlier iterations of
SDSS). The carton provides a list of galaxies which are candidate
members of clusters selected from early reductions of the eROSITA
performance verification survey in the eFEDS field. The X-ray clusters
have been associated by the eROSITA-DE team to
\href{https://www.legacysurvey.org/dr8/}{legacysurvey.org/dr8}
optical/IR counterparts. All targets in this carton are located within
the eFEDS field (approx 126\textless RA\textless146,
-3\textless Dec\textless+6~deg).

\noindent\textbf{Simplified description of selection criteria:} Starting from a
parent catalog of eFEDS $\rightarrow$
\href{https://www.legacysurvey.org/dr8}{legacysurvey.org/dr8} cluster
associations (eROMAPPER,
\citealt{Rykoff2014};
\citealt{IderChitham2020};
\citealt{Liu2022}), select targets which meet all of the following criteria:
i) have 13.5\textless{}\texttt{fibertotmag\_r}\textless21.0 or
13.5\textless{}\texttt{fibertotmag\_z}\textless20.0, ii) if detected by
Gaia DR2 then have \emph{G}\textgreater13.5 and
\emph{RP}\textgreater13.5~Vega, iii) do not have existing good quality
SDSS spectroscopy. We assign a range of priorities to targets in this
carton, with BCGs top ranked, followed by candidate member galaxies
according their probability of membership. We assign cadences (exposure
time requests) based on optical brightness.


\noindent\textbf{Target priority options:} 1500,1600-1629

\noindent\textbf{Cadence options:} \texttt{dark\_1x2,dark\_1x4}

\noindent\textbf{Implementation:}
\href{https://github.com/sdss/target_selection/blob/0.3.0/python/target_selection/cartons/bhm_spiders_clusters.py}{bhm\_spiders\_clusters.py}

\noindent\textbf{Number of targets:} 3060

\begin{center}\rule{0.5\linewidth}{0.5pt}\end{center}

\hypertarget{bhm_spiders_agn-efeds_plan0.1.0}{%
\subsection{bhm\_spiders\_agn-efeds}\label{bhm_spiders_agn-efeds_plan0.1.0}}

\noindent\textbf{target\_selection plan:} 0.1.0

\noindent\textbf{target\_selection tag:}
\href{https://github.com/sdss/target_selection/tree/0.1.0/}{0.1.0}

\noindent\textbf{crossmatch plan:} 0.1.0

\noindent\textbf{Summary:} A carton used during SDSS-V plate-mode observations,
that contains candidate AGN targets found in the eROSITA/eFEDS X-ray
survey field. This carton provides optical counterparts to point-like
(unresolved) X-ray sources detected in early reductions ('c940/V2T') of
the eROSITA/eFEDS performance validation field. The sample is expected
to contain a mixture of QSOs, AGN, stars and compact objects. The X-ray
sources have been cross-matched by the eROSITA-DE team to
\href{https://www.legacysurvey.org/dr8/}{legacysurvey.org/dr8}
optical/IR counterparts. All targets in this carton are located within
the eFEDS field (approx 126\textless RA\textless146,
-3\textless Dec\textless+6~deg).

\noindent\textbf{Simplified description of selection criteria:} Starting from a
parent catalog of eFEDS point source $\rightarrow$
\href{https://www.legacysurvey.org/dr8}{legacysurvey.org/dr8}
associations (primarily via NWAY assisted by optical/IR priors computed
via a pre-trained Random Forest, see
\citealt{Salvato2022}, supplemented by counterparts selected via a Likelihood
Ratio using r-band magnitudes), select targets which meet all of the
following criteria: i) have eROSITA detection likelihood\textgreater6.0,
ii) have an X-ray $\rightarrow$ optical/IR cross-match probability of either
\texttt{p\_any}\textgreater0.1 (NWAY associations) or
\emph{LR}\textgreater0.2 (Likelihood Ratio associations), iii) have
\texttt{fibermag\_r}\textgreater16.5 and
\texttt{fibermag\_r}\textless22.0 or \texttt{fibermag\_z}\textless21.0,
and iv) did not receive high quality spectroscopy during the SDSS-IV=
observations of the eFEDS field
(\citealt{Abdurrouf_2021_sdssDR17}). We deprioritise targets if any of the following criteria
are met: a) the target already has existing good quality SDSS
spectroscopy in SDSS DR16, b) the X-ray detection likelihood is
\textless8.0, c) the target is a secondary X-ray$\rightarrow$optical/IR association,
or d) if the optical/IR counterpart was only chosen by the LR method.
All targets were assigned a nominal cadence of:
\texttt{bhm\_spiders\_1x8} (8x15mins dark time).


\noindent\textbf{Target priority options:} 1510-1519

\noindent\textbf{Cadence options:} \texttt{bhm\_spiders\_1x8}

\noindent\textbf{Implementation:}
\href{https://github.com/sdss/target_selection/blob/0.1.0/python/target_selection/cartons/bhm_spiders_agn.py}{bhm\_spiders\_agn.py}

\noindent\textbf{Number of targets:} 12459

\begin{center}\rule{0.5\linewidth}{0.5pt}\end{center}

\hypertarget{bhm_spiders_clusters-efeds-ls-redmapper_plan0.1.0}{%
\subsection{bhm\_spiders\_clusters-efeds-ls-redmapper}\label{bhm_spiders_clusters-efeds-ls-redmapper_plan0.1.0}}

\noindent\textbf{target\_selection plan:} 0.1.0

\noindent\textbf{target\_selection tag:}
\href{https://github.com/sdss/target_selection/tree/0.1.0/}{0.1.0}

\noindent\textbf{crossmatch plan:} 0.1.0

\noindent\textbf{Summary:} A carton used during SDSS-V plate-mode observations,
that contains galaxy cluster targets found in the eROSITA/eFEDS X-ray
survey field. The carton provides a list of galaxies which are candidate
members of clusters selected from early reductions ('c940') of the
eROSITA performance verification survey in the eFEDS field. The parent
sample of galaxy clusters and their member galaxies have been selected
via a joint analysis of X-ray and (several) optical/IR datasets using
the eROMAPPER red-sequence finder algorithm
(\citealt{Rykoff2014};
\citealt{IderChitham2020}). This particular carton relies on optical/IR data
from \href{https://www.legacysurvey.org/dr8/}{legacysurvey.org/dr8}. All
targets in this carton are located within the eFEDS field (approx
126\textless RA\textless146, -3\textless Dec\textless+6~deg).

\noindent\textbf{Simplified description of selection criteria:} Starting from a
parent catalog of eFEDS $\rightarrow$ optical/IR cluster associations, select
targets which meet all of the following criteria: i) are selected by
eROMAPPER applied to
\href{https://www.legacysurvey.org/dr8}{legacysurvey.org/dr8}
photometric data, ii) have eROSITA X-ray detection likelihood
\textgreater{} 8.0, iii) have \texttt{fibermag\_r}\textgreater16.5 and
\texttt{fibermag\_r}\textless21.0 or \texttt{fibermag\_z}\textless20.0,
iv) do not have existing good quality (SDSS or external) spectroscopy.
We assign a range of priorities to targets in this carton, with BCGs top
ranked, followed by candidate member galaxies according their
probability of membership. All targets were assigned a nominal cadence
of: \texttt{bhm\_spiders\_1x8} (8x15mins dark time).


\noindent\textbf{Target priority options:} 1500, 1511-\/-1610

\noindent\textbf{Cadence options:} \texttt{dark\_1x8}

\noindent\textbf{Implementation:}
\href{https://github.com/sdss/target_selection/blob/0.1.0/python/target_selection/cartons/bhm_spiders_clusters.py}{bhm\_spiders\_clusters.py}

\noindent\textbf{Number of targets:} 4432

\begin{center}\rule{0.5\linewidth}{0.5pt}\end{center}

\hypertarget{bhm_spiders_clusters-efeds-sdss-redmapper_plan0.1.0}{%
\subsection{bhm\_spiders\_clusters-efeds-sdss-redmapper}\label{bhm_spiders_clusters-efeds-sdss-redmapper_plan0.1.0}}

\noindent\textbf{target\_selection plan:} 0.1.0

\noindent\textbf{target\_selection tag:}
\href{https://github.com/sdss/target_selection/tree/0.1.0/}{0.1.0}

\noindent\textbf{crossmatch plan:} 0.1.0

\noindent\textbf{Summary:} A carton used during SDSS-V plate-mode observations,
that contains galaxy cluster targets found in the eROSITA/eFEDS X-ray
survey field. The carton provides a list of galaxies which are candidate
members of clusters selected from early reductions ('c940') of the
eROSITA performance verification survey in the eFEDS field. The parent
sample of galaxy clusters and their member galaxies have been selected
via a joint analysis of X-ray and (several) optical/IR datasets using
the eROMAPPER red-sequence finder algorithm
(\citealt{Rykoff2014};
\citealt{IderChitham2020}). This particular carton relies on optical/IR data
from \href{https://www.sdss.org/dr13/}{SDSS/dr13}. All targets in this
carton are located within the eFEDS field (approx
126\textless RA\textless146, -3\textless Dec\textless+6~deg).

\noindent\textbf{Simplified description of selection criteria:} Starting from a
parent catalog of eFEDS $\rightarrow$ optical/IR cluster associations, select
targets which meet all of the following criteria: i) are selected by
eROMAPPER applied to \href{https://www.sdss.org/dr13/}{SDSS/dr13}
photometric data, ii) have eROSITA X-ray detection likelihood
\textgreater{} 8.0, iii) have \texttt{fibermag\_r}\textgreater16.5 and
\texttt{fibermag\_r}\textless21.0 or \texttt{fibermag\_z}\textless20.0,
iv) do not have existing good quality (SDSS or external) spectroscopy.
We assign a range of priorities to targets in this carton, with BCGs top
ranked, followed by candidate member galaxies according their
probability of membership. All targets were assigned a nominal cadence
of: \texttt{bhm\_spiders\_1x8} (8x15mins dark time).


\noindent\textbf{Target priority options:} 1500, 1511-1610

\noindent\textbf{Cadence options:} \texttt{dark\_1x8}

\noindent\textbf{Implementation:}
\href{https://github.com/sdss/target_selection/blob/0.1.0/python/target_selection/cartons/bhm_spiders_clusters.py}{bhm\_spiders\_clusters.py}

\noindent\textbf{Number of targets:} 4304

\begin{center}\rule{0.5\linewidth}{0.5pt}\end{center}

\hypertarget{bhm_spiders_clusters-efeds-hsc-redmapper_plan0.1.0}{%
\subsection{bhm\_spiders\_clusters-efeds-hsc-redmapper}\label{bhm_spiders_clusters-efeds-hsc-redmapper_plan0.1.0}}

\noindent\textbf{target\_selection plan:} 0.1.0

\noindent\textbf{target\_selection tag:}
\href{https://github.com/sdss/target_selection/tree/0.1.0/}{0.1.0}

\noindent\textbf{crossmatch plan:} 0.1.0

\noindent\textbf{Summary:} A carton used during SDSS-V plate-mode observations,
that contains galaxy cluster targets found in the eROSITA/eFEDS X-ray
survey field. The carton provides a list of galaxies which are candidate
members of clusters selected from early reductions ('c940') of the
eROSITA performance verification survey in the eFEDS field. The parent
sample of galaxy clusters and their member galaxies have been selected
via a joint analysis of X-ray and (several) optical/IR datasets using
the eROMAPPER red-sequence finder algorithm
(\citealt{Rykoff2014};
\citealt{IderChitham2020}). This particular carton relies on optical/IR data
from the \href{https://hsc.mtk.nao.ac.jp/ssp/}{Hyper Suprime-Cam Subaru
Strategic Program (HSC-SSP)}. All targets in this carton are located
within the eFEDS field (approx 126\textless RA\textless146,
-3\textless Dec\textless+6~deg).

\noindent\textbf{Simplified description of selection criteria:} Starting from a
parent catalog of eFEDS $\rightarrow$ optical/IR cluster associations, select
targets which meet all of the following criteria: i) are selected by
eROMAPPER applied to \href{https://hsc.mtk.nao.ac.jp/ssp/}{HSC-SSP}
photometric data, ii) have eROSITA X-ray detection likelihood
\textgreater{} 8.0, iii) have \texttt{fibermag\_r}\textgreater16.5 and
\texttt{fibermag\_r}\textless21.0 or \texttt{fibermag\_z}\textless20.0,
iv) do not have existing good quality (SDSS or external) spectroscopy.
We assign a range of priorities to targets in this carton, with BCGs top
ranked, followed by candidate member galaxies according their
probability of membership. All targets were assigned a nominal cadence
of: \texttt{bhm\_spiders\_1x8} (8x15mins dark time).


\noindent\textbf{Target priority options:} 1500, 1511-1610

\noindent\textbf{Cadence options:} \texttt{dark\_1x8}

\noindent\textbf{Implementation:}
\href{https://github.com/sdss/target_selection/blob/0.1.0/python/target_selection/cartons/bhm_spiders_clusters.py}{bhm\_spiders\_clusters.py}

\noindent\textbf{Number of targets:} 924

\begin{center}\rule{0.5\linewidth}{0.5pt}\end{center}

\hypertarget{bhm_spiders_clusters-efeds-erosita_plan0.1.0}{%
\subsection{bhm\_spiders\_clusters-efeds-erosita}\label{bhm_spiders_clusters-efeds-erosita_plan0.1.0}}

\noindent\textbf{target\_selection plan:} 0.1.0

\noindent\textbf{target\_selection tag:}
\href{https://github.com/sdss/target_selection/tree/0.1.0/}{0.1.0}

\noindent\textbf{crossmatch plan:} 0.1.0

\noindent\textbf{Summary:} A carton used during SDSS-V plate-mode observations,
that contains galaxy cluster targets found in the eROSITA/eFEDS X-ray
survey field. The carton provides a list of galaxies which are candidate
members of clusters selected from early reductions ('c940') of the
eROSITA performance verification survey in the eFEDS field. The parent
sample of galaxy clusters and their member galaxies have been selected
via a joint analysis of X-ray and (several) optical/IR datasets. This
particular carton includes counterparts to X-ray extended sources that
were not selected by the eROMAPPER red seuence finder algorithm when
applied to any of the legacysurvey/dr8, SDSS/dr13 or HSC-SSP datasets
(i.e. complementary to the cartons:
\texttt{bhm\_spiders\_clusters-efeds-ls-redmapper},
\texttt{bhm\_spiders\_clusters-efeds-sdss-redmapper} and
\texttt{bhm\_spiders\_clusters-efeds-hsc-redmapper}). All targets in
this carton are located within the eFEDS field (approx
126\textless RA\textless146, -3\textless Dec\textless+6~deg).

\noindent\textbf{Simplified description of selection criteria:} Starting from a
parent catalog of eFEDS $\rightarrow$ optical/IR cluster associations, select
targets which meet all of the following criteria: i) are identified as
being X-ray extended but not selected via the eROMAPPER algorithm, ii)
have eROSITA X-ray detection likelihood \textgreater{} 8.0, iii) have
\texttt{fibermag\_r}\textgreater16.5 and
\texttt{fibermag\_r}\textless21.0 or \texttt{fibermag\_z}\textless20.0,
iv) do not have existing good quality (SDSS or external) spectroscopy.
We assign a range of priorities to targets in this carton, with BCGs top
ranked, followed by candidate member galaxies according their
probability of membership. All targets were assigned a nominal cadence
of: \texttt{bhm\_spiders\_1x8} (8x15mins dark time).


\noindent\textbf{Target priority options:} 1500, 1511-1535

\noindent\textbf{Cadence options:} \texttt{dark\_1x8}

\noindent\textbf{Implementation:}
\href{https://github.com/sdss/target_selection/blob/0.1.0/python/target_selection/cartons/bhm_spiders_clusters.py}{bhm\_spiders\_clusters.py}

\noindent\textbf{Number of targets:} 15

\begin{center}\rule{0.5\linewidth}{0.5pt}\end{center}

\hypertarget{bhm_aqmes_med_plan0.1.0}{%
\subsection{bhm\_aqmes\_med}\label{bhm_aqmes_med_plan0.1.0}}

\noindent\textbf{target\_selection plan:} 0.1.0

\noindent\textbf{target\_selection tag:}
\href{https://github.com/sdss/target_selection/tree/0.1.0/}{0.1.0}

\noindent\textbf{crossmatch plan:} 0.1.0

\noindent\textbf{Summary:} Spectroscopically confirmed optically bright SDSS
QSOs, selected from the SDSS QSO catalog (DR16Q,
\citealt{Lyke2020}). Located in 36 mostly disjoint fields within the SDSS QSO
footprint that were pre-selected to contain higher than average numbers
of bright QSOs and CSC targets. The list of field centres can be found
within
\href{https://github.com/sdss/target_selection/blob/0.1.0/python/target_selection/masks/candidate_target_fields_bhm_aqmes_med_v0.2.1.fits}{the
target\_selection repository}.

\noindent\textbf{Simplified description of selection criteria:} Select all
objects from SDSS DR16 QSO catalog that have
16.0\textless{}\texttt{sdss\_psfmag\_i}\textless19.1~AB, that lie within
1.49~degrees of at least one AQMES-medium field location.


\noindent\textbf{Target priority options:} 1100

\noindent\textbf{Cadence options:} \texttt{bhm\_aqmes\_medium\_10x4}

\noindent\textbf{Implementation:}
\href{https://github.com/sdss/target_selection/blob/0.1.0/python/target_selection/cartons/bhm_aqmes.py}{bhm\_aqmes.py}

\noindent\textbf{Number of targets:} 2663

\begin{center}\rule{0.5\linewidth}{0.5pt}\end{center}

\hypertarget{bhm_aqmes_med-faint_plan0.1.0}{%
\subsection{bhm\_aqmes\_med-faint}\label{bhm_aqmes_med-faint_plan0.1.0}}

\noindent\textbf{target\_selection plan:} 0.1.0

\noindent\textbf{target\_selection tag:}
\href{https://github.com/sdss/target_selection/tree/0.1.0/}{0.1.0}

\noindent\textbf{crossmatch plan:} 0.1.0

\noindent\textbf{Summary:} Spectroscopically confirmed optically faint SDSS QSOs,
selected from the SDSS QSO catalog (DR16Q,
\citealt{Lyke2020}). Located in 36 mostly disjoint fields within the SDSS QSO
footprint that were pre-selected to contain higher than average numbers
of bright QSOs and CSC targets. The list of field centres can be found
within
\href{https://github.com/sdss/target_selection/blob/0.1.0/python/target_selection/masks/candidate_target_fields_bhm_aqmes_med_v0.2.1.fits}{the
target\_selection repository}.

\noindent\textbf{Simplified description of selection criteria:} Select all
objects from SDSS DR16 QSO catalog that have
19.1\textless{}\texttt{sdss\_psfmag\_i}\textless21.0~AB, that lie within
1.49~degrees of at least one AQMES-medium field location.


\noindent\textbf{Target priority options:} 3100

\noindent\textbf{Cadence options:} \texttt{bhm\_aqmes\_medium\_10x4}

\noindent\textbf{Implementation:}
\href{https://github.com/sdss/target_selection/blob/0.1.0/python/target_selection/cartons/bhm_aqmes.py}{bhm\_aqmes.py}

\noindent\textbf{Number of targets:} 16853

\begin{center}\rule{0.5\linewidth}{0.5pt}\end{center}

\hypertarget{bhm_aqmes_wide2_plan0.1.0}{%
\subsection{bhm\_aqmes\_wide2}\label{bhm_aqmes_wide2_plan0.1.0}}

\noindent\textbf{target\_selection plan:} 0.1.0

\noindent\textbf{target\_selection tag:}
\href{https://github.com/sdss/target_selection/tree/0.1.0/}{0.1.0}

\noindent\textbf{crossmatch plan:} 0.1.0

\noindent\textbf{Summary:} Spectroscopically confirmed optically bright SDSS
QSOs, selected from the SDSS QSO catalog (DR16Q,
\citealt{Lyke2020}). Located in 330 fields within the SDSS QSO footprint,
where the choice of survey area prioritized fields that overlapped with
the SPIDERS footprint (approx 180\textless{}\emph{l}\textless360~deg),
and/or had higher than average numbers of bright QSOs and CSC targets.
The list of field centres can be found
\href{https://github.com/sdss/target_selection/blob/0.1.0/python/target_selection/masks/candidate_target_fields_bhm_aqmes_wide_v0.2.1.fits}{within
the target\_selection repository}.

\noindent\textbf{Simplified description of selection criteria:} Select all
objects from SDSS DR16 QSO catalog that have
16.0\textless{}\texttt{sdss\_psfmag\_i}\textless19.1~AB, and that lie
within 1.49~degrees of at least one AQMES-wide2 field location.


\noindent\textbf{Target priority options:} 1210

\noindent\textbf{Cadence options:} \texttt{bhm\_aqmes\_wide\_2x4}

\noindent\textbf{Implementation:}
\href{https://github.com/sdss/target_selection/blob/0.1.0/python/target_selection/cartons/bhm_aqmes.py}{bhm\_aqmes.py}

\noindent\textbf{Number of targets:} 18376

\begin{center}\rule{0.5\linewidth}{0.5pt}\end{center}

\hypertarget{bhm_aqmes_wide2-faint_plan0.1.0}{%
\subsection{bhm\_aqmes\_wide2-faint}\label{bhm_aqmes_wide2-faint_plan0.1.0}}

\noindent\textbf{target\_selection plan:} 0.1.0

\noindent\textbf{target\_selection tag:}
\href{https://github.com/sdss/target_selection/tree/0.1.0/}{0.1.0}

\noindent\textbf{crossmatch plan:} 0.1.0

\noindent\textbf{Summary:} Spectroscopically confirmed optically faint SDSS QSOs,
selected from the SDSS QSO catalog (DR16Q,
\citealt{Lyke2020}). Located in 330 fields within the SDSS QSO footprint,
where the choice of survey area prioritized field that overlapped with
the SPIDERS footprint (approx 180\textless{}\emph{l}\textless360~deg),
and/or had higher than average numbers of bright QSOs and CSC targets.
The list of field centres can be found
\href{https://github.com/sdss/target_selection/blob/0.1.0/python/target_selection/masks/candidate_target_fields_bhm_aqmes_wide_v0.2.1.fits}{within
the target\_selection repository}.

\noindent\textbf{Simplified description of selection criteria:} Select all
objects from SDSS DR16 QSO catalog that have
19.1\textless{}\texttt{sdss\_psfmag\_i}\textless21.0~AB, and that lie
within 1.49~degrees of at least one AQMES-wide2 field location.


\noindent\textbf{Target priority options:} 3210

\noindent\textbf{Cadence options:} \texttt{bhm\_aqmes\_wide\_2x4}

\noindent\textbf{Implementation:}
\href{https://github.com/sdss/target_selection/blob/0.1.0/python/target_selection/cartons/bhm_aqmes.py}{bhm\_aqmes.py}

\noindent\textbf{Number of targets:}

\begin{center}\rule{0.5\linewidth}{0.5pt}\end{center}

\hypertarget{bhm_aqmes_wide3_plan0.1.0}{%
\subsection{bhm\_aqmes\_wide3}\label{bhm_aqmes_wide3_plan0.1.0}}

\noindent\textbf{target\_selection plan:} 0.1.0

\noindent\textbf{target\_selection tag:}
\href{https://github.com/sdss/target_selection/tree/0.1.0/}{0.1.0}

\noindent\textbf{crossmatch plan:} 0.1.0

\noindent\textbf{Summary:} Spectroscopically confirmed optically bright SDSS
QSOs, selected from the SDSS QSO catalog (DR16Q,
\citealt{Lyke2020}). Located in 95 fields within the SDSS QSO footprint,
where the choice of survey area prioritized fields having higher than
average numbers of bright QSOs and CSC targets. The list of field
centres can be found
\href{https://github.com/sdss/target_selection/blob/0.1.0/python/target_selection/masks/candidate_target_fields_bhm_aqmes_wide_v0.2.1.fits}{within
the target\_selection repository}.

\noindent\textbf{Simplified description of selection criteria:} Select all
objects from SDSS DR16 QSO catalog that have
16.0\textless{}\texttt{sdss\_psfmag\_i}\textless19.1~AB, and that lie
within 1.49~degrees of at least one AQMES-wide3 field location.


\noindent\textbf{Target priority options:} 1210

\noindent\textbf{Cadence options:} \texttt{bhm\_aqmes\_wide\_3x4}

\noindent\textbf{Implementation:}
\href{https://github.com/sdss/target_selection/blob/0.1.0/python/target_selection/cartons/bhm_aqmes.py}{bhm\_aqmes.py}

\noindent\textbf{Number of targets:} 5785

\begin{center}\rule{0.5\linewidth}{0.5pt}\end{center}

\hypertarget{bhm_aqmes_wide3-faint_plan0.1.0}{%
\subsection{bhm\_aqmes\_wide3-faint}\label{bhm_aqmes_wide3-faint_plan0.1.0}}

\noindent\textbf{target\_selection plan:} 0.1.0

\noindent\textbf{target\_selection tag:}
\href{https://github.com/sdss/target_selection/tree/0.1.0/}{0.1.0}

\noindent\textbf{crossmatch plan:} 0.1.0

\noindent\textbf{Summary:} Spectroscopically confirmed optically faint SDSS QSOs,
selected from the SDSS QSO catalog (DR16Q,
\citealt{Lyke2020}). Located in 95 fields within the SDSS QSO footprint,
where the choice of survey area prioritized fields having higher than
average numbers of bright QSOs and CSC targets. The list of field
centres can be found
\href{https://github.com/sdss/target_selection/blob/0.1.0/python/target_selection/masks/candidate_target_fields_bhm_aqmes_wide_v0.2.1.fits}{within
the target\_selection repository}.

\noindent\textbf{Simplified description of selection criteria:} Select all
objects from SDSS DR16 QSO catalog that have
19.1\textless{}\texttt{sdss\_psfmag\_i}\textless21.0~AB, and that lie
within 1.49~degrees of at least one AQMES-wide2 field location.


\noindent\textbf{Target priority options:} 3210

\noindent\textbf{Cadence options:} \texttt{dark\_2x4}

\noindent\textbf{Implementation:}
\href{https://github.com/sdss/target_selection/blob/0.1.0/python/target_selection/cartons/bhm_aqmes.py}{bhm\_aqmes.py}

\noindent\textbf{Number of targets:} 35803

\begin{center}\rule{0.5\linewidth}{0.5pt}\end{center}

\hypertarget{bhm_aqmes_bonus-dark_plan0.1.0}{%
\subsection{bhm\_aqmes\_bonus-dark}\label{bhm_aqmes_bonus-dark_plan0.1.0}}

\noindent\textbf{target\_selection plan:} 0.1.0

\noindent\textbf{target\_selection tag:}
\href{https://github.com/sdss/target_selection/tree/0.1.0/}{0.1.0}

\noindent\textbf{crossmatch plan:} 0.1.0

\noindent\textbf{Summary:} Spectroscopically confirmed SDSS QSOs, with magnitudes
in the range appropriate for observations in dark time, selected from
the SDSS QSO catalog (DR16Q,
\citealt{Lyke2020}). Located anywhere within the SDSS DR16Q footprint.

\noindent\textbf{Simplified description of selection criteria:} Select all
objects from SDSS DR16 QSO catalog that have
16.5\textless{}\texttt{sdss\_psfmag\_i}\textless21.5~AB


\noindent\textbf{Target priority options:} 3300

\noindent\textbf{Cadence options:} \texttt{dark\_spiders\_1x4}

\noindent\textbf{Implementation:}
\href{https://github.com/sdss/target_selection/blob/0.1.0/python/target_selection/cartons/bhm_aqmes.py}{bhm\_aqmes.py}

\noindent\textbf{Number of targets:} 579590

\begin{center}\rule{0.5\linewidth}{0.5pt}\end{center}

\hypertarget{bhm_aqmes_bonus-bright_plan0.1.0}{%
\subsection{bhm\_aqmes\_bonus-bright}\label{bhm_aqmes_bonus-bright_plan0.1.0}}

\noindent\textbf{target\_selection plan:} 0.1.0

\noindent\textbf{target\_selection tag:}
\href{https://github.com/sdss/target_selection/tree/0.1.0/}{0.1.0}

\noindent\textbf{crossmatch plan:} 0.1.0

\noindent\textbf{Summary:} Spectroscopically confirmed SDSS QSOs, with magnitudes
in the range appropriate for observations in bright time, selected from
the SDSS QSO catalog (DR16Q,
\citealt{Lyke2020}). Located anywhere within the SDSS DR16Q footprint.

\noindent\textbf{Simplified description of selection criteria:} Select all
objects from SDSS DR16 QSO catalog that have
14.0\textless{}\texttt{sdss\_psfmag\_i}\textless18.0~AB


\noindent\textbf{Target priority options:} 4040

\noindent\textbf{Cadence options:} \texttt{bhm\_boss\_bright\_3x1}

\noindent\textbf{Implementation:}
\href{https://github.com/sdss/target_selection/blob/0.1.0/python/target_selection/cartons/bhm_aqmes.py}{bhm\_aqmes.py}

\noindent\textbf{Number of targets:} 10848

\begin{center}\rule{0.5\linewidth}{0.5pt}\end{center}

\hypertarget{bhm_rm_ancillary_plan0.1.0}{%
\subsection{bhm\_rm\_ancillary}\label{bhm_rm_ancillary_plan0.1.0}}

\noindent\textbf{target\_selection plan:} 0.1.0

\noindent\textbf{target\_selection tag:}
\href{https://github.com/sdss/target_selection/tree/0.1.0/}{0.1.0}

\noindent\textbf{crossmatch plan:} 0.1.0

\noindent\textbf{Summary:} A supporting sample of candidate QSOs which have been
selected by the Gaia-unWISE AGN catalog
(\citealt{Shu2019}) and/or the SDSS XDQSO catalog
(\citealt{Bovy2011}). These targets are located within five (+1 backup) well
known survey fields (SDSS-RM, COSMOS, XMM-LSS, ECDFS, CVZ-S/SEP, and
ELIAS-S1).

\noindent\textbf{Simplified description of selection criteria:} Starting from a
parent catalog of optically selected objects in the RM fields (as
presented by
\citealt{Yang2022}), select candidate QSOs that satisfy all of the
following: i) are identified via external ancillary methods
(\texttt{photo\_bitmask\ \&\ 3\ !=\ 0}); ii) have
15\textless{}\texttt{psfmag\_i}\textless21.5~AB; iii) do not have
significant detections (\textgreater3$\sigma$) of parallax and/or proper motion
in Gaia DR2; iv) are not classified as a STAR in SDSS DR16 spectroscopy;
and v) do not lie in the SDSS-RM field


\noindent\textbf{Target priority options:} 1004-1050

\noindent\textbf{Cadence options:} \texttt{bhm\_rm\_174x8}

\noindent\textbf{Implementation:}
\href{https://github.com/sdss/target_selection/blob/0.1.0/python/target_selection/cartons/bhm_rm.py}{bhm\_rm.py}

\noindent\textbf{Number of targets:} 948

\begin{center}\rule{0.5\linewidth}{0.5pt}\end{center}

\hypertarget{bhm_rm_core_plan0.1.0}{%
\subsection{bhm\_rm\_core}\label{bhm_rm_core_plan0.1.0}}

\noindent\textbf{target\_selection plan:} 0.1.0

\noindent\textbf{target\_selection tag:}
\href{https://github.com/sdss/target_selection/tree/0.1.0/}{0.1.0}

\noindent\textbf{crossmatch plan:} 0.1.0

\noindent\textbf{Summary:} A sample of candidate QSOs selected via the methods
presented by
\citet{Yang2022}. These targets are located within five (+1 backup) well
known survey fields (SDSS-RM, COSMOS, XMM-LSS, ECDFS, CVZ-S/SEP, and
ELIAS-S1).

\noindent\textbf{Simplified description of selection criteria:} Starting from a
parent catalog of optically selected objects in the RM fields (as
presented by
\citealt{Yang2022}), select candidate QSOs that satisfy all of the
following: i) are identified via the Skew-T algorithm
(\texttt{skewt\_qso\ ==\ 1}, and \texttt{skewt\_qso\_prior\ ==\ 1} for
targets in the CVZ-S/SEP field); ii) have
17\textless{}\texttt{psfmag\_i}\textless21.5~AB; iii) do not have
significant detections (\textgreater3$\sigma$) of parallax and/or proper motion
in Gaia DR2; iv) are not classified as a STAR in SDSS DR16 spectroscopy;
vi) have detections in all of the \emph{gri} bands (a Gaia detection is
sufficient in the CVZ-S/SEP field); and vii) do not lie in the SDSS-RM
field


\noindent\textbf{Target priority options:} 1002-1050

\noindent\textbf{Cadence options:} \texttt{bhm\_rm\_174x8}

\noindent\textbf{Implementation:}
\href{https://github.com/sdss/target_selection/blob/0.1.0/python/target_selection/cartons/bhm_rm.py}{bhm\_rm.py}

\noindent\textbf{Number of targets:}

\begin{center}\rule{0.5\linewidth}{0.5pt}\end{center}

\hypertarget{bhm_rm_var_plan0.1.0}{%
\subsection{bhm\_rm\_var}\label{bhm_rm_var_plan0.1.0}}

\noindent\textbf{target\_selection plan:} 0.1.0

\noindent\textbf{target\_selection tag:}
\href{https://github.com/sdss/target_selection/tree/0.1.0/}{0.1.0}

\noindent\textbf{crossmatch plan:} 0.1.0

\noindent\textbf{Summary:} A sample of candidate QSOs selected via their optical
variability properties, as presented by
\citet{Yang2022}. These targets are located within five (+1 backup) well
known survey fields (SDSS-RM, COSMOS, XMM-LSS, ECDFS, CVZ-S/SEP, and
ELIAS-S1).

\noindent\textbf{Simplified description of selection criteria:} Starting from a
parent catalog of optically selected objects in the RM fields (as
presented by
\citealt{Yang2022}), select candidate QSOs that satisfy all of the
following: i) have significant variability in the DES or PanSTARRS1
multi-epoch photometry (\texttt{var\_sn{[}g{]}}\textgreater3 and
\texttt{var\_rms{[}g{]}}\textgreater0.05); ii) have
17\textless{}\texttt{psfmag\_i}\textless20.5~AB; iii) do not have
significant detections (\textgreater3$\sigma$) of parallax and/or proper motion
in Gaia DR2; iv) are not classified as a STAR in SDSS DR16 spectroscopy;
and vi) do not lie in the SDSS-RM field


\noindent\textbf{Target priority options:} 1003-1050

\noindent\textbf{Cadence options:} \texttt{bhm\_rm\_174x8}

\noindent\textbf{Implementation:}
\href{https://github.com/sdss/target_selection/blob/0.1.0/python/target_selection/cartons/bhm_rm.py}{bhm\_rm.py}

\noindent\textbf{Number of targets:} 992

\begin{center}\rule{0.5\linewidth}{0.5pt}\end{center}

\hypertarget{bhm_rm_known-spec_plan0.1.0}{%
\subsection{bhm\_rm\_known-spec}\label{bhm_rm_known-spec_plan0.1.0}}

\noindent\textbf{target\_selection plan:} 0.1.0

\noindent\textbf{target\_selection tag:}
\href{https://github.com/sdss/target_selection/tree/0.1.0/}{0.1.0}

\noindent\textbf{crossmatch plan:} 0.1.0

\noindent\textbf{Summary:} A sample of known QSOs identified through optical
spectroscopy from various projects, as collated by
\citet{Yang2022}. These targets are located within five (+1 backup) well
known survey fields (SDSS-RM, COSMOS, XMM-LSS, ECDFS, CVZ-S/SEP, and
ELIAS-S1).

\noindent\textbf{Simplified description of selection criteria:} Starting from a
parent catalog of optically selected objects in the RM fields (as
presented by
\citealt{Yang2022}), select targets which satisfy all of the following: i)
are flagged as having a spectroscopic identification in the parent
catalog; ii) have 15\textless{}\texttt{psfmag\_i}\textless21.7~AB
(SDSS-RM, CDFS, ELIAS-S1, CVZ-S/SEP fields), or
15\textless{}\texttt{psfmag\_i}\textless21.5~AB (COSMOS and XMM-LSS
fields); iii) have a spectroscopic redshift in the range
0.005\textless{}\emph{z}\textless7; iv) are not classified as a STAR in
SDSS DR16 spectroscopy;


\noindent\textbf{Target priority options:} 1001-1050

\noindent\textbf{Cadence options:} \texttt{bhm\_rm\_174x8}

\noindent\textbf{Implementation:}
\href{https://github.com/sdss/target_selection/blob/0.1.0/python/target_selection/cartons/bhm_rm.py}{bhm\_rm.py}

\noindent\textbf{Number of targets:} 2992

\begin{center}\rule{0.5\linewidth}{0.5pt}\end{center}

\hypertarget{bhm_csc_apogee_plan0.1.0}{%
\subsection{bhm\_csc\_apogee}\label{bhm_csc_apogee_plan0.1.0}}

\noindent\textbf{target\_selection plan:} 0.1.0

\noindent\textbf{target\_selection tag:}
\href{https://github.com/sdss/target_selection/tree/0.1.0/}{0.1.0}

\noindent\textbf{crossmatch plan:} 0.1.0

\noindent\textbf{Summary:} X-ray sources from the CSC2.0 source catalog with NIR
counterparts in 2MASS PSC

\noindent\textbf{Simplified description of selection criteria:} Starting from the
parent catalog of CSC2.0 sources with optical/IR counterparts
(\texttt{bhm\_csc}). Select entries satisfying the following criteria:
i) counterpart is from the 2MASS catalog, ii) 2MASS \emph{H}-band
magnitude measurement is in the accepted range for SDSS-V:
10.0\textless{}\emph{H}\textless15.0.


\noindent\textbf{Target priority options:} 4000

\noindent\textbf{Cadence options:} \texttt{bhm\_csc\_apogee\_3x1}

\noindent\textbf{Implementation:}
\href{https://github.com/sdss/target_selection/blob/0.1.0/python/target_selection/cartons/bhm_csc.py}{bhm\_csc.py}

\noindent\textbf{Number of targets:} 10633

\begin{center}\rule{0.5\linewidth}{0.5pt}\end{center}

\hypertarget{bhm_csc_boss-dark_plan0.1.0}{%
\subsection{bhm\_csc\_boss-dark}\label{bhm_csc_boss-dark_plan0.1.0}}

\noindent\textbf{target\_selection plan:} 0.1.0

\noindent\textbf{target\_selection tag:}
\href{https://github.com/sdss/target_selection/tree/0.1.0/}{0.1.0}

\noindent\textbf{crossmatch plan:} 0.1.0

\noindent\textbf{Summary:} X-ray sources from the CSC2.0 source catalog with
counterparts in PanSTARRS1-DR1

\noindent\textbf{Simplified description of selection criteria:} Starting from the
parent catalog of CSC sources with optical/IR counterparts
(\texttt{bhm\_csc}). Select entries satisfying the following criteria:
i) counterpart is from the PanSTARRS1 catalog, ii) optical
flux/magnitude is in a range suited to SDSS-V dark time observations:
16.0\textless{}\texttt{psfmag\_i}\textless22.0~AB.


\noindent\textbf{Target priority options:} 3000

\noindent\textbf{Cadence options:} \texttt{bhm\_csc\_boss\_1x4}

\noindent\textbf{Implementation:}
\href{https://github.com/sdss/target_selection/blob/0.1.0/python/target_selection/cartons/bhm_csc.py}{bhm\_csc.py}

\noindent\textbf{Number of targets:} 65350

\begin{center}\rule{0.5\linewidth}{0.5pt}\end{center}

\hypertarget{bhm_csc_boss-bright_plan0.1.0}{%
\subsection{bhm\_csc\_boss-bright}\label{bhm_csc_boss-bright_plan0.1.0}}

\noindent\textbf{target\_selection plan:} 0.1.0

\noindent\textbf{target\_selection tag:}
\href{https://github.com/sdss/target_selection/tree/0.1.0/}{0.1.0}

\noindent\textbf{crossmatch plan:} 0.1.0

\noindent\textbf{Summary:} X-ray sources from the CSC2.0 source catalog with
counterparts in PanSTARRS1-DR1

\noindent\textbf{Simplified description of selection criteria:} Starting from the
parent catalog of CSC sources with optical/IR counterparts
(\texttt{bhm\_csc}). Select entries satisfying the following criteria:
i) counterpart is from the PanSTARRS1 catalog, ii) optical
flux/magnitude is in a range suited to SDSS-V bright time observations:
13.0\textless{}\texttt{psfmag\_i}\textless18.0~AB.


\noindent\textbf{Target priority options:} 4000

\noindent\textbf{Cadence options:} \texttt{bhm\_csc\_boss\_1x1}

\noindent\textbf{Implementation:}
\href{https://github.com/sdss/target_selection/blob/0.1.0/python/target_selection/cartons/bhm_csc.py}{bhm\_csc.py}

\noindent\textbf{Number of targets:} 22173

\begin{center}\rule{0.5\linewidth}{0.5pt}\end{center}

\hypertarget{bhm_gua_dark_plan0.1.0}{%
\subsection{bhm\_gua\_dark}\label{bhm_gua_dark_plan0.1.0}}

\noindent\textbf{target\_selection plan:} 0.1.0

\noindent\textbf{target\_selection tag:}
\href{https://github.com/sdss/target_selection/tree/0.1.0/}{0.1.0}

\noindent\textbf{crossmatch plan:} 0.1.0

\noindent\textbf{Summary:} A sample of optically faint candidate AGN lacking
spectroscopic confirmations, derived from the parent sample presented by
\citet{Shu2019}, who applied a machine-learning approach to select QSO
candidates from a combination of the Gaia DR2 and unWISE catalogs.

\noindent\textbf{Simplified description of selection criteria:} Starting with the
\citet{Shu2019} catalog, select targets which satisfy the following
criteria: i) have a Random Forest probability of being a QSO
of\textgreater0.8, ii) are in the (dereddened) magnitude range suitable
for BOSS spectroscopy in dark time
(\emph{G}\textsubscript{dered}\textgreater16.5 and
\emph{RP}\textsubscript{dered}\textgreater16.5, as well as
\emph{G}\textsubscript{dered}\textless21.2 or
\emph{RP}\textsubscript{dered}\textless21.0,~Vega mag), iii) do not have
good quality optical spectroscopic measurements in SDSS DR16


\noindent\textbf{Target priority options:} 3400

\noindent\textbf{Cadence options:} \texttt{bhm\_spiders\_1x4}

\noindent\textbf{Implementation:}
\href{https://github.com/sdss/target_selection/blob/0.1.0/python/target_selection/cartons/bhm_gua.py}{bhm\_gua.py}

\noindent\textbf{Number of targets:} 2085729

\begin{center}\rule{0.5\linewidth}{0.5pt}\end{center}

\hypertarget{bhm_gua_bright_plan0.1.0}{%
\subsection{bhm\_gua\_bright}\label{bhm_gua_bright_plan0.1.0}}

\noindent\textbf{target\_selection plan:} 0.1.0

\noindent\textbf{target\_selection tag:}
\href{https://github.com/sdss/target_selection/tree/0.1.0/}{0.1.0}

\noindent\textbf{crossmatch plan:} 0.1.0

\noindent\textbf{Summary:} A sample of optically bright candidate AGN lacking
spectroscopic confirmations, derived from the parent sample presented by
\citet{Shu2019}, who applied a machine-learning approach to select QSO
candidates from a combination of the Gaia DR2 and unWISE catalogs.

\noindent\textbf{Simplified description of selection criteria:} Starting with the
\citet{Shu2019} catalog, select targets which satisfy the following
criteria: i) have a Random Forest probability of being a QSO
of\textgreater0.8, ii) are in the (dereddened) magnitude range suitable
for BOSS spectroscopy in bright time
(\emph{G}\textsubscript{dered}\textgreater13.0 and
\emph{RP}\textsubscript{dered}\textgreater13.5, as well as
\emph{G}\textsubscript{dered}\textless18.5 or
\emph{RP}\textsubscript{dered}\textless18.5,~Vega mag), iii) do not have
good quality optical spectroscopic measurements in SDSS DR16.


\noindent\textbf{Target priority options:} 4040

\noindent\textbf{Cadence options:} \texttt{bhm\_boss\_bright\_3x1}

\noindent\textbf{Implementation:}
\href{https://github.com/sdss/target_selection/blob/0.1.0/python/target_selection/cartons/bhm_gua.py}{bhm\_gua.py}

\noindent\textbf{Number of targets:} 237236

\begin{center}\rule{0.5\linewidth}{0.5pt}\end{center}

\hypertarget{bhm_csc_apogee_plan0.5.0}{%
\subsection{bhm\_csc\_apogee}\label{bhm_csc_apogee_plan0.5.0}}

\noindent\textbf{target\_selection plan:} 0.5.0

\noindent\textbf{target\_selection tag:}
\href{https://github.com/sdss/target_selection/tree/0.3.0/}{0.3.0}

\noindent\textbf{crossmatch plan:} 0.5.0

\noindent\textbf{Summary:} X-ray sources from the CSC2.0 source catalog with NIR
counterparts in 2MASS PSC

\noindent\textbf{Simplified description of selection criteria:} Starting from the
parent catalog of CSC2.0 sources with optical/IR counterparts
(\texttt{bhm\_csc}). Select entries satisfying the following criteria:
i) counterpart is from the 2MASS catalog, ii) 2MASS \emph{H}-band
magnitude measurement is in the accepted range for SDSS-V:
10.0\textless{}\emph{H}\textless15.0.


\noindent\textbf{Target priority options:} 4000

\noindent\textbf{Cadence options:} \texttt{bhm\_csc\_apogee\_3x1}

\noindent\textbf{Implementation:}
\href{https://github.com/sdss/target_selection/blob/0.3.0/python/target_selection/cartons/bhm_csc.py}{bhm\_csc.py}

\noindent\textbf{Number of targets:} 13437

\begin{center}\rule{0.5\linewidth}{0.5pt}\end{center}

\hypertarget{bhm_csc_boss_dark_plan0.5.0}{%
\subsection{bhm\_csc\_boss\_dark}\label{bhm_csc_boss_dark_plan0.5.0}}

\noindent\textbf{target\_selection plan:} 0.5.0

\noindent\textbf{target\_selection tag:}
\href{https://github.com/sdss/target_selection/tree/0.3.0/}{0.3.0}

\noindent\textbf{crossmatch plan:} 0.5.0

\noindent\textbf{Summary:} X-ray sources from the CSC2.0 source catalog with
counterparts in PanSTARRS1-DR1

\noindent\textbf{Simplified description of selection criteria:} Starting from the
parent catalog of CSC sources with optical/IR counterparts
(\texttt{bhm\_csc}). Select entries satisfying the following criteria:
i) counterpart is from the PanSTARRS1 catalog, ii) optical
flux/magnitude is in a range suited to SDSS-V dark time observations:
16.0\textless{}\texttt{psfmag\_i}\textless22.0~AB .


\noindent\textbf{Target priority options:} 3000

\noindent\textbf{Cadence options:} \texttt{bhm\_csc\_boss\_1x4}

\noindent\textbf{Implementation:}
\href{https://github.com/sdss/target_selection/blob/0.3.0/python/target_selection/cartons/bhm_csc.py}{bhm\_csc.py}

\noindent\textbf{Number of targets:} 65296

\begin{center}\rule{0.5\linewidth}{0.5pt}\end{center}

\hypertarget{bhm_csc_boss_bright_plan0.5.0}{%
\subsection{bhm\_csc\_boss\_bright}\label{bhm_csc_boss_bright_plan0.5.0}}

\noindent\textbf{target\_selection plan:} 0.5.0

\noindent\textbf{target\_selection tag:}
\href{https://github.com/sdss/target_selection/tree/0.3.0/}{0.3.0}

\noindent\textbf{crossmatch plan:} 0.5.0

\noindent\textbf{Summary:} X-ray sources from the CSC2.0 source catalog with
counterparts in PanSTARRS1-DR1

\noindent\textbf{Simplified description of selection criteria:} Starting from the
parent catalog of CSC sources with optical/IR counterparts
(\texttt{bhm\_csc}). Select entries satisfying the following criteria:
i) counterpart is from the PanSTARRS1 catalog, ii) optical
flux/magnitude is in a range suited to SDSS-V bright time observations:
13.0\textless{}\texttt{psfmag\_i}\textless18.0~AB .


\noindent\textbf{Target priority options:} 4000

\noindent\textbf{Cadence options:} \texttt{bhm\_csc\_boss\_1x1}

\noindent\textbf{Implementation:}
\href{https://github.com/sdss/target_selection/blob/0.3.0/python/target_selection/cartons/bhm_csc.py}{bhm\_csc.py}

\noindent\textbf{Number of targets:} 22146

\begin{center}\rule{0.5\linewidth}{0.5pt}\end{center}

\hypertarget{openfibertargets_nov2020_11_plan0.5.3}{%
\subsection{openfibertargets\_nov2020\_11}\label{openfibertargets_nov2020_11_plan0.5.3}}

\noindent\textbf{target\_selection plan:} 0.5.3

\noindent\textbf{target\_selection tag:}
\href{https://github.com/sdss/target_selection/tree/0.3.4/}{0.3.4}

\noindent\textbf{crossmatch plan:} 0.5.0

\noindent\textbf{Summary:} Quasars observed during SDSS-I and -II which had only
a single epoch to changes in emission lines and absorption line troughs
over timescales of years to decades

\noindent\textbf{Simplified description of selection criteria:} manual


\noindent\textbf{Target priority options:} 6080

\noindent\textbf{Cadence options:} \texttt{bright\_1x1,\ dark\_1x1}

\noindent\textbf{Implementation:} n/a

\noindent\textbf{Number of targets:} 39684

\begin{center}\rule{0.5\linewidth}{0.5pt}\end{center}

\hypertarget{openfibertargets_nov2020_18_plan0.5.3}{%
\subsection{openfibertargets\_nov2020\_18}\label{openfibertargets_nov2020_18_plan0.5.3}}

\noindent\textbf{target\_selection plan:} 0.5.3

\noindent\textbf{target\_selection tag:}
\href{https://github.com/sdss/target_selection/tree/0.3.4/}{0.3.4}

\noindent\textbf{crossmatch plan:} 0.5.0

\noindent\textbf{Summary:} Hard X-ray emitting sources (primarily AGNs) from
XMM-Newton and Swift serendipitous surveys

\noindent\textbf{Simplified description of selection criteria:} manual


\noindent\textbf{Target priority options:} 6080

\noindent\textbf{Cadence options:} \texttt{bright\_1x1,\ dark\_1x1}

\noindent\textbf{Implementation:} n/a

\noindent\textbf{Number of targets:} 10012

\begin{center}\rule{0.5\linewidth}{0.5pt}\end{center}

\hypertarget{openfibertargets_nov2020_26_plan0.5.3}{%
\subsection{openfibertargets\_nov2020\_26}\label{openfibertargets_nov2020_26_plan0.5.3}}

\noindent\textbf{target\_selection plan:} 0.5.3

\noindent\textbf{target\_selection tag:}
\href{https://github.com/sdss/target_selection/tree/0.3.4/}{0.3.4}

\noindent\textbf{crossmatch plan:} 0.5.0

\noindent\textbf{Summary:} Quasars monitored in SDSS-I and II as a part of
ultraviolet Broad Absorption Line (BAL) Quasar Variability Survey
(\citealt{Gibson2009}) to receive an additional epoch of observations to monitor
wind evolution.

\noindent\textbf{Simplified description of selection criteria:} manual


\noindent\textbf{Target priority options:} 6080

\noindent\textbf{Cadence options:} \texttt{dark\_1x1}

\noindent\textbf{Implementation:} n/a

\noindent\textbf{Number of targets:} 2105

\begin{center}\rule{0.5\linewidth}{0.5pt}\end{center}

\hypertarget{openfibertargets_nov2020_27_plan0.5.3}{%
\subsection{openfibertargets\_nov2020\_27}\label{openfibertargets_nov2020_27_plan0.5.3}}

\noindent\textbf{target\_selection plan:} 0.5.3

\noindent\textbf{target\_selection tag:}
\href{https://github.com/sdss/target_selection/tree/0.3.4/}{0.3.4}

\noindent\textbf{crossmatch plan:} 0.5.0

\noindent\textbf{Summary:} Bright quasar candidates with quasar-like Gaia and
unWISE photometry, supplemented with multi-epoch variability observed by
DES, excluding regions of high confusion

\noindent\textbf{Simplified description of selection criteria:} manual


\noindent\textbf{Target priority options:} 6085

\noindent\textbf{Cadence options:} \texttt{bright\_1x1,\ dark\_1x1}

\noindent\textbf{Implementation:} n/a

\noindent\textbf{Number of targets:} 814963

\begin{center}\rule{0.5\linewidth}{0.5pt}\end{center}

\hypertarget{openfibertargets_nov2020_30_plan0.5.3}{%
\subsection{openfibertargets\_nov2020\_30}\label{openfibertargets_nov2020_30_plan0.5.3}}

\noindent\textbf{target\_selection plan:} 0.5.3

\noindent\textbf{target\_selection tag:}
\href{https://github.com/sdss/target_selection/tree/0.3.4/}{0.3.4}

\noindent\textbf{crossmatch plan:} 0.5.0

\noindent\textbf{Summary:} Bright sources in the JWST North Ecliptic Pole
Time-Domain Field

\noindent\textbf{Simplified description of selection criteria:} manual


\noindent\textbf{Target priority options:} 6080

\noindent\textbf{Cadence options:} \texttt{dark\_1x1}

\noindent\textbf{Implementation:} n/a

\noindent\textbf{Number of targets:} 237

\begin{center}\rule{0.5\linewidth}{0.5pt}\end{center}

\hypertarget{openfibertargets_nov2020_33_plan0.5.3}{%
\subsection{openfibertargets\_nov2020\_33}\label{openfibertargets_nov2020_33_plan0.5.3}}

\noindent\textbf{target\_selection plan:} 0.5.3

\noindent\textbf{target\_selection tag:}
\href{https://github.com/sdss/target_selection/tree/0.3.4/}{0.3.4}

\noindent\textbf{crossmatch plan:} 0.5.0

\noindent\textbf{Summary:} Candidate AGN, QSOs, and Blazars identified based on
their ZTF variability flagged by ALeRCE light curve classifier
(\citealt{SanchezSaez2020})

\noindent\textbf{Simplified description of selection criteria:} manual


\noindent\textbf{Target priority options:} 6080

\noindent\textbf{Cadence options:} \texttt{bright\_1x1,\ dark\_1x1}

\noindent\textbf{Implementation:} n/a

\noindent\textbf{Number of targets:} 22249

\begin{center}\rule{0.5\linewidth}{0.5pt}\end{center}

\hypertarget{openfibertargets_nov2020_11_plan1.0.7}{%
\subsection{openfibertargets\_nov2020\_11}\label{openfibertargets_nov2020_11_plan1.0.7}}

\noindent\textbf{target\_selection plan:} 1.0.7

\noindent\textbf{target\_selection tag:}
\href{https://github.com/sdss/target_selection/tree/1.0.7/}{1.0.7}

\noindent\textbf{crossmatch plan:} 1.0.0

\noindent\textbf{Summary:} Quasars observed during SDSS-I and -II which had only
a single epoch to changes in emission lines and absorption line troughs
over timescales of years to decades

\noindent\textbf{Simplified description of selection criteria:} manual


\noindent\textbf{Target priority options:} 6080

\noindent\textbf{Cadence options:} \texttt{bright\_1x1,\ dark\_1x1}

\noindent\textbf{Implementation:} n/a

\noindent\textbf{Number of targets:} 39684

\begin{center}\rule{0.5\linewidth}{0.5pt}\end{center}

\hypertarget{openfibertargets_nov2020_18_plan1.0.7}{%
\subsection{openfibertargets\_nov2020\_18}\label{openfibertargets_nov2020_18_plan1.0.7}}

\noindent\textbf{target\_selection plan:} 1.0.7

\noindent\textbf{target\_selection tag:}
\href{https://github.com/sdss/target_selection/tree/1.0.7/}{1.0.7}

\noindent\textbf{crossmatch plan:} 1.0.0

\noindent\textbf{Summary:} Hard X-ray emitting sources (primarily AGNs) from
XMM-Newton and Swift serendipitous surveys

\noindent\textbf{Simplified description of selection criteria:} manual


\noindent\textbf{Target priority options:} 6080

\noindent\textbf{Cadence options:} \texttt{bright\_1x1,\ dark\_1x1}

\noindent\textbf{Implementation:} n/a

\noindent\textbf{Number of targets:} 10090

\begin{center}\rule{0.5\linewidth}{0.5pt}\end{center}

\hypertarget{openfibertargets_nov2020_26_plan1.0.7}{%
\subsection{openfibertargets\_nov2020\_26}\label{openfibertargets_nov2020_26_plan1.0.7}}

\noindent\textbf{target\_selection plan:} 1.0.7

\noindent\textbf{target\_selection tag:}
\href{https://github.com/sdss/target_selection/tree/1.0.7/}{1.0.7}

\noindent\textbf{crossmatch plan:} 1.0.0

\noindent\textbf{Summary:} Quasars monitored in SDSS-I and II as a part of
ultraviolet Broad Absorption Line (BAL) Quasar Variability Survey
(\citealt{Gibson2009}) to receive an additional epoch of observations to monitor
wind evolution

\noindent\textbf{Simplified description of selection criteria:} manual


\noindent\textbf{Target priority options:} 6080

\noindent\textbf{Cadence options:} \texttt{dark\_1x1}

\noindent\textbf{Implementation:} n/a

\noindent\textbf{Number of targets:} 2105

\begin{center}\rule{0.5\linewidth}{0.5pt}\end{center}

\hypertarget{openfibertargets_nov2020_27_plan1.0.7}{%
\subsection{openfibertargets\_nov2020\_27}\label{openfibertargets_nov2020_27_plan1.0.7}}

\noindent\textbf{target\_selection plan:} 1.0.7

\noindent\textbf{target\_selection tag:}
\href{https://github.com/sdss/target_selection/tree/1.0.7/}{1.0.7}

\noindent\textbf{crossmatch plan:} 1.0.0

\noindent\textbf{Summary:} Bright quasar candidates with quasar-like Gaia and
unWISE photometry, supplemented with multi-epoch variability observed by
DES, excluding regions of high confusion

\noindent\textbf{Simplified description of selection criteria:} manual


\noindent\textbf{Target priority options:} 6085

\noindent\textbf{Cadence options:} \texttt{bright\_1x1,\ dark\_1x1}

\noindent\textbf{Implementation:} n/a

\noindent\textbf{Number of targets:} 814965

\begin{center}\rule{0.5\linewidth}{0.5pt}\end{center}

\hypertarget{openfibertargets_nov2020_30_plan1.0.7}{%
\subsection{openfibertargets\_nov2020\_30}\label{openfibertargets_nov2020_30_plan1.0.7}}

\noindent\textbf{target\_selection plan:} 1.0.7

\noindent\textbf{target\_selection tag:}
\href{https://github.com/sdss/target_selection/tree/1.0.7/}{1.0.7}

\noindent\textbf{crossmatch plan:} 1.0.0

\noindent\textbf{Summary:} Bright sources in the JWST North Ecliptic Pole
Time-Domain Field

\noindent\textbf{Simplified description of selection criteria:} manual


\noindent\textbf{Target priority options:} 6080

\noindent\textbf{Cadence options:} \texttt{dark\_1x1}

\noindent\textbf{Implementation:} n/a

\noindent\textbf{Number of targets:} 237

\begin{center}\rule{0.5\linewidth}{0.5pt}\end{center}

\hypertarget{openfibertargets_nov2020_33_plan1.0.10}{%
\subsection{openfibertargets\_nov2020\_33}\label{openfibertargets_nov2020_33_plan1.0.10}}

\noindent\textbf{target\_selection plan:} 1.0.10

\noindent\textbf{target\_selection tag:}
\href{https://github.com/sdss/target_selection/tree/1.0.10/}{1.0.10}

\noindent\textbf{crossmatch plan:} 1.0.0

\noindent\textbf{Summary:} Candidate AGN, QSOs, and Blazars identified based on
their ZTF variability flagged by ALeRCE light curve classifier
(\citealt{SanchezSaez2020})

\noindent\textbf{Simplified description of selection criteria:} manual


\noindent\textbf{Target priority options:} 6080

\noindent\textbf{Cadence options:} \texttt{bright\_1x1,\ dark\_1x1}

\noindent\textbf{Implementation:} n/a

\noindent\textbf{Number of targets:} 22246

\begin{center}\rule{0.5\linewidth}{0.5pt}\end{center}

\hypertarget{bhm_rm_core_plan1.0.12}{%
\subsection{bhm\_rm\_core}\label{bhm_rm_core_plan1.0.12}}

\noindent\textbf{target\_selection plan:} 1.0.12

\noindent\textbf{target\_selection tag:}
\href{https://github.com/sdss/target_selection/tree/1.0.12/}{1.0.12}

\noindent\textbf{crossmatch plan:} 1.0.0

\noindent\textbf{Summary:} A sample of candidate QSOs selected via the methods
presented by
\citet{Yang2022}. These targets are located within five (+1 backup) well
known survey fields (SDSS-RM, COSMOS, XMM-LSS, ECDFS, CVZ-S/SEP, and
ELIAS-S1).

\noindent\textbf{Simplified description of selection criteria:} Starting from a
parent catalog of optically selected objects in the RM fields (as
presented by
\citealt{Yang2022}), select candidate QSOs that satisfy all of the
following: i) are identified via the Skew-T algorithm
(\texttt{skewt\_qso\ ==\ 1}); ii) have
17\textless{}\texttt{psfmag\_i}\textless21.5~AB
(16\textless{}\emph{G}\textless21.7~AB in the CVZ-S/SEP field); iii) do
not have significant detections (\textgreater3$\sigma$) of parallax and/or
proper motion in Gaia DR2; iv) are not vetoed due to results of visual
inspections of recent spectroscopy; vi) have detections in all of the
gri bands (a Gaia detection is sufficient in the CVZ-S/SEP field);and
vii) do not lie in the SDSS-RM field.


\noindent\textbf{Target priority options:} 1000-1100

\noindent\textbf{Cadence options:} \texttt{dark\_100x8,\ dark\_174x8}

\noindent\textbf{Implementation:}
\href{https://github.com/sdss/target_selection/blob/1.0.12/python/target_selection/cartons/bhm_rm.py}{bhm\_rm.py}

\noindent\textbf{Number of targets:} 3741

\begin{center}\rule{0.5\linewidth}{0.5pt}\end{center}

\hypertarget{bhm_rm_known_spec_plan1.0.12}{%
\subsection{bhm\_rm\_known\_spec}\label{bhm_rm_known_spec_plan1.0.12}}

\noindent\textbf{target\_selection plan:} 1.0.12

\noindent\textbf{target\_selection tag:}
\href{https://github.com/sdss/target_selection/tree/1.0.12/}{1.0.12}

\noindent\textbf{crossmatch plan:} 1.0.0

\noindent\textbf{Summary:} A sample of known QSOs identified through optical
spectroscopy from various projects, as collated by
\citet{Yang2022}. These targets are located within five (+1 backup) well
known survey fields (SDSS-RM, COSMOS, XMM-LSS, ECDFS, CVZ-S/SEP, and
ELIAS-S1).

\noindent\textbf{Simplified description of selection criteria:} Starting from a
parent catalog of optically selected objects in the RM fields (as
presented by
\citealt{Yang2022}), select targets which satisfy all of the following: i)
are flagged as having a spectroscopic identification (in the parent
catalog); ii) have 15\textless{}\texttt{psfmag\_i}\textless21.7~AB
(SDSS-RM, CDFS, ELIAS-S1 field), 16\textless{}\emph{G}\textless21.7~Vega
(CVZ-S/SEP field), 15\textless{}\texttt{psfmag\_i}\textless21.5~AB
(COSMOS and XMM-LSS fields); iii) have a spectroscopic redshift in the
range 0.005\textless{}\emph{z}\textless7; iv) are not vetoed due to
results of visual inspections of recent spectroscopy.


\noindent\textbf{Target priority options:} 1000-1100

\noindent\textbf{Cadence options:} \texttt{dark\_100x8,\ dark\_174x8}

\noindent\textbf{Implementation:}
\href{https://github.com/sdss/target_selection/blob/1.0.12/python/target_selection/cartons/bhm_rm.py}{bhm\_rm.py}

\noindent\textbf{Number of targets:} 3206

\begin{center}\rule{0.5\linewidth}{0.5pt}\end{center}

\hypertarget{bhm_rm_var_plan1.0.12}{%
\subsection{bhm\_rm\_var}\label{bhm_rm_var_plan1.0.12}}

\noindent\textbf{target\_selection plan:} 1.0.12

\noindent\textbf{target\_selection tag:}
\href{https://github.com/sdss/target_selection/tree/1.0.12/}{1.0.12}

\noindent\textbf{crossmatch plan:} 1.0.0

\noindent\textbf{Summary:} A sample of candidate QSOs selected via their optical
variability properties, as presented by
\citet{Yang2022}. These targets are located within five (+1 backup) well
known survey fields (SDSS-RM, COSMOS, XMM-LSS, ECDFS, CVZ-S/SEP, and
ELIAS-S1).

\noindent\textbf{Simplified description of selection criteria:} Starting from a
parent catalog of optically selected objects in the RM fields (as
presented by
\citealt{Yang2022}), select candidate QSOs that satisfy all of the
following: i) have significant variability in the DES or PanSTARRS1
multi-epoch photometry (\texttt{var\_sn{[}g{]}}\textgreater3 and
\texttt{var\_rms{[}g{]}}\textgreater0.05); ii) have
17\textless{}\texttt{psfmag\_i}\textless20.5~AB
(16\textless{}\emph{G}\textless21.7~AB in the CVZ-S/SEP field); iii) do
not have significant detections (\textgreater3$\sigma$) of parallax and/or
proper motion in Gaia DR2; iv) are not vetoed due to results of visual
inspections of recent spectroscopy; and vi) do not lie in the SDSS-RM
field.


\noindent\textbf{Target priority options:} 1000-1100

\noindent\textbf{Cadence options:} \texttt{dark\_174x8}

\noindent\textbf{Implementation:}
\href{https://github.com/sdss/target_selection/blob/1.0.12/python/target_selection/cartons/bhm_rm.py}{bhm\_rm.py}

\noindent\textbf{Number of targets:} 871

\begin{center}\rule{0.5\linewidth}{0.5pt}\end{center}

\hypertarget{bhm_rm_ancillary_plan1.0.12}{%
\subsection{bhm\_rm\_ancillary}\label{bhm_rm_ancillary_plan1.0.12}}

\noindent\textbf{target\_selection plan:} 1.0.12

\noindent\textbf{target\_selection tag:}
\href{https://github.com/sdss/target_selection/tree/1.0.12/}{1.0.12}

\noindent\textbf{crossmatch plan:} 1.0.0

\noindent\textbf{Summary:} A supporting sample of candidate QSOs which have been
selected by the Gaia-unWISE AGN catalog
(\citealt{Shu2019}) and/or the SDSS XDQSO catalog
(\citealt{Bovy2011}). These targets are located within five (+1 backup) well
known survey fields (SDSS-RM, COSMOS, XMM-LSS, ECDFS, CVZ-S/SEP, and
ELIAS-S1).

\noindent\textbf{Simplified description of selection criteria:} Starting from a
parent catalog of optically selected objects in the RM fields (as
presented by
\citealt{Yang2022}), select candidate QSOs that satisfy all of the
following: i) are identified via external ancillary methods
(\texttt{photo\_bitmask\ \&\ 3\ !=\ 0}); ii) have
15\textless{}\texttt{psfmag\_i}\textless21.5~AB
(16\textless{}\emph{G}\textless21.7~AB in the CVZ-S/SEP field); iii) do
not have significant detections (\textgreater3$\sigma$) of parallax and/or
proper motion in Gaia DR2; iv) are not vetoed due to results of visual
inspections of recent spectroscopy; and v) do not lie in the SDSS-RM
field.


\noindent\textbf{Target priority options:} 1000-1100

\noindent\textbf{Cadence options:} \texttt{dark\_100x8,\ dark\_174x8}

\noindent\textbf{Implementation:}
\href{https://github.com/sdss/target_selection/blob/1.0.12/python/target_selection/cartons/bhm_rm.py}{bhm\_rm.py}

\noindent\textbf{Number of targets:} 1114

\begin{center}\rule{0.5\linewidth}{0.5pt}\end{center}

\hypertarget{bhm_rm_xrayqso_plan1.0.12}{%
\subsection{bhm\_rm\_xrayqso}\label{bhm_rm_xrayqso_plan1.0.12}}

\noindent\textbf{target\_selection plan:} 1.0.12

\noindent\textbf{target\_selection tag:}
\href{https://github.com/sdss/target_selection/tree/1.0.12/}{1.0.12}

\noindent\textbf{crossmatch plan:} 1.0.0

\noindent\textbf{Summary:} A small supporting sample of candidate QSOs in two of
the SDSS-V Reverberation Mapping fields (CDFS, ELIAS-S1), identified via
X-ray detection and SED modelling by
\citealt{Ni2021}.

\noindent\textbf{Simplified description of selection criteria:} Starting from the
sample of candidate QSOs identified by
\citealt{Ni2021}, select all that satisfy the following criteria: i) are
\textless1.5~deg from the centre of at least one RM field, ii) have
optical magnitudes in the range
16.5\textless{}\texttt{mag\_i}\textless21.5, iii) are not vetoed due to
results of visual inspections of recent spectroscopy.


\noindent\textbf{Target priority options:} 1000-1100

\noindent\textbf{Cadence options:} \texttt{dark\_174x8}

\noindent\textbf{Implementation:}
\href{https://github.com/sdss/target_selection/blob/1.0.12/python/target_selection/cartons/bhm_rm.py}{bhm\_rm.py}

\noindent\textbf{Number of targets:} 43

\begin{center}\rule{0.5\linewidth}{0.5pt}\end{center}

\hypertarget{bhm_csc_boss_plan1.0.37}{%
\subsection{bhm\_csc\_boss}\label{bhm_csc_boss_plan1.0.37}}

\noindent\textbf{target\_selection plan:} 1.0.37

\noindent\textbf{target\_selection tag:}
\href{https://github.com/sdss/target_selection/tree/1.0.37/}{1.0.37}

\noindent\textbf{crossmatch plan:} 1.0.0

\noindent\textbf{Summary:} X-ray sources from the CSC2.1 source catalog with
counterparts in LegacySurvey DR10, Panstarrs1-DR1 or Gaia DR3.

\noindent\textbf{Simplified description of selection criteria:} Starting from the
parent catalog of CSC sources with optical/IR counterparts
(\texttt{bhm\_csc\_v3}). Select entries satisfying the following
criteria: i) optical counterpart is from the LegacySurvey DR10,
PanSTARRS1 or Gaia DR3 catalogs, ii) optical flux/magnitude is in the
accepted range for SDSS-V: all \texttt{fibermag\_{[}g,r,i,z{]}}
\textgreater14.5 and at least one \texttt{fibermag\_{[}g,r,i,z{]}}
\textless21.0~AB (objects with LegacySurvey DR10 counterparts); all
\texttt{psfmag\_{[}g,r,i,z{]}} \textgreater{} 14.0 and at least one
\texttt{psfmag\_{[}g,r,i,z{]}} \textless{} 22.0~AB (objects with
PanSTARRS1 counterparts); 14.0 \textless{} \emph{G}\textless20.0~Vega
(Gaia DR3 counterparts). De-prioritize targets which already have good
quality SDSS spectroscopy. Allocate cadence (exposure time) requests
based on optical brightness (LegacySurvey \texttt{fibermag\_r}, PS1
\texttt{psfmag\_i} or Gaia \emph{G}).


\noindent\textbf{Target priority options:} 1920-1939, 2920-2939

\noindent\textbf{Cadence options:}
\texttt{bright\_1x1,\ dark\_flexible\_2x1,\ dark\_flexible\_2x2}

\noindent\textbf{Implementation:}
\href{https://github.com/sdss/target_selection/blob/1.0.37/python/target_selection/cartons/bhm_csc.py}{bhm\_csc.py}

\noindent\textbf{Number of targets:} 135154

\begin{center}\rule{0.5\linewidth}{0.5pt}\end{center}

\hypertarget{bhm_csc_apogee_plan1.0.37}{%
\subsection{bhm\_csc\_apogee}\label{bhm_csc_apogee_plan1.0.37}}

\noindent\textbf{target\_selection plan:} 1.0.37

\noindent\textbf{target\_selection tag:}
\href{https://github.com/sdss/target_selection/tree/1.0.37/}{1.0.37}

\noindent\textbf{crossmatch plan:} 1.0.0

\noindent\textbf{Summary:} X-ray sources from the CSC2.1 source catalog with NIR
counterparts in 2MASS PSC

\noindent\textbf{Simplified description of selection criteria:} Starting from the
parent catalog of CSC sources with optical/IR counterparts
(\texttt{bhm\_csc\_v2}). Select entries satisfying the following
criteria: i) NIR counterpart is from the 2MASS catalog, ii) 2MASS
\emph{H}-band magnitude measurement is not null and is within a
reasonable range for SDSS-V: 7.0\textless{}\emph{H}\textless14.0.
Allocate cadence (exposure time) requests based on \emph{H}-band
magnitude.


\noindent\textbf{Target priority options:} 2930-2939

\noindent\textbf{Cadence options:} \texttt{bright\_1x1,\ bright\_flexible\_3x1}

\noindent\textbf{Implementation:}
\href{https://github.com/sdss/target_selection/blob/1.0.37/python/target_selection/cartons/bhm_csc.py}{bhm\_csc.py}

\noindent\textbf{Number of targets:} 54593

\begin{center}\rule{0.5\linewidth}{0.5pt}\end{center}

\hypertarget{bhm_gua_dark_plan1.0.37}{%
\subsection{bhm\_gua\_dark}\label{bhm_gua_dark_plan1.0.37}}

\noindent\textbf{target\_selection plan:} 1.0.37

\noindent\textbf{target\_selection tag:}
\href{https://github.com/sdss/target_selection/tree/1.0.37/}{1.0.37}

\noindent\textbf{crossmatch plan:} 1.0.0

\noindent\textbf{Summary:} A sample of optically faint candidate AGN lacking
spectroscopic confirmations, derived from the parent sample presented by
\citet{Shu2019}, who applied a machine-learning approach to select QSO
candidates from a combination of the Gaia DR2 and unWISE catalogs.

\noindent\textbf{Simplified description of selection criteria:} Starting with the
\citet{Shu2019} catalog, select targets which satisfy the following
criteria: i) have a Random Forest probability of being a QSO
of\textgreater0.8, ii) are in the magnitude range suitable for BOSS
spectroscopy in dark time (\emph{G}\textgreater16.5 and
\emph{RP}\textgreater16.5, as well as \emph{G}\textless21.2 or
\emph{RP}\textless21.0,~Vega mag), iii) do not yet have good SDSS
optical spectroscopic measurements. Note that the selection criteria are
based on apparent magnitudes, rather than the dereddened magnitudes that
were used in an earlier iteration of this carton.


\noindent\textbf{Target priority options:} 3400

\noindent\textbf{Cadence options:} \texttt{dark\_flexible\_2x2}

\noindent\textbf{Implementation:}
\href{https://github.com/sdss/target_selection/blob/1.0.37/python/target_selection/cartons/bhm_gua.py}{bhm\_gua.py}

\noindent\textbf{Number of targets:} 1936784

\begin{center}\rule{0.5\linewidth}{0.5pt}\end{center}

\hypertarget{bhm_gua_bright_plan1.0.37}{%
\subsection{bhm\_gua\_bright}\label{bhm_gua_bright_plan1.0.37}}

\noindent\textbf{target\_selection plan:} 1.0.37

\noindent\textbf{target\_selection tag:}
\href{https://github.com/sdss/target_selection/tree/1.0.37/}{1.0.37}

\noindent\textbf{crossmatch plan:} 1.0.0

\noindent\textbf{Summary:} A sample of optically bright candidate AGN lacking
spectroscopic confirmations, derived from the parent sample presented by
\citet{Shu2019}, who applied a machine-learning approach to select QSO
candidates from a combination of the Gaia DR2 and unWISE catalogs.

\noindent\textbf{Simplified description of selection criteria:} Starting with the
\citet{Shu2019} catalog, select targets which satisfy the following
criteria: i) have a Random Forest probability of being a QSO
of\textgreater0.8, ii) are in the magnitude range suitable for BOSS
spectroscopy in bright time (\emph{G}\textgreater13.0 and
\emph{RP}\textgreater13.5, as well as \emph{G}\textless18.5 or
\emph{RP}\textless18.5,~Vega mag), iii) do not yet have good SDSS
optical spectroscopic measurements. Note that the selection criteria are
based on apparent magnitudes, rather than the dereddened magnitudes that
were used in an earlier iteration of this carton.


\noindent\textbf{Target priority options:} 4040

\noindent\textbf{Cadence options:} \texttt{bright\_flexible\_2x1}

\noindent\textbf{Implementation:}
\href{https://github.com/sdss/target_selection/blob/1.0.37/python/target_selection/cartons/bhm_gua.py}{bhm\_gua.py}

\noindent\textbf{Number of targets:} 185644

\begin{center}\rule{0.5\linewidth}{0.5pt}\end{center}

\hypertarget{bhm_spiders_clusters_lsdr10_plan1.0.37}{%
\subsection{bhm\_spiders\_clusters\_lsdr10}\label{bhm_spiders_clusters_lsdr10_plan1.0.37}}

\noindent\textbf{target\_selection plan:} 1.0.37

\noindent\textbf{target\_selection tag:}
\href{https://github.com/sdss/target_selection/tree/1.0.37/}{1.0.37}

\noindent\textbf{crossmatch plan:} 1.0.0

\noindent\textbf{Summary:} This is the main carton for SPIDERS galaxy cluster
wide area follow up. The carton provides a catalog of galaxies which are
candidate members of clusters selected from early reductions of the
first 24-months of eROSITA all sky survey data (eRASS:4). The X-ray
clusters have been associated by the eROSITA-DE team to
\href{https://www.legacysurvey.org/dr10/}{legacysurvey.org/dr10}
optical/IR counterparts using the eROMAPPER red-sequence finder
algorithm
(\citealt{Rykoff2014};
\citealt{IderChitham2020}). All targets are located in the sky hemisphere
where MPE controls the data rights (approx.
180\textless{}\emph{l}\textless360~deg). Targets in this carton are
located at high Galactic latitudes
\textbar{}\emph{b}\textbar\textgreater20~deg.

\noindent\textbf{Simplified description of selection criteria:} Starting from a
parent catalog of eRASS:4 $\rightarrow$ legacysurvey.org/dr10 eROMAPPER cluster
associations, select targets which meet all of the following criteria:
i) have 13.5\textless{}\texttt{fibertotmag\_i}\textless21.0 or
13.5\textless{}\texttt{fibertotmag\_z}\textless20.5~AB, ii) if detected
by Gaia DR3 then have \emph{G}\textgreater13.5 and
\emph{RP}\textgreater13.5~Vega, and iii) do not have existing good
quality SDSS spectroscopy. We assign a range of priorities to targets in
this carton, with Brightest Cluster Galaxies (BCGs) top ranked, followed
by candidate member galaxies according their probability of membership
and whether they are within the 'core'
16\textless{}\texttt{fibertotmag\_i}\textless20.5 or
16\textless{}\texttt{fibertotmag\_z}\textless20.0~AB). We assign
cadences (exposure time requests) based on optical brightness.


\noindent\textbf{Target priority options:} 1501, 1631-1634, 3501, 3631-3635

\noindent\textbf{Cadence options:}
\texttt{bright\_flexible\_2x1,\ dark\_flexible\_2x1,\ dark\_flexible\_2x2}

\noindent\textbf{Implementation:}
\href{https://github.com/sdss/target_selection/blob/1.0.37/python/target_selection/cartons/bhm_spiders_clusters.py}{bhm\_spiders\_clusters.py}

\noindent\textbf{Number of targets:} 440276

\begin{center}\rule{0.5\linewidth}{0.5pt}\end{center}

\hypertarget{bhm_spiders_agn_lsdr10_plan1.0.37}{%
\subsection{bhm\_spiders\_agn\_lsdr10}\label{bhm_spiders_agn_lsdr10_plan1.0.37}}

\noindent\textbf{target\_selection plan:} 1.0.37

\noindent\textbf{target\_selection tag:}
\href{https://github.com/sdss/target_selection/tree/1.0.37/}{1.0.37}

\noindent\textbf{crossmatch plan:} 1.0.0

\noindent\textbf{Summary:} This is the highest priority carton for SPIDERS AGN
wide area follow up. The carton provides optical counterparts to
point-like (unresolved) X-ray sources detected in early reductions of
the first 18~months of eROSITA all sky survey data (eRASS:3). The sample
is expected to contain a mixture of QSOs, AGN, stars and compact
objects. The X-ray sources have been cross-matched by the eROSITA-DE
team to \href{https://www.legacysurvey.org/dr10/}{legacysurvey.org/DR10}
optical/IR counterparts (supplemented with the DR9 catalog at
Dec\textgreater+32.375~deg). All targets are located in the sky
hemisphere where MPE controls the data rights (approx.
180\textless{}\emph{l}\textless360~deg). Due to the LegacySurvey
footprint, nearly all targets in this carton are located at high
Galactic latitudes \textbar{}\emph{b}\textbar\textgreater15~deg.

\noindent\textbf{Simplified description of selection criteria:} Starting from a
parent catalog of eRASS:3 point source $\rightarrow$ legacysurvey.org/DR10(+DR9)
associations (method: NWAY assisted by optical/IR priors computed via a
pre-trained Random Forest, building on
\citealt{Salvato2022}), select targets which meet all of the following criteria:
i) have eROSITA detection likelihood\textgreater6.0, ii) have an X-ray $\rightarrow$
optical/IR cross-match probability (NWAY) of
\texttt{p\_any}\textgreater0.1, iii) have
13.5\textless{}\texttt{fibermag\_r}\textless22.5 or
13.5\textless{}\texttt{fibermag\_i}\textless22.3 or
13.5\textless{}\texttt{fibermag\_z}\textless21.5, iv) are not saturated
in LegacySurvey imaging, v) if detected by Gaia DR2/DR3 then have
\emph{G}\textgreater13.5 and \emph{RP}\textgreater13.5~Vega. We
\emph{deprioritise} targets if any of the following criteria are met: a)
the target already has existing good quality SDSS spectroscopy, b) the
X-ray detection likelihood is \textless8.0, c) the target is a secondary
X-ray$\rightarrow$optical/IR association, d) the target falls outside the core
magnitude range (i.e. doesn't satisfy one of
16.5\textless{}\texttt{fibermag\_r}\textless22.0 or
16.5\textless{}\texttt{fibermag\_i}\textless21.8 or
16.5\textless{}\texttt{fibermag\_z}\textless21.0), e) the target lies at
\textbar{}\emph{b}\textbar\textless15~deg or at Dec\textless-75~deg, f)
the X-ray flux is
\textless2x10\textsuperscript{-14}~erg~s\textsuperscript{-1}~cm\textsuperscript{-2}.
We slightly boost the priorities of targets which have hard X-ray
detections. Cadence choices (exposure time requests per target) are
based on optical brightness.


\noindent\textbf{Target priority options:} 1520-1527, 1720-1727, 3520-3527,
3720-3727

\noindent\textbf{Cadence options:}
\texttt{bright\_flexible\_2x1,\ dark\_flexible\_2x1,\ dark\_flexible\_2x2}

\noindent\textbf{Implementation:}
\href{https://github.com/sdss/target_selection/blob/1.0.37/python/target_selection/cartons/bhm_spiders_agn.py}{bhm\_spiders\_agn.py}

\noindent\textbf{Number of targets:} 975978

\begin{center}\rule{0.5\linewidth}{0.5pt}\end{center}

\hypertarget{bhm_spiders_agn_hard_plan1.0.37}{%
\subsection{bhm\_spiders\_agn\_hard}\label{bhm_spiders_agn_hard_plan1.0.37}}

\noindent\textbf{target\_selection plan:} 1.0.37

\noindent\textbf{target\_selection tag:}
\href{https://github.com/sdss/target_selection/tree/1.0.37/}{1.0.37}

\noindent\textbf{crossmatch plan:} 1.0.0

\noindent\textbf{Summary:} This small supplementary carton provides a sample of
point-like (unresolved) hard-band (2.3-5~keV) X-ray sources detected in
the first 6~months of eROSITA all sky survey data (eRASS:1). The
selection of the eROSITA hard band sample is described by
\citealt{Waddell2026}. The X-ray sources have been cross-matched by the
eROSITA-DE team to
\href{https://www.legacysurvey.org/dr10/}{LegacySurvey DR10} optical/IR
counterparts. All targets are located in the sky hemisphere where MPE
controls the data rights (approx.
180\textless{}\emph{l}\textless360~deg).

\noindent\textbf{Simplified description of selection criteria:} Starting from a
parent catalog of multi-band eRASS:1 X-ray point source $\rightarrow$ LegacySurvey
DR10 associations (method: NWAY assisted by optical/IR priors computed
via a pre-trained Random Forest, building on
\citealt{Salvato2022}), select targets which meet all of the following criteria:
i) have eROSITA 2.3-5keV band detection likelihood\textgreater12.0, ii)
have an X-ray $\rightarrow$ optical/IR cross-match probability (NWAY) of
\texttt{p\_any}\textgreater0.01, iii) have
13.5\textless{}\texttt{fibermag\_r}\textless22.5 or
13.5\textless{}\texttt{fibermag\_i}\textless22.3 or
13.5\textless{}\texttt{fibermag\_z}\textless21.5, iv) are not saturated
in LegacySurvey imaging, v) if detected by Gaia DR2/DR3 then have
\emph{G}\textgreater13.5 and \emph{RP}\textgreater13.5~Vega. We
\emph{deprioritise} targets if any of the following criteria are met: a)
the target already has existing good quality SDSS spectroscopy, b) the
target is a secondary X-ray$\rightarrow$optical/IR association, c) the target falls
outside the core magnitude range (i.e. doesn't satisfy one of
16.5\textless{}\texttt{fibermag\_r}\textless22.0 or
16.5\textless{}\texttt{fibermag\_i}\textless21.8 or
16.5\textless{}\texttt{fibermag\_z}\textless21.0), d) the target lies at
\textbar{}\emph{b}\textbar\textless10~deg or at Dec\textless-75~deg.
Cadence choices (exposure time requests per target) are based on optical
brightness.


\noindent\textbf{Target priority options:} 1510-1512, 1710, 3510-3512, 3710

\noindent\textbf{Cadence options:}
\texttt{bright\_flexible\_2x1,\ dark\_flexible\_2x1,\ dark\_flexible\_2x2}

\noindent\textbf{Implementation:}
\href{https://github.com/sdss/target_selection/blob/1.0.37/python/target_selection/cartons/bhm_spiders_agn.py}{bhm\_spiders\_agn.py}

\noindent\textbf{Number of targets:} 2603

\begin{center}\rule{0.5\linewidth}{0.5pt}\end{center}

\hypertarget{bhm_spiders_agn_gaiadr3_plan1.0.37}{%
\subsection{bhm\_spiders\_agn\_gaiadr3}\label{bhm_spiders_agn_gaiadr3_plan1.0.37}}

\noindent\textbf{target\_selection plan:} 1.0.37

\noindent\textbf{target\_selection tag:}
\href{https://github.com/sdss/target_selection/tree/1.0.37/}{1.0.37}

\noindent\textbf{crossmatch plan:} 1.0.0

\noindent\textbf{Summary:} This is the second highest priority carton for SPIDERS
AGN wide area follow up, it is included to expand the survey footprint
beyond the LegacySurvey footprint. This carton provides optical
counterparts to point-like (unresolved) X-ray sources detected in early
reductions of the first 18~months of eROSITA all sky survey data
(eRASS:3). The sample is expected to contain a mixture of QSOs, AGN,
stars and compact objects. The X-ray sources have been cross-matched by
the eROSITA-DE team to Gaia DR3 optical counterparts. All targets are
located in the sky hemisphere where MPE controls the data rights
(approx. 180\textless{}\emph{l}\textless360~deg). The targets in this
carton are distributed over a wide range of Galactic latitudes, but
targets at low Galactic latitudes
\textbar{}\emph{b}\textbar\textless15~deg do not drive survey strategy.

\noindent\textbf{Simplified description of selection criteria:} Starting from a
parent catalog of eRASS:3 point source $\rightarrow$ Gaia DR3 associations (method:
NWAY assisted by Gaia photometric+astrometric priors computed via a
pre-trained Random Forest, building on
\citealt{Salvato2022}), select targets which meet all of the following criteria:
i) have eROSITA detection likelihood\textgreater6.0, ii) have an
X-ray$\rightarrow$Gaia cross-match probability of \texttt{p\_any}\textgreater0.1,
iii) have \emph{G}\textgreater13.5 and \emph{RP}\textgreater13.5~Vega.
We deprioritise targets if any of the following criteria are met: a) the
target already has existing good quality SDSS spectroscopy, b) the X-ray
detection likelihood is \textless8.0, c) the target is a secondary
X-ray$\rightarrow$Gaia association, d) the target falls outside the core magnitude
range (i.e. doesn't satisfy 16.0\textless{}\emph{G}\textless21.0), e)
the target lies at \textbar{}\emph{b}\textbar\textless15~deg or at
Dec\textless-75~deg, f) the X-ray flux is
\textless2x10\textsuperscript{-14}~erg~s\textsuperscript{-1}~cm\textsuperscript{-2}.
We slightly boost the priorities of targets which have hard X-ray
detections. Cadence choices (exposure time requests per target) are
based on optical brightness.


\noindent\textbf{Target priority options:} 1540-1547, 1740-1747, 3540-3547,
3740-3747

\noindent\textbf{Cadence options:}
\texttt{bright\_flexible\_2x1,\ dark\_flexible\_2x1,\ dark\_flexible\_2x2}

\noindent\textbf{Implementation:}
\href{https://github.com/sdss/target_selection/blob/1.0.37/python/target_selection/cartons/bhm_spiders_agn.py}{bhm\_spiders\_agn.py}

\noindent\textbf{Number of targets:} 717805

\begin{center}\rule{0.5\linewidth}{0.5pt}\end{center}

\hypertarget{bhm_spiders_agn_tda_plan1.0.37}{%
\subsection{bhm\_spiders\_agn\_tda}\label{bhm_spiders_agn_tda_plan1.0.37}}

\noindent\textbf{target\_selection plan:} 1.0.37

\noindent\textbf{target\_selection tag:}
\href{https://github.com/sdss/target_selection/tree/1.0.37/}{1.0.37}

\noindent\textbf{crossmatch plan:} 1.0.0

\noindent\textbf{Summary:} This is a small supplementary carton of
transient/variable X-ray sources selected by comparing X-ray
measurements between several independent 6-month eROSITA sky surveys.
All targets are located in the sky hemisphere where MPE controls the
data rights (approx. 180\textless{}\emph{l}\textless360~deg).

\noindent\textbf{Simplified description of selection criteria:} We start from a
parent catalog of transient/variable X-ray sources selected from the
first three eROSITA All Sky surveys. These X-ray sources have been
associated with LegacySurvey DR10 optical counterparts (method: NWAY
assisted by optical/IR priors computed via a pre-trained Random Forest,
building on
\citealt{Salvato2022}). We then select targets which meet all of the following
criteria: i) have 13.5\textless{}\texttt{fibermag\_r}\textless22.5 or
13.5\textless{}\texttt{fibermag\_i}\textless22.3 or
13.5\textless{}\texttt{fibermag\_z}\textless21.5, ii) if detected by
Gaia DR2/DR3 then have \emph{G}\textgreater13.5 and
\emph{RP}\textgreater13.5~Vega. Cadence choices (exposure time requests
per target) are based on optical brightness.


\noindent\textbf{Target priority options:} 1690

\noindent\textbf{Cadence options:}
\texttt{bright\_flexible\_2x1,\ dark\_flexible\_2x1,\ dark\_flexible\_2x2}

\noindent\textbf{Implementation:}
\href{https://github.com/sdss/target_selection/blob/1.0.37/python/target_selection/cartons/bhm_spiders_agn.py}{bhm\_spiders\_agn.py}

\noindent\textbf{Number of targets:} 9627

\begin{center}\rule{0.5\linewidth}{0.5pt}\end{center}

\hypertarget{bhm_spiders_agn_sep_plan1.0.37}{%
\subsection{bhm\_spiders\_agn\_sep}\label{bhm_spiders_agn_sep_plan1.0.37}}

\noindent\textbf{target\_selection plan:} 1.0.37

\noindent\textbf{target\_selection tag:}
\href{https://github.com/sdss/target_selection/tree/1.0.37/}{1.0.37}

\noindent\textbf{crossmatch plan:} 1.0.0

\noindent\textbf{Summary:} This is a small supplementary catalog of eROSITA
point-like (unresolved) X-ray sources detected in the Continous Viewing
Zone South/South Ecliptic Pole (CVZ-S/SEP) field. Owing to the deepest
exposure of the eROSITA All-Sky Survey (eRASS) in this region, the
CVZ-S/SEP field suffers from severe source confusion, requiring a
specialized source-detection procedure (Liu et al., in prep.). The X-ray
catalog used here was derived from an early data reduction of
0.2-2.3~keV band data from the 2.3-year eRASS:5 dataset, covering an
approximately 6~deg~x~6~deg sky region centered on the CVZ-S/SEP. The
X-ray sources have been cross-matched by the eROSITA-DE team with Gaia
DR3.

\noindent\textbf{Simplified description of selection criteria:} Starting from a
parent catalog of eRASS:5 point source $\rightarrow$ Gaia DR3 associations from teh
SEP/CVZ-S region, select targets which meet all of the following
criteria: i) have eROSITA detection likelihood\textgreater6.0, ii) have
an X-ray$\rightarrow$Gaia cross-match probability of \texttt{p\_any}\textgreater0.1,
iii) have \emph{G}\textgreater13.5 and \emph{RP}\textgreater13.5~Vega.
We deprioritise targets if any of the following criteria are met: a) the
target already has existing good quality SDSS spectroscopy, b) the X-ray
detection likelihood is \textless8.0, c) the target is a secondary
X-ray$\rightarrow$Gaia association, d) the target falls outside the core magnitude
range (i.e. doesn't satisfy 16.0\textless{}\emph{G}\textless21.5), e)
the X-ray flux is
\textless1x10\textsuperscript{-14}~erg~s\textsuperscript{-1}~cm\textsuperscript{-2}.
Cadence choices (exposure time requests per target) are based on optical
brightness.


\noindent\textbf{Target priority options:} 1501-1503, 3501-3507

\noindent\textbf{Cadence options:}
\texttt{bright\_flexible\_2x1,\ dark\_flexible\_2x1,\ dark\_flexible\_2x2}

\noindent\textbf{Implementation:}
\href{https://github.com/sdss/target_selection/blob/1.0.37/python/target_selection/cartons/bhm_spiders_agn.py}{bhm\_spiders\_agn.py}

\noindent\textbf{Number of targets:} 1681

\begin{center}\rule{0.5\linewidth}{0.5pt}\end{center}

\hypertarget{bhm_aqmes_med_plan1.0.37}{%
\subsection{bhm\_aqmes\_med}\label{bhm_aqmes_med_plan1.0.37}}

\noindent\textbf{target\_selection plan:} 1.0.37

\noindent\textbf{target\_selection tag:}
\href{https://github.com/sdss/target_selection/tree/1.0.37/}{1.0.37}

\noindent\textbf{crossmatch plan:} 1.0.0

\noindent\textbf{Summary:} Spectroscopically confirmed optically bright SDSS
QSOs, selected from the SDSS QSO catalog (DR16Q,
\citealt{Lyke2020}). Located in 36 mostly disjoint fields within the SDSS QSO
footprint that were pre-selected to contain higher than average numbers
of bright QSOs and CSC targets. The list of field centres can be found
within
\href{https://github.com/sdss/target_selection/blob/0.3.0/python/target_selection/masks/candidate_target_fields_bhm_aqmes_med_v0.3.1.fits}{the
target\_selection repository}.

\noindent\textbf{Simplified description of selection criteria:} Select all
objects from SDSS DR16 QSO catalog that have
16.0\textless{}\texttt{sdss\_psfmag\_i}\textless19.1~AB, that lie within
1.49~degrees of at least one AQMES-medium field location.


\noindent\textbf{Target priority options:} 1100

\noindent\textbf{Cadence options:} \texttt{dark\_10x4\_4yr}

\noindent\textbf{Implementation:}
\href{https://github.com/sdss/target_selection/blob/1.0.37/python/target_selection/cartons/bhm_aqmes.py}{bhm\_aqmes.py}

\noindent\textbf{Number of targets:} 2663

\begin{center}\rule{0.5\linewidth}{0.5pt}\end{center}

\hypertarget{bhm_aqmes_med_faint_plan1.0.37}{%
\subsection{bhm\_aqmes\_med\_faint}\label{bhm_aqmes_med_faint_plan1.0.37}}

\noindent\textbf{target\_selection plan:} 1.0.37

\noindent\textbf{target\_selection tag:}
\href{https://github.com/sdss/target_selection/tree/1.0.37/}{1.0.37}

\noindent\textbf{crossmatch plan:} 1.0.0

\noindent\textbf{Summary:} Spectroscopically confirmed optically faint SDSS QSOs,
selected from the SDSS QSO catalog (DR16Q,
\citealt{Lyke2020}). Located in 36 mostly disjoint fields within the SDSS QSO
footprint that were pre-selected to contain higher than average numbers
of bright QSOs and CSC targets. The list of field centres can be found
within
\href{https://github.com/sdss/target_selection/blob/0.3.0/python/target_selection/masks/candidate_target_fields_bhm_aqmes_med_v0.3.1.fits}{the
target\_selection repository}.

\noindent\textbf{Simplified description of selection criteria:} Select all
objects from SDSS DR16 QSO catalog that have
19.1\textless{}\texttt{sdss\_psfmag\_i}\textless21.0~AB, that lie within
1.49~degrees of at least one AQMES-medium field location.


\noindent\textbf{Target priority options:} 3100

\noindent\textbf{Cadence options:} \texttt{dark\_10x4\_4yr}

\noindent\textbf{Implementation:}
\href{https://github.com/sdss/target_selection/blob/1.0.37/python/target_selection/cartons/bhm_aqmes.py}{bhm\_aqmes.py}

\noindent\textbf{Number of targets:} 16853

\begin{center}\rule{0.5\linewidth}{0.5pt}\end{center}

\hypertarget{bhm_aqmes_wide2_plan1.0.37}{%
\subsection{bhm\_aqmes\_wide2}\label{bhm_aqmes_wide2_plan1.0.37}}

\noindent\textbf{target\_selection plan:} 1.0.37

\noindent\textbf{target\_selection tag:}
\href{https://github.com/sdss/target_selection/tree/1.0.37/}{1.0.37}

\noindent\textbf{crossmatch plan:} 1.0.0

\noindent\textbf{Summary:} Spectroscopically confirmed optically bright SDSS
QSOs, selected from the SDSS QSO catalog (DR16Q,
\citealt{Lyke2020}). Located in 425 fields within the SDSS QSO footprint,
where the choice of survey area prioritized field that overlapped with
the SPIDERS footprint (approx 180\textless{}\emph{b}\textless360~deg),
and/or had higher than average numbers of bright QSOs and CSC targets.
The list of field centres can be found
\href{https://github.com/sdss/target_selection/blob/0.3.0/python/target_selection/masks/candidate_target_fields_bhm_aqmes_wide_v0.3.1.fits}{within
the target\_selection repository}. The targets in this carton request
two epochs of SDSS-V spectroscopy.

\noindent\textbf{Simplified description of selection criteria:} Select all
objects from SDSS DR16 QSO catalog that have
16.0\textless{}\texttt{sdss\_psfmag\_i}\textless19.1~AB, and that lie
within 1.49~degrees of at least one AQMES-wide field location.


\noindent\textbf{Target priority options:} 1210

\noindent\textbf{Cadence options:} \texttt{dark\_2x4}

\noindent\textbf{Implementation:}
\href{https://github.com/sdss/target_selection/blob/1.0.37/python/target_selection/cartons/bhm_aqmes.py}{bhm\_aqmes.py}

\noindent\textbf{Number of targets:} 24142

\begin{center}\rule{0.5\linewidth}{0.5pt}\end{center}

\hypertarget{bhm_aqmes_wide2_faint_plan1.0.37}{%
\subsection{bhm\_aqmes\_wide2\_faint}\label{bhm_aqmes_wide2_faint_plan1.0.37}}

\noindent\textbf{target\_selection plan:} 1.0.37

\noindent\textbf{target\_selection tag:}
\href{https://github.com/sdss/target_selection/tree/1.0.37/}{1.0.37}

\noindent\textbf{crossmatch plan:} 1.0.0

\noindent\textbf{Summary:} Spectroscopically confirmed optically faint SDSS QSOs,
selected from the SDSS QSO catalog (DR16Q,
\citealt{Lyke2020}). Located in 425 fields within the SDSS QSO footprint,
where the choice of survey area prioritized field that overlapped with
the SPIDERS footprint (approx 180\textless{}\emph{b}\textless360~deg),
and/or had higher than average numbers of bright QSOs and CSC targets.
The list of field centres can be found
\href{https://github.com/sdss/target_selection/blob/0.3.0/python/target_selection/masks/candidate_target_fields_bhm_aqmes_wide_v0.3.1.fits}{within
the target\_selection repository}. The targets in this carton request
two epochs of SDSS-V spectroscopy.

\noindent\textbf{Simplified description of selection criteria:} Select all
objects from SDSS DR16 QSO catalog that have
19.1\textless{}\texttt{sdss\_psfmag\_i}\textless21.0~AB, and that lie
within 1.49~degrees of at least one AQMES-wide field location.


\noindent\textbf{Target priority options:} 3210

\noindent\textbf{Cadence options:} \texttt{dark\_2x4}

\noindent\textbf{Implementation:}
\href{https://github.com/sdss/target_selection/blob/1.0.37/python/target_selection/cartons/bhm_aqmes.py}{bhm\_aqmes.py}

\noindent\textbf{Number of targets:} 99586

\begin{center}\rule{0.5\linewidth}{0.5pt}\end{center}

\hypertarget{bhm_aqmes_bonus_core_plan1.0.37}{%
\subsection{bhm\_aqmes\_bonus\_core}\label{bhm_aqmes_bonus_core_plan1.0.37}}

\noindent\textbf{target\_selection plan:} 1.0.37

\noindent\textbf{target\_selection tag:}
\href{https://github.com/sdss/target_selection/tree/1.0.37/}{1.0.37}

\noindent\textbf{crossmatch plan:} 1.0.0

\noindent\textbf{Summary:} Spectroscopically confirmed optically bright SDSS
QSOs, selected from the SDSS QSO catalog (DR16Q,
\citealt{Lyke2020}). Located anywhere within the SDSS DR16Q footprint.

\noindent\textbf{Simplified description of selection criteria:} Select all
objects from SDSS DR16 QSO catalog that have
16.0\textless{}\texttt{sdss\_psfmag\_i}\textless19.1~AB


\noindent\textbf{Target priority options:} 3300

\noindent\textbf{Cadence options:} \texttt{dark\_flexible\_2x2}

\noindent\textbf{Implementation:}
\href{https://github.com/sdss/target_selection/blob/1.0.37/python/target_selection/cartons/bhm_aqmes.py}{bhm\_aqmes.py}

\noindent\textbf{Number of targets:} 83163

\begin{center}\rule{0.5\linewidth}{0.5pt}\end{center}

\hypertarget{bhm_aqmes_bonus_bright_plan1.0.37}{%
\subsection{bhm\_aqmes\_bonus\_bright}\label{bhm_aqmes_bonus_bright_plan1.0.37}}

\noindent\textbf{target\_selection plan:} 1.0.37

\noindent\textbf{target\_selection tag:}
\href{https://github.com/sdss/target_selection/tree/1.0.37/}{1.0.37}

\noindent\textbf{crossmatch plan:} 1.0.0

\noindent\textbf{Summary:} Spectroscopically confirmed, extremely optically
bright SDSS QSOs, selected from the SDSS QSO catalog (DR16Q,
\citealt{Lyke2020}). Located anywhere within the SDSS DR16Q footprint.

\noindent\textbf{Simplified description of selection criteria:} Select all
objects from SDSS DR16 QSO catalog that have
14.0\textless{}\texttt{sdss\_psfmag\_i}\textless18.0~AB


\noindent\textbf{Target priority options:} 4040

\noindent\textbf{Cadence options:} \texttt{bright\_flexible\_2x1}

\noindent\textbf{Implementation:}
\href{https://github.com/sdss/target_selection/blob/1.0.37/python/target_selection/cartons/bhm_aqmes.py}{bhm\_aqmes.py}

\noindent\textbf{Number of targets:} 10848

\begin{center}\rule{0.5\linewidth}{0.5pt}\end{center}

\hypertarget{bhm_aqmes_bonus_faint_plan1.0.37}{%
\subsection{bhm\_aqmes\_bonus\_faint}\label{bhm_aqmes_bonus_faint_plan1.0.37}}

\noindent\textbf{target\_selection plan:} 1.0.37

\noindent\textbf{target\_selection tag:}
\href{https://github.com/sdss/target_selection/tree/1.0.37/}{1.0.37}

\noindent\textbf{crossmatch plan:} 1.0.0

\noindent\textbf{Summary:} Spectroscopically confirmed optically faint SDSS QSOs,
selected from the SDSS QSO catalog (DR16Q,
\citealt{Lyke2020}). Located anywhere within the SDSS DR16Q footprint.

\noindent\textbf{Simplified description of selection criteria:} Select all
objects from SDSS DR16 QSO catalog that have
19.1\textless{}\texttt{sdss\_psfmag\_i}\textless21.0~AB


\noindent\textbf{Target priority options:} 3302

\noindent\textbf{Cadence options:} \texttt{dark\_flexible\_2x2}

\noindent\textbf{Implementation:}
\href{https://github.com/sdss/target_selection/blob/1.0.37/python/target_selection/cartons/bhm_aqmes.py}{bhm\_aqmes.py}

\noindent\textbf{Number of targets:} 424163

\begin{center}\rule{0.5\linewidth}{0.5pt}\end{center}

\hypertarget{bhm_colr_galaxies_lsdr10_plan1.0.38}{%
\subsection{bhm\_colr\_galaxies\_lsdr10}\label{bhm_colr_galaxies_lsdr10_plan1.0.38}}

\noindent\textbf{target\_selection plan:} 1.0.38

\noindent\textbf{target\_selection tag:}
\href{https://github.com/sdss/target_selection/tree/1.0.38/}{1.0.38}

\noindent\textbf{crossmatch plan:} 1.0.0

\noindent\textbf{Summary:} A supplementary magnitude limited sample of optically
bright galaxies selected from the LegacySurvey DR10 optical/IR imaging
catalog (supplemented with the LegacySurvey DR9 catalog at
Dec\textgreater+32.375~deg). Selection is based on optical morphology,
lack of Gaia parallax, and several magnitude cuts.

\noindent\textbf{Simplified description of selection criteria:} Starting from the
\texttt{legacy\_survey\_dr10} catalog (lsdr10), select entries
satisfying all of the following criteria: i) lsdr10 morphological
\texttt{type} != 'PSF', ii) lsdr10
\texttt{shape\_r}\textgreater1.0~arcsec, iii) no detection of parallax
by Gaia, iv) dereddened \emph{z}-band model mag\textless19.0~AB, and
dereddened \emph{z}-band fiber mag\textless19.5~AB, v) within all the
following apparent fiber mag limits: 16\textless{}\emph{g}\textless22.5~
and 16\textless{}\emph{r}\textless21.5~ and
16\textless{}\emph{z}\textless20.0~AB, vi) if detected in Gaia, then
\emph{G}\textgreater15.0~ and \emph{RP}\textgreater15.0~Vega, vii) at
moderate to high Galactic latitude
\textbar{}\emph{b}\textbar\textgreater15~deg.


\noindent\textbf{Target priority options:} 7100

\noindent\textbf{Cadence options:}
\texttt{bright\_1x1,\ dark\_1x1,\ dark\_flexible\_2x2}

\noindent\textbf{Implementation:}
\href{https://github.com/sdss/target_selection/blob/1.0.38/python/target_selection/cartons/bhm_galaxies.py}{bhm\_galaxies.py}

\noindent\textbf{Number of targets:} 6246703

\begin{center}\rule{0.5\linewidth}{0.5pt}\end{center}

\hypertarget{openfibertargets_bhm_racsradio_boss_plan1.2.0}{%
\subsection{openfibertargets\_bhm\_racsradio\_boss}\label{openfibertargets_bhm_racsradio_boss_plan1.2.0}}

\noindent\textbf{target\_selection plan:} 1.2.0

\noindent\textbf{target\_selection tag:}
\href{https://github.com/sdss/target_selection/tree/1.2.7/}{1.2.7}

\noindent\textbf{crossmatch plan:} 1.0.0

\noindent\textbf{Summary:} Optically bright counterparts to bright 888~MHz radio
sources identified by in the ASKAP/RACS-low survey
(\citealt{Hale2021}).

\noindent\textbf{Simplified description of selection criteria:} Subset of radio
sources from the RACS-low ASKAP survey (entire sky region between
declination -80 deg and +30~deg), with radio (888~MHz) flux exceeding
3~mJy and optical counterparts (with
14.5\textless{}\emph{r}\textless20.5) identified in the footprints of
LegacySurvey DR10-South using NWAY (adopting separation, positional
errors, number density and a prior based on various features).


\noindent\textbf{Target priority options:} 6085

\noindent\textbf{Cadence options:} \texttt{bright\_1x1,\ dark\_1x1}

\noindent\textbf{Implementation:} n/a

\noindent\textbf{Number of targets:} 149989

\begin{center}\rule{0.5\linewidth}{0.5pt}\end{center}

\hypertarget{openfibertargets_bhm_varagn_boss_plan1.2.0}{%
\subsection{openfibertargets\_bhm\_varagn\_boss}\label{openfibertargets_bhm_varagn_boss_plan1.2.0}}

\noindent\textbf{target\_selection plan:} 1.2.0

\noindent\textbf{target\_selection tag:}
\href{https://github.com/sdss/target_selection/tree/1.2.7/}{1.2.7}

\noindent\textbf{crossmatch plan:} 1.0.0

\noindent\textbf{Summary:} Optical variability selected AGN candidates. We have
used custom made forced photometry in ZTF difference images to build
accurate light curves of complicated objects such as low redshift AGN,
where the variable AGN point source is contaminated by the host galaxy
in a PSF dependent way. With these light curves and complementary color
information we built a random forest algorithm to classify all objects,
obtaining 350k AGN candidates, some of which fall outside the usual
locus of AGN in color-color diagrams.

\noindent\textbf{Simplified description of selection criteria:} All ZTF
forced-photometry based AGN candidates that are not included in the
Gaia-unWISE color-selected AGN candidate catalog of
\citet{Shu2019} and that: a) have -30\textless Dec\textless15 and
\textbar{}\emph{b}\textbar\textgreater20~deg, have zero proper motion
according to Gaia (at 2 sigma), are not in the labeled set of
well-characterized objects used in
\citealt{SanchezSaez2023}, and are classified as either low-z AGN, mid-z AGN or
Blazar by our classifier. For candidates classified as high-z AGN the
same limits apply except that in this case
\textbar{}\emph{b}\textbar\textgreater30~deg and the Gaia proper motion
must be within 1 sigma of zero.


\noindent\textbf{Target priority options:} 6085

\noindent\textbf{Cadence options:} \texttt{bright\_1x1,\ dark\_1x1}

\noindent\textbf{Implementation:} n/a

\noindent\textbf{Number of targets:} 6966

\begin{center}\rule{0.5\linewidth}{0.5pt}\end{center}

\hypertarget{openfibertargets_bhm_morexray_boss_plan1.2.0}{%
\subsection{openfibertargets\_bhm\_morexray\_boss}\label{openfibertargets_bhm_morexray_boss_plan1.2.0}}

\noindent\textbf{target\_selection plan:} 1.2.0

\noindent\textbf{target\_selection tag:}
\href{https://github.com/sdss/target_selection/tree/1.2.7/}{1.2.7}

\noindent\textbf{crossmatch plan:} 1.0.0

\noindent\textbf{Summary:} Counterparts of X-ray sources identified by Swift-XRT
and XMM-Newton

\noindent\textbf{Simplified description of selection criteria:} Select sources
from ExSeSS (Swift,
\citealt{Delaney2023}) and
\href{http://xmmssc.irap.omp.eu/Catalog/4XMM-DR13/4XMM_DR13.html}{4XMM-DR13}
(\citealt{Webb2020}), in any X-ray band. Cross-match to LegacySurvey DR10
counterparts (using NWAY informed by WISE \emph{W1-W2} colors, with
\texttt{p\_any}\textgreater0.3), and then apply magnitude cuts
(14.5\textless{}\emph{r}\textless20~AB for bright time and
20\textless{}\emph{r}\textless21~AB for dark time)


\noindent\textbf{Target priority options:} 6085

\noindent\textbf{Cadence options:} \texttt{bright\_1x1,\ dark\_1x1}

\noindent\textbf{Implementation:} n/a

\noindent\textbf{Number of targets:} 103105

\begin{center}\rule{0.5\linewidth}{0.5pt}\end{center}

\hypertarget{openfibertargets_bhm_hightderates_psb_boss_plan1.2.0}{%
\subsection{openfibertargets\_bhm\_hightderates\_psb\_boss}\label{openfibertargets_bhm_hightderates_psb_boss_plan1.2.0}}

\noindent\textbf{target\_selection plan:} 1.2.0

\noindent\textbf{target\_selection tag:}
\href{https://github.com/sdss/target_selection/tree/1.2.7/}{1.2.7}

\noindent\textbf{crossmatch plan:} 1.0.0

\noindent\textbf{Summary:} Galaxies that have been photometrically identified as
or post-starburst (E+A). Such galaxies are known to have an elevated
rate of tidal disruption events of stars by their supermassive black
holes.

\noindent\textbf{Simplified description of selection criteria:} All or
post-starburst galaxies identified by the machine learning algorithm in
\citealt{French2018}


\noindent\textbf{Target priority options:} 6085

\noindent\textbf{Cadence options:} \texttt{dark\_1x1}

\noindent\textbf{Implementation:} n/a

\noindent\textbf{Number of targets:} 8934

\begin{center}\rule{0.5\linewidth}{0.5pt}\end{center}

\hypertarget{openfibertargets_bhm_hightderates_qbs_boss_plan1.2.0}{%
\subsection{openfibertargets\_bhm\_hightderates\_qbs\_boss}\label{openfibertargets_bhm_hightderates_qbs_boss_plan1.2.0}}

\noindent\textbf{target\_selection plan:} 1.2.0

\noindent\textbf{target\_selection tag:}
\href{https://github.com/sdss/target_selection/tree/1.2.7/}{1.2.7}

\noindent\textbf{crossmatch plan:} 1.0.0

\noindent\textbf{Summary:} Galaxies that have been photometrically identified as
quiescent Balmer-strong (QBS). Such galaxies are known to have an
elevated rate of tidal disruption events of stars by their supermassive
black holes.

\noindent\textbf{Simplified description of selection criteria:} All quiescent
Balmer-strong galaxies identified by the machine learning algorithm in
\citealt{French2018}


\noindent\textbf{Target priority options:} 6085

\noindent\textbf{Cadence options:} \texttt{dark\_1x1}

\noindent\textbf{Implementation:} n/a

\noindent\textbf{Number of targets:} 29659

\begin{center}\rule{0.5\linewidth}{0.5pt}\end{center}

\hypertarget{openfibertargets_bhm_quaia_boss_plan1.2.0}{%
\subsection{openfibertargets\_bhm\_quaia\_boss}\label{openfibertargets_bhm_quaia_boss_plan1.2.0}}

\noindent\textbf{target\_selection plan:} 1.2.0

\noindent\textbf{target\_selection tag:}
\href{https://github.com/sdss/target_selection/tree/1.2.7/}{1.2.7}

\noindent\textbf{crossmatch plan:} 1.0.0

\noindent\textbf{Summary:} Quasar candidates from the Quaia Catalog (the
Gaia--unWISE Quasar Catalog) that have not previously been
spectroscopically observed by any generation of SDSS

\noindent\textbf{Simplified description of selection criteria:} All Quaia
G\textless20.5 quasars not previously spectroscopically observed by SDSS


\noindent\textbf{Target priority options:} 6085

\noindent\textbf{Cadence options:} \texttt{bright\_1x1}

\noindent\textbf{Implementation:} n/a

\noindent\textbf{Number of targets:} 999906

\begin{center}\rule{0.5\linewidth}{0.5pt}\end{center}

\hypertarget{openfibertargets_bhm_lofarradio_boss_plan1.2.1}{%
\subsection{openfibertargets\_bhm\_lofarradio\_boss}\label{openfibertargets_bhm_lofarradio_boss_plan1.2.1}}

\noindent\textbf{target\_selection plan:} 1.2.1

\noindent\textbf{target\_selection tag:}
\href{https://github.com/sdss/target_selection/tree/1.3.0/}{1.3.0}

\noindent\textbf{crossmatch plan:} 1.0.0

\noindent\textbf{Summary:} Optically bright counterparts of low-frequency radio
sources identified by LOFAR Two Metre Sky Survey (LOTSS DR2,
\citealt{Hardcastle2023})

\noindent\textbf{Simplified description of selection criteria:} Extract LOFAR
radio detected sources from the LOTSS survey with published matches to
LegacySurvey DR8 counterparts
(\citealt{Hardcastle2023}). Apply a radio flux cut
(\texttt{Total\ flux}\textgreater2~mJy), and apply magnitude limits
(14.5\textless{}\emph{r}\textless20.0)


\noindent\textbf{Target priority options:} 6085

\noindent\textbf{Cadence options:} \texttt{bright\_1x1}

\noindent\textbf{Implementation:} n/a

\noindent\textbf{Number of targets:} 210483

\begin{center}\rule{0.5\linewidth}{0.5pt}\end{center}

\hypertarget{bhm_csc_boss_d3_plan1.2.4}{%
\subsection{bhm\_csc\_boss\_d3}\label{bhm_csc_boss_d3_plan1.2.4}}

\noindent\textbf{target\_selection plan:} 1.2.4

\noindent\textbf{target\_selection tag:}
\href{https://github.com/sdss/target_selection/tree/1.3.3/}{1.3.3}

\noindent\textbf{crossmatch plan:} 1.0.0

\noindent\textbf{Summary:} X-ray sources from the CSC2.1 source catalog with
counterparts in LegacySurvey DR10, Panstarrs1-DR1 or Gaia DR3.

\noindent\textbf{Simplified description of selection criteria:} Starting from the
parent catalog of CSC sources with optical/IR counterparts
(\texttt{bhm\_csc\_v3}). Select entries satisfying the following
criteria: i) optical counterpart is from the LegacySurvey DR10,
PanSTARRS1 or Gaia DR3 catalogs, ii) optical flux/magnitude is in the
accepted range for SDSS-V: all \texttt{fibermag\_{[}g,r,i,z{]}}
\textgreater14.5 and at least one \texttt{fibermag\_{[}g,r,i,z{]}}
\textless21.0~AB (objects with LegacySurvey DR10 counterparts); all
\texttt{psfmag\_{[}g,r,i,z{]}} \textgreater{} 14.0 and at least one
\texttt{psfmag\_{[}g,r,i,z{]}} \textless{} 22.0~AB (objects with
PanSTARRS1 counterparts); 14.0 \textless{} \emph{G}\textless20.0~Vega
(Gaia DR3 counterparts). De-prioritize targets which already have good
quality SDSS spectroscopy. The '\_d3' variation of this carton provides
a shallower cadence option for targets toward the optically faint end of
the sample.


\noindent\textbf{Target priority options:} 1921-1930

\noindent\textbf{Cadence options:} \texttt{dark\_flexible\_3x1}

\noindent\textbf{Implementation:}
\href{https://github.com/sdss/target_selection/blob/1.3.3/python/target_selection/cartons/bhm_csc.py}{bhm\_csc.py}

\noindent\textbf{Number of targets:} 97267

\begin{center}\rule{0.5\linewidth}{0.5pt}\end{center}

\hypertarget{bhm_spiders_clusters_lsdr10_d3_plan1.2.4}{%
\subsection{bhm\_spiders\_clusters\_lsdr10\_d3}\label{bhm_spiders_clusters_lsdr10_d3_plan1.2.4}}

\noindent\textbf{target\_selection plan:} 1.2.4

\noindent\textbf{target\_selection tag:}
\href{https://github.com/sdss/target_selection/tree/1.3.3/}{1.3.3}

\noindent\textbf{crossmatch plan:} 1.0.0

\noindent\textbf{Summary:} This is the main carton for SPIDERS galaxy cluster
wide area follow up. The carton provides a catalog of galaxies which are
candidate members of clusters selected from early reductions of the
first 24-months of eROSITA all sky survey data (eRASS:4). The X-ray
clusters have been associated by the eROSITA-DE team to
\href{https://www.legacysurvey.org/dr10/}{legacysurvey.org/dr10}
optical/IR counterparts using the eROMAPPER red-sequence finder
algorithm
(\citealt{Rykoff2014};
\citealt{IderChitham2020}). All targets are located in the sky hemisphere
where MPE controls the data rights (approx.
180\textless{}\emph{l}\textless360~deg). Targets in this carton are
located at high Galactic latitudes
\textbar{}\emph{b}\textbar\textgreater20~deg.

\noindent\textbf{Simplified description of selection criteria:} Starting from a
parent catalog of eRASS:4 $\rightarrow$ legacysurvey.org/dr10 eROMAPPER cluster
associations, select targets which meet all of the following criteria:
i) have 13.5\textless{}\texttt{fibertotmag\_i}\textless21.0 or
13.5\textless{}\texttt{fibertotmag\_z}\textless20.5~AB, ii) if detected
by Gaia DR3 then have \emph{G}\textgreater13.5 and
\emph{RP}\textgreater13.5~Vega, and iii) do not have existing good
quality SDSS spectroscopy. We assign a range of priorities to targets in
this carton, with Brightest Cluster Galaxies (BCGs) top ranked, followed
by candidate member galaxies according their probability of membership
and whether they are within the 'core'
16\textless{}\texttt{fibertotmag\_i}\textless20.5 or
16\textless{}\texttt{fibertotmag\_z}\textless20.0~AB). The '\_d3'
variation of this carton provides a shallower cadence option for targets
toward the optically faint end of the sample.


\noindent\textbf{Target priority options:} 1502, 1632-1635, 3502, 3632-3636

\noindent\textbf{Cadence options:} \texttt{dark\_flexible\_3x1}

\noindent\textbf{Implementation:}
\href{https://github.com/sdss/target_selection/blob/1.3.3/python/target_selection/cartons/bhm_spiders_clusters.py}{bhm\_spiders\_clusters.py}

\noindent\textbf{Number of targets:} 427445

\begin{center}\rule{0.5\linewidth}{0.5pt}\end{center}

\hypertarget{bhm_spiders_agn_gaiadr3_d3_plan1.2.4}{%
\subsection{bhm\_spiders\_agn\_gaiadr3\_d3}\label{bhm_spiders_agn_gaiadr3_d3_plan1.2.4}}

\noindent\textbf{target\_selection plan:} 1.2.4

\noindent\textbf{target\_selection tag:}
\href{https://github.com/sdss/target_selection/tree/1.3.3/}{1.3.3}

\noindent\textbf{crossmatch plan:} 1.0.0

\noindent\textbf{Summary:} This is the second highest priority carton for SPIDERS
AGN wide area follow up, it is included to expand the survey footprint
beyond the LegacySurvey footprint. This carton provides optical
counterparts to point-like (unresolved) X-ray sources detected in early
reductions of the first 18~months of eROSITA all sky survey data
(eRASS:3). The sample is expected to contain a mixture of QSOs, AGN,
stars and compact objects. The X-ray sources have been cross-matched by
the eROSITA-DE team to Gaia DR3 optical counterparts. All targets are
located in the sky hemisphere where MPE controls the data rights
(approx. 180\textless{}\emph{l}\textless360~deg). The targets in this
carton are distributed over a wide range of Galactic latitudes, but
targets at low Galactic latitudes
\textbar{}\emph{b}\textbar\textless15~deg do not drive survey strategy.

\noindent\textbf{Simplified description of selection criteria:} Starting from a
parent catalog of eRASS:3 point source $\rightarrow$ Gaia DR3 associations (method:
NWAY assisted by Gaia photometric+astrometric priors computed via a
pre-trained Random Forest, building on
\citealt{Salvato2022}), select targets which meet all of the following criteria:
i) have eROSITA detection likelihood\textgreater6.0, ii) have an
X-ray$\rightarrow$Gaia cross-match probability of \texttt{p\_any}\textgreater0.1,
iii) have \emph{G}\textgreater13.5 and \emph{RP}\textgreater13.5~Vega.
We deprioritise targets if any of the following criteria are met: a) the
target already has existing good quality SDSS spectroscopy, b) the X-ray
detection likelihood is \textless8.0, c) the target is a secondary
X-ray$\rightarrow$Gaia association, d) the target falls outside the core magnitude
range (i.e. doesn't satisfy 16.0\textless{}\emph{G}\textless21.0), e)
the target lies at \textbar{}\emph{b}\textbar\textless15~deg or at
Dec\textless-75~deg, f) the X-ray flux is
\textless2x10\textsuperscript{-14}~erg~s\textsuperscript{-1}~cm\textsuperscript{-2}.
We slightly boost the priorities of targets which have hard X-ray
detections. The '\_d3' variation of this carton provides a shallower
cadence option for targets toward the optically faint end of the sample.


\noindent\textbf{Target priority options:} 1541-1548, 3541-3548

\noindent\textbf{Cadence options:} \texttt{dark\_flexible\_3x1}

\noindent\textbf{Implementation:}
\href{https://github.com/sdss/target_selection/blob/1.3.3/python/target_selection/cartons/bhm_spiders_agn.py}{bhm\_spiders\_agn.py}

\noindent\textbf{Number of targets:} 398493

\begin{center}\rule{0.5\linewidth}{0.5pt}\end{center}

\hypertarget{bhm_spiders_agn_sep_d3_plan1.2.4}{%
\subsection{bhm\_spiders\_agn\_sep\_d3}\label{bhm_spiders_agn_sep_d3_plan1.2.4}}

\noindent\textbf{target\_selection plan:} 1.2.4

\noindent\textbf{target\_selection tag:}
\href{https://github.com/sdss/target_selection/tree/1.3.3/}{1.3.3}

\noindent\textbf{crossmatch plan:} 1.0.0

\noindent\textbf{Summary:} This is a small supplementary catalog of eROSITA
point-like (unresolved) X-ray sources detected in the Continous Viewing
Zone South/South Ecliptic Pole (CVZ-S/SEP) field. Owing to the deepest
exposure of the eROSITA All-Sky Survey (eRASS) in this region, the
CVZ-S/SEP field suffers from severe source confusion, requiring a
specialized source-detection procedure (Liu et al., in prep.). The X-ray
catalog used here was derived from an early data reduction of
0.2-2.3~keV band data from the 2.3-year eRASS:5 dataset, covering an
approximately 6~deg~x~6~deg sky region centered on the CVZ-S/SEP. The
X-ray sources have been cross-matched by the eROSITA-DE team with Gaia
DR3.

\noindent\textbf{Simplified description of selection criteria:} Starting from a
parent catalog of eRASS:5 point source $\rightarrow$ Gaia DR3 associations from teh
SEP/CVZ-S region, select targets which meet all of the following
criteria: i) have eROSITA detection likelihood\textgreater6.0, ii) have
an X-ray$\rightarrow$Gaia cross-match probability of \texttt{p\_any}\textgreater0.1,
iii) have \emph{G}\textgreater13.5 and \emph{RP}\textgreater13.5~Vega.
We deprioritise targets if any of the following criteria are met: a) the
target already has existing good quality SDSS spectroscopy, b) the X-ray
detection likelihood is \textless8.0, c) the target is a secondary
X-ray$\rightarrow$Gaia association, d) the target falls outside the core magnitude
range (i.e. doesn't satisfy 16.0\textless{}\emph{G}\textless21.5), e)
the X-ray flux is
\textless1x10\textsuperscript{-14}~erg~s\textsuperscript{-1}~cm\textsuperscript{-2}.
Cadence choices (exposure time requests per target) are based on optical
brightness. The '\_d3' variation of this carton provides a shallower
cadence option for targets toward the optically faint end of the sample.


\noindent\textbf{Target priority options:} 1502-1504, 3502-3508

\noindent\textbf{Cadence options:} \texttt{dark\_flexible\_3x1}

\noindent\textbf{Implementation:}
\href{https://github.com/sdss/target_selection/blob/1.3.3/python/target_selection/cartons/bhm_spiders_agn.py}{bhm\_spiders\_agn.py}

\noindent\textbf{Number of targets:} 863

\begin{center}\rule{0.5\linewidth}{0.5pt}\end{center}

\hypertarget{bhm_aqmes_wide1_plan1.2.4}{%
\subsection{bhm\_aqmes\_wide1}\label{bhm_aqmes_wide1_plan1.2.4}}

\noindent\textbf{target\_selection plan:} 1.2.4

\noindent\textbf{target\_selection tag:}
\href{https://github.com/sdss/target_selection/tree/1.3.3/}{1.3.3}

\noindent\textbf{crossmatch plan:} 1.0.0

\noindent\textbf{Summary:} Spectroscopically confirmed optically bright SDSS
QSOs, selected from the SDSS QSO catalog (DR16Q,
\citealt{Lyke2020}). Located in 425 fields within the SDSS QSO footprint,
where the choice of survey area prioritized field that overlapped with
the SPIDERS footprint (approx 180\textless{}\emph{b}\textless360~deg),
and/or had higher than average numbers of bright QSOs and CSC targets.
The list of field centres can be found
\href{https://github.com/sdss/target_selection/blob/0.3.0/python/target_selection/masks/candidate_target_fields_bhm_aqmes_wide_v0.3.1.fits}{within
the target\_selection repository}. The targets in this carton request a
single epoch of SDSS-V spectroscopy.

\noindent\textbf{Simplified description of selection criteria:} Select all
objects from SDSS DR16 QSO catalog that have
16.0\textless{}\texttt{sdss\_psfmag\_i}\textless19.1~AB, and that lie
within 1.49~degrees of at least one AQMES-wide field location.


\noindent\textbf{Target priority options:} 1211

\noindent\textbf{Cadence options:} \texttt{dark\_1x4}

\noindent\textbf{Implementation:}
\href{https://github.com/sdss/target_selection/blob/1.3.3/python/target_selection/cartons/bhm_aqmes.py}{bhm\_aqmes.py}

\noindent\textbf{Number of targets:} 24142

\begin{center}\rule{0.5\linewidth}{0.5pt}\end{center}

\hypertarget{bhm_aqmes_wide1_faint_plan1.2.4}{%
\subsection{bhm\_aqmes\_wide1\_faint}\label{bhm_aqmes_wide1_faint_plan1.2.4}}

\noindent\textbf{target\_selection plan:} 1.2.4

\noindent\textbf{target\_selection tag:}
\href{https://github.com/sdss/target_selection/tree/1.3.3/}{1.3.3}

\noindent\textbf{crossmatch plan:} 1.0.0

\noindent\textbf{Summary:} Spectroscopically confirmed optically faint SDSS QSOs,
selected from the SDSS QSO catalog (DR16Q,
\citealt{Lyke2020}). Located in 425 fields within the SDSS QSO footprint,
where the choice of survey area prioritized field that overlapped with
the SPIDERS footprint (approx 180\textless{}\emph{b}\textless360~deg),
and/or had higher than average numbers of bright QSOs and CSC targets.
The list of field centres can be found
\href{https://github.com/sdss/target_selection/blob/0.3.0/python/target_selection/masks/candidate_target_fields_bhm_aqmes_wide_v0.3.1.fits}{within
the target\_selection repository}. The targets in this carton request a
single epoch of SDSS-V spectroscopy.

\noindent\textbf{Simplified description of selection criteria:} Select all
objects from SDSS DR16 QSO catalog that have
19.1\textless{}\texttt{sdss\_psfmag\_i}\textless21.0~AB, and that lie
within 1.49~degrees of at least one AQMES-wide field location.


\noindent\textbf{Target priority options:} 3211

\noindent\textbf{Cadence options:} \texttt{dark\_1x4}

\noindent\textbf{Implementation:}
\href{https://github.com/sdss/target_selection/blob/1.3.3/python/target_selection/cartons/bhm_aqmes.py}{bhm\_aqmes.py}

\noindent\textbf{Number of targets:} 99586

\begin{center}\rule{0.5\linewidth}{0.5pt}\end{center}

\hypertarget{bhm_gua_dark_d3_plan1.2.4}{%
\subsection{bhm\_gua\_dark\_d3}\label{bhm_gua_dark_d3_plan1.2.4}}

\noindent\textbf{target\_selection plan:} 1.2.4

\noindent\textbf{target\_selection tag:}
\href{https://github.com/sdss/target_selection/tree/1.3.3/}{1.3.3}

\noindent\textbf{crossmatch plan:} 1.0.0

\noindent\textbf{Summary:} A sample of optically faint candidate AGN lacking
spectroscopic confirmations, derived from the parent sample presented by
\citet{Shu2019}, who applied a machine-learning approach to select QSO
candidates from a combination of the Gaia DR2 and unWISE catalogs.

\noindent\textbf{Simplified description of selection criteria:} Starting with the
\citet{Shu2019} catalog, select targets which satisfy the following
criteria: i) have a Random Forest probability of being a QSO
of\textgreater0.8, ii) are in the magnitude range suitable for BOSS
spectroscopy in dark time (\emph{G}\textgreater16.5 and
\emph{RP}\textgreater16.5, as well as \emph{G}\textless21.2 or
\emph{RP}\textless21.0,~Vega mag), iii) do not yet have good SDSS
optical spectroscopic measurements. Note that the selection criteria are
based on apparent magnitudes, rather than the dereddened magnitudes that
were used in an earlier iteration of this carton. The '\_d3' variation
of this carton provides a shallower cadence option for targets in this
sample.


\noindent\textbf{Target priority options:} 3401

\noindent\textbf{Cadence options:} \texttt{dark\_flexible\_3x1}

\noindent\textbf{Implementation:}
\href{https://github.com/sdss/target_selection/blob/1.3.3/python/target_selection/cartons/bhm_gua.py}{bhm\_gua.py}

\noindent\textbf{Number of targets:} 1936784

\begin{center}\rule{0.5\linewidth}{0.5pt}\end{center}

\hypertarget{bhm_spiders_agn_hard_d3_plan1.2.4}{%
\subsection{bhm\_spiders\_agn\_hard\_d3}\label{bhm_spiders_agn_hard_d3_plan1.2.4}}

\noindent\textbf{target\_selection plan:} 1.2.4

\noindent\textbf{target\_selection tag:}
\href{https://github.com/sdss/target_selection/tree/1.3.3/}{1.3.3}

\noindent\textbf{crossmatch plan:} 1.0.0

\noindent\textbf{Summary:} This small supplementary carton provides a sample of
point-like (unresolved) hard-band (2.3-5~keV) X-ray sources detected in
the first 6~months of eROSITA all sky survey data (eRASS:1). The
selection of the eROSITA hard band sample is described by
\citealt{Waddell2026}. The X-ray sources have been cross-matched by the
eROSITA-DE team to
\href{https://www.legacysurvey.org/dr10/}{LegacySurvey DR10} optical/IR
counterparts. All targets are located in the sky hemisphere where MPE
controls the data rights (approx.
180\textless{}\emph{l}\textless360~deg).

\noindent\textbf{Simplified description of selection criteria:} Starting from a
parent catalog of multi-band eRASS:1 X-ray point source $\rightarrow$ LegacySurvey
DR10 associations (method: NWAY assisted by optical/IR priors computed
via a pre-trained Random Forest, building on
\citealt{Salvato2022}), select targets which meet all of the following criteria:
i) have eROSITA 2.3-5keV band detection likelihood\textgreater12.0, ii)
have an X-ray $\rightarrow$ optical/IR cross-match probability (NWAY) of
\texttt{p\_any}\textgreater0.01, iii) have
13.5\textless{}\texttt{fibermag\_r}\textless22.5 or
13.5\textless{}\texttt{fibermag\_i}\textless22.3 or
13.5\textless{}\texttt{fibermag\_z}\textless21.5, iv) are not saturated
in LegacySurvey imaging, v) if detected by Gaia DR2/DR3 then have
\emph{G}\textgreater13.5 and \emph{RP}\textgreater13.5~Vega. We
\emph{deprioritise} targets if any of the following criteria are met: a)
the target already has existing good quality SDSS spectroscopy, b) the
target is a secondary X-ray$\rightarrow$optical/IR association, c) the target falls
outside the core magnitude range (i.e. doesn't satisfy one of
16.5\textless{}\texttt{fibermag\_r}\textless22.0 or
16.5\textless{}\texttt{fibermag\_i}\textless21.8 or
16.5\textless{}\texttt{fibermag\_z}\textless21.0), d) the target lies at
\textbar{}\emph{b}\textbar\textless10~deg or at Dec\textless-75~deg.
Cadence choices (exposure time requests per target) are based on optical
brightness. The '\_d3' variation of this carton provides a shallower
cadence option for targets in this sample.


\noindent\textbf{Target priority options:} 1511-1513, 1711, 3511-3513, 3711

\noindent\textbf{Cadence options:} \texttt{dark\_flexible\_3x1}

\noindent\textbf{Implementation:}
\href{https://github.com/sdss/target_selection/blob/1.3.3/python/target_selection/cartons/bhm_spiders_agn.py}{bhm\_spiders\_agn.py}

\noindent\textbf{Number of targets:} 798

\begin{center}\rule{0.5\linewidth}{0.5pt}\end{center}

\hypertarget{bhm_spiders_agn_tda_d3_plan1.2.4}{%
\subsection{bhm\_spiders\_agn\_tda\_d3}\label{bhm_spiders_agn_tda_d3_plan1.2.4}}

\noindent\textbf{target\_selection plan:} 1.2.4

\noindent\textbf{target\_selection tag:}
\href{https://github.com/sdss/target_selection/tree/1.3.3/}{1.3.3}

\noindent\textbf{crossmatch plan:} 1.0.0

\noindent\textbf{Summary:} This is a small supplementary carton of
transient/variable X-ray sources selected by comparing X-ray
measurements between several independent 6-month eROSITA sky surveys.
All targets are located in the sky hemisphere where MPE controls the
data rights (approx. 180\textless{}\emph{l}\textless360~deg).

\noindent\textbf{Simplified description of selection criteria:} We start from a
parent catalog of transient/variable X-ray sources selected from the
first three eROSITA All Sky surveys. These X-ray sources have been
associated with LegacySurvey DR10 optical counterparts (method: NWAY
assisted by optical/IR priors computed via a pre-trained Random Forest,
building on
\citealt{Salvato2022}). We then select targets which meet all of the following
criteria: i) have 13.5\textless{}\texttt{fibermag\_r}\textless22.5 or
13.5\textless{}\texttt{fibermag\_i}\textless22.3 or
13.5\textless{}\texttt{fibermag\_z}\textless21.5, ii) if detected by
Gaia DR2/DR3 then have \emph{G}\textgreater13.5 and
\emph{RP}\textgreater13.5~Vega. Cadence choices (exposure time requests
per target) are based on optical brightness. The '\_d3' variation of
this carton provides a shallower cadence option for targets in this
sample.


\noindent\textbf{Target priority options:} 1691

\noindent\textbf{Cadence options:} \texttt{dark\_flexible\_3x1}

\noindent\textbf{Implementation:}
\href{https://github.com/sdss/target_selection/blob/1.3.3/python/target_selection/cartons/bhm_spiders_agn.py}{bhm\_spiders\_agn.py}

\noindent\textbf{Number of targets:} 8141

\begin{center}\rule{0.5\linewidth}{0.5pt}\end{center}

\hypertarget{bhm_spiders_agn_lsdr10_d3_plan1.2.4}{%
\subsection{bhm\_spiders\_agn\_lsdr10\_d3}\label{bhm_spiders_agn_lsdr10_d3_plan1.2.4}}

\noindent\textbf{target\_selection plan:} 1.2.4

\noindent\textbf{target\_selection tag:}
\href{https://github.com/sdss/target_selection/tree/1.3.3/}{1.3.3}

\noindent\textbf{crossmatch plan:} 1.0.0

\noindent\textbf{Summary:} This is the highest priority carton for SPIDERS AGN
wide area follow up. The carton provides optical counterparts to
point-like (unresolved) X-ray sources detected in early reductions of
the first 18~months of eROSITA all sky survey data (eRASS:3). The sample
is expected to contain a mixture of QSOs, AGN, stars and compact
objects. The X-ray sources have been cross-matched by the eROSITA-DE
team to \href{https://www.legacysurvey.org/dr10/}{legacysurvey.org/DR10}
optical/IR counterparts (supplemented with the DR9 catalog at
Dec\textgreater+32.375~deg). All targets are located in the sky
hemisphere where MPE controls the data rights (approx.
180\textless{}\emph{l}\textless360~deg). Due to the LegacySurvey
footprint, nearly all targets in this carton are located at high
Galactic latitudes \textbar{}\emph{b}\textbar\textgreater15~deg.

\noindent\textbf{Simplified description of selection criteria:} Starting from a
parent catalog of eRASS:3 point source $\rightarrow$ legacysurvey.org/DR10(+DR9)
associations (method: NWAY assisted by optical/IR priors computed via a
pre-trained Random Forest, building on
\citealt{Salvato2022}), select targets which meet all of the following criteria:
i) have eROSITA detection likelihood\textgreater6.0, ii) have an X-ray $\rightarrow$
optical/IR cross-match probability (NWAY) of
\texttt{p\_any}\textgreater0.1, iii) have
13.5\textless{}\texttt{fibermag\_r}\textless22.5 or
13.5\textless{}\texttt{fibermag\_i}\textless22.3 or
13.5\textless{}\texttt{fibermag\_z}\textless21.5, iv) are not saturated
in LegacySurvey imaging, v) if detected by Gaia DR2/DR3 then have
\emph{G}\textgreater13.5 and \emph{RP}\textgreater13.5~Vega. We
\emph{deprioritise} targets if any of the following criteria are met: a)
the target already has existing good quality SDSS spectroscopy, b) the
X-ray detection likelihood is \textless8.0, c) the target is a secondary
X-ray$\rightarrow$optical/IR association, d) the target falls outside the core
magnitude range (i.e. doesn't satisfy one of
16.5\textless{}\texttt{fibermag\_r}\textless22.0 or
16.5\textless{}\texttt{fibermag\_i}\textless21.8 or
16.5\textless{}\texttt{fibermag\_z}\textless21.0), e) the target lies at
\textbar{}\emph{b}\textbar\textless15~deg or at Dec\textless-75~deg, f)
the X-ray flux is
\textless2x10\textsuperscript{-14}~erg~s\textsuperscript{-1}~cm\textsuperscript{-2}.
We slightly boost the priorities of targets which have hard X-ray
detections. The '\_d3' variation of this carton provides a shallower
cadence option for targets toward the optically faint end of the sample.


\noindent\textbf{Target priority options:} 1521-1528, 1721-1728, 3521-3528,
3721-3728

\noindent\textbf{Cadence options:} \texttt{dark\_flexible\_3x1}

\noindent\textbf{Implementation:}
\href{https://github.com/sdss/target_selection/blob/1.3.3/python/target_selection/cartons/bhm_spiders_agn.py}{bhm\_spiders\_agn.py}

\noindent\textbf{Number of targets:} 851160

\begin{center}\rule{0.5\linewidth}{0.5pt}\end{center}

\hypertarget{bhm_colr_galaxies_lsdr10_d3_plan1.2.16}{%
\subsection{bhm\_colr\_galaxies\_lsdr10\_d3}\label{bhm_colr_galaxies_lsdr10_d3_plan1.2.16}}

\noindent\textbf{target\_selection plan:} 1.2.16

\noindent\textbf{target\_selection tag:}
\href{https://github.com/sdss/target_selection/tree/1.3.15/}{1.3.15}

\noindent\textbf{crossmatch plan:} 1.0.0

\noindent\textbf{Summary:} A supplementary magnitude limited sample of optically
bright galaxies selected from the LegacySurvey DR10 optical/IR imaging
catalog (supplemented with the LegacySurvey DR9 catalog at
Dec\textgreater+32.375~deg). Selection is based on optical morphology,
lack of Gaia parallax, and several magnitude cuts.

\noindent\textbf{Simplified description of selection criteria:} Starting from the
\texttt{legacy\_survey\_dr10} catalog (lsdr10), select entries
satisfying all of the following criteria: i) lsdr10 morphological
\texttt{type} != 'PSF', ii) lsdr10
\texttt{shape\_r}\textgreater1.0~arcsec, iii) no detection of parallax
by Gaia, iv) dereddened \emph{z}-band model mag\textless19.0~AB, and
dereddened \emph{z}-band fiber mag\textless19.5~AB, v) within all the
following apparent fiber mag limits: 16\textless{}\emph{g}\textless22.5~
and 16\textless{}\emph{r}\textless21.5~ and
16\textless{}\emph{z}\textless20.0~AB, vi) if detected in Gaia, then
\emph{G}\textgreater15.0~ and \emph{RP}\textgreater15.0~Vega, vii) at
moderate to high Galactic latitude
\textbar{}\emph{b}\textbar\textgreater15~deg. The '\_d3' variation of
this carton provides a shallower cadence option for targets toward the
optically faint end of the sample.


\noindent\textbf{Target priority options:} 7101

\noindent\textbf{Cadence options:} \texttt{dark\_flexible\_3x1}

\noindent\textbf{Implementation:}
\href{https://github.com/sdss/target_selection/blob/1.3.15/python/target_selection/cartons/bhm_galaxies.py}{bhm\_galaxies.py}

\noindent\textbf{Number of targets:} 6091854

\begin{center}\rule{0.5\linewidth}{0.5pt}\end{center}

\hypertarget{bhm_rm_core_plan1.0.48}{%
\subsection{bhm\_rm\_core}\label{bhm_rm_core_plan1.0.48}}

\noindent\textbf{target\_selection plan:} 1.0.48

\noindent\textbf{target\_selection tag:}
\href{https://github.com/sdss/target_selection/tree/1.0.48/}{1.0.48}

\noindent\textbf{crossmatch plan:} 1.0.0

\noindent\textbf{Summary:} A sample of candidate QSOs selected via the methods
presented by
\citet{Yang2022}. These targets are located within five (+1 backup) well
known survey fields (SDSS-RM, COSMOS, XMM-LSS, ECDFS, CVZ-S/SEP, and
ELIAS-S1).

\noindent\textbf{Simplified description of selection criteria:} Starting from a
parent catalog of optically selected objects in the RM fields (as
presented by
\citealt{Yang2022}), select candidate QSOs that satisfy all of the
following: i) are identified via the Skew-T algorithm
(\texttt{skewt\_qso\ ==\ 1}); ii) have
17\textless{}\texttt{psfmag\_i}\textless21.5~AB
(16\textless{}\emph{G}\textless21.7~AB in the CVZ-S/SEP field); iii) do
not have significant detections (\textgreater3$\sigma$) of parallax and/or
proper motion in Gaia DR2; iv) are not vetoed due to results of visual
inspections of recent spectroscopy; vi) have detections in all of the
gri bands (a Gaia detection is sufficient in the CVZ-S/SEP field);and
vii) do not lie in the SDSS-RM field. This carton version starts with an
updated parent catalog, and fixes an issue which could have caused a few
targets to be unintentionally excluded.


\noindent\textbf{Target priority options:} 1000-1100

\noindent\textbf{Cadence options:} \texttt{dark\_100x8,\ dark\_174x8}

\noindent\textbf{Implementation:}
\href{https://github.com/sdss/target_selection/blob/1.0.48/python/target_selection/cartons/bhm_rm.py}{bhm\_rm.py}

\noindent\textbf{Number of targets:} 3757

\begin{center}\rule{0.5\linewidth}{0.5pt}\end{center}

\hypertarget{bhm_rm_known_spec_plan1.0.48}{%
\subsection{bhm\_rm\_known\_spec}\label{bhm_rm_known_spec_plan1.0.48}}

\noindent\textbf{target\_selection plan:} 1.0.48

\noindent\textbf{target\_selection tag:}
\href{https://github.com/sdss/target_selection/tree/1.0.48/}{1.0.48}

\noindent\textbf{crossmatch plan:} 1.0.0

\noindent\textbf{Summary:} A sample of known QSOs identified through optical
spectroscopy from various projects, as collated by
\citet{Yang2022}. These targets are located within five (+1 backup) well
known survey fields (SDSS-RM, COSMOS, XMM-LSS, ECDFS, CVZ-S/SEP, and
ELIAS-S1).

\noindent\textbf{Simplified description of selection criteria:} Starting from a
parent catalog of optically selected objects in the RM fields (as
presented by
\citealt{Yang2022}), select targets which satisfy all of the following: i)
are flagged as having a spectroscopic identification (in the parent
catalog); ii) have 15\textless{}\texttt{psfmag\_i}\textless21.7~AB
(SDSS-RM, CDFS, ELIAS-S1 field), 16\textless{}\emph{G}\textless21.7~Vega
(CVZ-S/SEP field), 15\textless{}\texttt{psfmag\_i}\textless21.5~AB
(COSMOS and XMM-LSS fields); iii) have a spectroscopic redshift in the
range 0.005\textless{}\emph{z}\textless7; iv) are not vetoed due to
results of visual inspections of recent spectroscopy. This carton
version starts with an updated parent catalog, and fixes an issue which
could have caused a few targets to be unintentionally excluded.


\noindent\textbf{Target priority options:} 1000-1100

\noindent\textbf{Cadence options:} \texttt{dark\_100x8,\ dark\_174x8}

\noindent\textbf{Implementation:}
\href{https://github.com/sdss/target_selection/blob/1.0.48/python/target_selection/cartons/bhm_rm.py}{bhm\_rm.py}

\noindent\textbf{Number of targets:} 3212

\begin{center}\rule{0.5\linewidth}{0.5pt}\end{center}

\hypertarget{bhm_rm_var_plan1.0.48}{%
\subsection{bhm\_rm\_var}\label{bhm_rm_var_plan1.0.48}}

\noindent\textbf{target\_selection plan:} 1.0.48

\noindent\textbf{target\_selection tag:}
\href{https://github.com/sdss/target_selection/tree/1.0.48/}{1.0.48}

\noindent\textbf{crossmatch plan:} 1.0.0

\noindent\textbf{Summary:} A sample of candidate QSOs selected via their optical
variability properties, as presented by
\citet{Yang2022}. These targets are located within five (+1 backup) well
known survey fields (SDSS-RM, COSMOS, XMM-LSS, ECDFS, CVZ-S/SEP, and
ELIAS-S1).

\noindent\textbf{Simplified description of selection criteria:} Starting from a
parent catalog of optically selected objects in the RM fields (as
presented by
\citealt{Yang2022}), select candidate QSOs that satisfy all of the
following: i) have significant variability in the DES or PanSTARRS1
multi-epoch photometry (\texttt{var\_sn{[}g{]}}\textgreater3 and
\texttt{var\_rms{[}g{]}}\textgreater0.05); ii) have
17\textless{}\texttt{psfmag\_i}\textless20.5~AB
(16\textless{}\emph{G}\textless21.7~AB in the CVZ-S/SEP field); iii) do
not have significant detections (\textgreater3$\sigma$) of parallax and/or
proper motion in Gaia DR2; iv) are not vetoed due to results of visual
inspections of recent spectroscopy; and vi) do not lie in the SDSS-RM
field. This carton version starts with an updated parent catalog, and
fixes an issue which could have caused a few targets to be
unintentionally excluded.


\noindent\textbf{Target priority options:} 1000-1100

\noindent\textbf{Cadence options:} \texttt{dark\_174x8}

\noindent\textbf{Implementation:}
\href{https://github.com/sdss/target_selection/blob/1.0.48/python/target_selection/cartons/bhm_rm.py}{bhm\_rm.py}

\noindent\textbf{Number of targets:} 871

\begin{center}\rule{0.5\linewidth}{0.5pt}\end{center}

\hypertarget{bhm_rm_ancillary_plan1.0.48}{%
\subsection{bhm\_rm\_ancillary}\label{bhm_rm_ancillary_plan1.0.48}}

\noindent\textbf{target\_selection plan:} 1.0.48

\noindent\textbf{target\_selection tag:}
\href{https://github.com/sdss/target_selection/tree/1.0.48/}{1.0.48}

\noindent\textbf{crossmatch plan:} 1.0.0

\noindent\textbf{Summary:} A supporting sample of candidate QSOs which have been
selected by the Gaia-unWISE AGN catalog
(\citealt{Shu2019}) and/or the SDSS XDQSO catalog
(\citealt{Bovy2011}). These targets are located within five (+1 backup) well
known survey fields (SDSS-RM, COSMOS, XMM-LSS, ECDFS, CVZ-S/SEP, and
ELIAS-S1).

\noindent\textbf{Simplified description of selection criteria:} Starting from a
parent catalog of optically selected objects in the RM fields (as
presented by
\citealt{Yang2022}), select candidate QSOs that satisfy all of the
following: i) are identified via external ancillary methods
(\texttt{photo\_bitmask\ \&\ 3\ !=\ 0}); ii) have
15\textless{}\texttt{psfmag\_i}\textless21.5~AB
(16\textless{}\emph{G}\textless21.7~AB in the CVZ-S/SEP field); iii) do
not have significant detections (\textgreater3$\sigma$) of parallax and/or
proper motion in Gaia DR2; iv) are not vetoed due to results of visual
inspections of recent spectroscopy; and v) do not lie in the SDSS-RM
field. This carton version starts with an updated parent catalog, and
fixes an issue which could have caused a few targets to be
unintentionally excluded.


\noindent\textbf{Target priority options:} 1000-1100

\noindent\textbf{Cadence options:} \texttt{dark\_100x8,\ dark\_174x8}

\noindent\textbf{Implementation:}
\href{https://github.com/sdss/target_selection/blob/1.0.48/python/target_selection/cartons/bhm_rm.py}{bhm\_rm.py}

\noindent\textbf{Number of targets:} 1114

\begin{center}\rule{0.5\linewidth}{0.5pt}\end{center}

\hypertarget{manual_bhm_spiders_comm_lco_plan0.5.23}{%
\subsection{manual\_bhm\_spiders\_comm\_lco}\label{manual_bhm_spiders_comm_lco_plan0.5.23}}

\noindent\textbf{target\_selection plan:} 0.5.23

\noindent\textbf{target\_selection tag:}
\href{https://github.com/sdss/target_selection/tree/0.3.22/}{0.3.22}

\noindent\textbf{crossmatch plan:} 0.5.0

\noindent\textbf{Summary:} A utility sample of X-ray selected targets located in
the XMM-XXL-N and Stripe-82X fields having reliable spectroscopic
redshifts. This carton was designed to support SDSS-V commissioning (of
FPS at LCO), as a reference sample to test the accuracy and completeness
of pipeline-derived redshifts extracted from new SDSS-V spectroscopic
data.

\noindent\textbf{Simplified description of selection criteria:} Starting from a
compilation of X-ray selected sources in the XMM-XXL-N and Stripe-82X
fields with optical counterparts in the LegacySurvey DR8 catalog, select
those having i) reliable spectroscopic redshifts (confirmed visually)
from previous SDSS observations, and ii) which lie in the magnitude
range 16\textless{}\texttt{mag\_r}\textless23.0~AB.


\noindent\textbf{Target priority options:} 1000

\noindent\textbf{Cadence options:} \texttt{dark\_1x4}

\noindent\textbf{Implementation:} n/a

\noindent\textbf{Number of targets:} 3160
